\documentclass[11pt]{report}
% Neuron document style: Charter + Noto Serif CJK KR, light outlined boxes, grayscale.
% Usage: \documentclass[11pt]{report} % Neuron document style: Charter + Noto Serif CJK KR, light outlined boxes, grayscale.
% Usage: \documentclass[11pt]{report} % Neuron document style: Charter + Noto Serif CJK KR, light outlined boxes, grayscale.
% Usage: \documentclass[11pt]{report} \input{neuron-style}  (compile with xelatex)
\usepackage[a4paper,margin=24mm,headheight=14pt]{geometry}
\usepackage{fontspec}
\setmainfont{Charter}[Path=fonts/,Extension=.otf,UprightFont=*-Regular,ItalicFont=*-Italic,
  BoldFont=*-Bold,BoldItalicFont=*-BoldItalic,Ligatures=TeX]
\setsansfont{lmsans10-regular.otf}[BoldFont=lmsans10-bold.otf,ItalicFont=lmsans10-oblique.otf]
\setmonofont{lmmono10-regular.otf}[ItalicFont=lmmono10-italic.otf]
\newfontfamily\krfont{Nanum Myeongjo}[Scale=0.9]
\XeTeXinterchartokenstate=1
\newXeTeXintercharclass\krclass
\count255="AC00 \loop\ifnum\count255<"D7A4 \XeTeXcharclass\count255=\krclass \advance\count255 by 1 \repeat
\count255="3000 \loop\ifnum\count255<"3040 \XeTeXcharclass\count255=\krclass \advance\count255 by 1 \repeat
\count255="3130 \loop\ifnum\count255<"3190 \XeTeXcharclass\count255=\krclass \advance\count255 by 1 \repeat
\count255="1100 \loop\ifnum\count255<"1200 \XeTeXcharclass\count255=\krclass \advance\count255 by 1 \repeat
\XeTeXinterchartoks 0 \krclass = {\begingroup\krfont}
\XeTeXinterchartoks 4095 \krclass = {\begingroup\krfont}
\XeTeXinterchartoks \krclass 0 = {\endgroup}
\XeTeXinterchartoks \krclass 4095 = {\endgroup}
\XeTeXlinebreaklocale "ko"
\XeTeXlinebreakskip=0pt plus 1pt
\usepackage{amsmath,amssymb,amsthm,mathtools,bm}
\usepackage[italic,defaultmathsizes,noasterisk]{mathastext}
\emergencystretch=2.5em

\usepackage{graphicx,booktabs,tabularx,longtable,array,multirow,enumitem,xcolor}
\usepackage{caption}
\captionsetup{font=small,labelfont=bf}
\usepackage{tikz}
\usetikzlibrary{shapes.geometric,arrows.meta,positioning,fit,calc,backgrounds,decorations.pathreplacing}
\usepackage[most]{tcolorbox}
\usepackage{listings}
\usepackage{fancyhdr}
\usepackage[hidelinks]{hyperref}

\setlist{itemsep=1pt,topsep=3pt,parsep=0pt}
\setlength{\parskip}{3pt}
\renewcommand{\arraystretch}{1.12}

\lstdefinestyle{sql}{language=SQL,basicstyle=\ttfamily\small,keywordstyle=\bfseries,
  commentstyle=\itshape\color{black!60},columns=fullflexible,keepspaces=true,
  frame=l,framerule=0.6pt,rulecolor=\color{black!40},xleftmargin=8pt,framexleftmargin=4pt,
  aboveskip=5pt,belowskip=5pt,showstringspaces=false,breaklines=true,
  morekeywords={OVER,PARTITION,RANK,DENSE_RANK,ROW_NUMBER,NTILE,LAG,LEAD,ROWS,RANGE,PRECEDING,FOLLOWING,UNBOUNDED,CURRENT,ROW,NVL,NVL2,REFERENCES,CASCADE,MINUS,INTERSECT,EXISTS,MATERIALIZED,VIEW,REPLACE,RATIO_TO_REPORT,FIRST_VALUE,LAST_VALUE,INTERVAL,NULLS,LAST,FIRST,CHECK,CONSTRAINT,UNIQUE,DEFERRABLE,COMMIT,BEGIN,ROLLBACK,TRANSACTION,INDEX,DEFAULT,ANY,ALL,SOME,ROUND}}
\lstset{style=sql}

\tcbset{
 nb/.style={enhanced,breakable,colback=white,coltitle=black,fonttitle=\bfseries,
   left=3mm,right=3mm,top=1.6mm,bottom=1.6mm,before skip=9pt,after skip=9pt,parbox=false},
 framed/.style={nb,boxrule=0.5pt,colframe=black!70,arc=1.2mm,colbacktitle=white,
   titlerule=0.4pt,titlerule style={black!25},toptitle=1mm,bottomtitle=0.6mm},
 light/.style={framed,colframe=black!35},
 ruled/.style={nb,boxrule=0pt,leftrule=1.1pt,colframe=black!45,sharp corners,left=3.5mm,
   top=0.5mm,bottom=0.5mm,detach title,before upper={\tcbtitle\par\smallskip}},
}
\DeclareTColorBox{keybox}{o}{framed,leftrule=2.2pt,title={\IfValueTF{#1}{#1}{Key idea}}}
\DeclareTColorBox{formulabox}{o}{framed,IfValueTF={#1}{title={#1}}{}}
\DeclareTColorBox{workout}{o}{light,title={Example\IfValueT{#1}{\ \textmd{---}\ #1}}}
\DeclareTColorBox{qa}{m}{light,title={#1}}
\DeclareTColorBox{warnbox}{o}{ruled,leftrule=2.2pt,colframe=black,title={Caution\IfValueT{#1}{\ \textmd{---}\ #1}}}
\DeclareTColorBox{whybox}{o}{ruled,title={\IfValueTF{#1}{#1}{Why it works}}}
\DeclareTColorBox{prof}{o}{ruled,title={Lecture note\IfValueT{#1}{\ \textmd{---}\ #1}}}
\DeclareTColorBox{career}{o}{ruled,title={Career note\IfValueT{#1}{\ \textmd{---}\ #1}}}
\DeclareTColorBox{insight}{o}{ruled,title={In practice\IfValueT{#1}{\ \textmd{---}\ #1}}}

\newcommand{\Step}[1]{\par\smallskip\noindent\textbf{Step #1.}\ }
\newcommand{\Concept}[1]{\hfill{\small[개념: #1]}}
\newcommand{\Use}[1]{\noindent\textit{사용 개념:} #1\par}
\newcommand{\Ans}[1]{\par\smallskip\noindent\fbox{\parbox{\dimexpr\linewidth-2\fboxsep-2\fboxrule\relax}{\textbf{답}\ \ #1}}\par}
\newcommand{\Trap}[1]{\par\smallskip\noindent\textit{흔한 실수:} #1}
\newcommand{\Verify}[1]{\par\smallskip\noindent\textit{검산:} #1}
\newcommand{\Intuit}[1]{\par\smallskip\textbf{직관: #1}}
\newcommand{\pk}[1]{\underline{#1}}
\newcommand{\NULL}{\textsc{null}}
\newcommand{\sqlk}[1]{\texttt{#1}}
\newcommand{\Nul}{{\color{black!50}\textit{null}}}

% relational algebra symbols
\newcommand{\sel}{\sigma}
\newcommand{\proj}{\pi}
\newcommand{\join}{\bowtie}

\pagestyle{fancy}
\fancyhf{}
\fancyhead[R]{\small\itshape\leftmark}
\fancyfoot[C]{\small\thepage}
\renewcommand{\headrulewidth}{0pt}
\renewcommand{\chaptermark}[1]{\markboth{#1}{}}

\usepackage{titlesec}
\titleformat{\chapter}[display]{\bfseries\Large}{\chaptertitlename\ \thechapter}{6pt}{\LARGE}
\titlespacing*{\chapter}{0pt}{-10pt}{14pt}
\titleformat{\section}{\bfseries\large}{\thesection}{0.8em}{}
\titleformat{\subsection}{\bfseries\normalsize}{\thesubsection}{0.8em}{}

% ---------- cheat-sheet / problem-sheet helpers ----------
\newcounter{qn}
\newcommand{\Q}[2]{\par\smallskip\refstepcounter{qn}\noindent\textbf{Q\theqn.}\ {\footnotesize[#1]}\ #2\par}
\newcommand{\A}[1]{\vspace{1pt}\noindent\begin{tcolorbox}[ruled,leftrule=1.6pt,colframe=black!45,before skip=1pt,after skip=4pt,top=0.5mm,bottom=0.5mm]\small\textbf{답} #1\end{tcolorbox}}
\newcommand{\Lab}[3]{\item \textbf{#1}\\ \textit{해 본 것:} #2\\ \textit{\ensuremath{\Rightarrow} 개념:} #3}
\newcommand{\eg}{\textit{예)}\ }
\newcommand{\bi}{\textit{비유)}\ }
  (compile with xelatex)
\usepackage[a4paper,margin=24mm,headheight=14pt]{geometry}
\usepackage{fontspec}
\setmainfont{Charter}[Path=fonts/,Extension=.otf,UprightFont=*-Regular,ItalicFont=*-Italic,
  BoldFont=*-Bold,BoldItalicFont=*-BoldItalic,Ligatures=TeX]
\setsansfont{lmsans10-regular.otf}[BoldFont=lmsans10-bold.otf,ItalicFont=lmsans10-oblique.otf]
\setmonofont{lmmono10-regular.otf}[ItalicFont=lmmono10-italic.otf]
\newfontfamily\krfont{Nanum Myeongjo}[Scale=0.9]
\XeTeXinterchartokenstate=1
\newXeTeXintercharclass\krclass
\count255="AC00 \loop\ifnum\count255<"D7A4 \XeTeXcharclass\count255=\krclass \advance\count255 by 1 \repeat
\count255="3000 \loop\ifnum\count255<"3040 \XeTeXcharclass\count255=\krclass \advance\count255 by 1 \repeat
\count255="3130 \loop\ifnum\count255<"3190 \XeTeXcharclass\count255=\krclass \advance\count255 by 1 \repeat
\count255="1100 \loop\ifnum\count255<"1200 \XeTeXcharclass\count255=\krclass \advance\count255 by 1 \repeat
\XeTeXinterchartoks 0 \krclass = {\begingroup\krfont}
\XeTeXinterchartoks 4095 \krclass = {\begingroup\krfont}
\XeTeXinterchartoks \krclass 0 = {\endgroup}
\XeTeXinterchartoks \krclass 4095 = {\endgroup}
\XeTeXlinebreaklocale "ko"
\XeTeXlinebreakskip=0pt plus 1pt
\usepackage{amsmath,amssymb,amsthm,mathtools,bm}
\usepackage[italic,defaultmathsizes,noasterisk]{mathastext}
\emergencystretch=2.5em

\usepackage{graphicx,booktabs,tabularx,longtable,array,multirow,enumitem,xcolor}
\usepackage{caption}
\captionsetup{font=small,labelfont=bf}
\usepackage{tikz}
\usetikzlibrary{shapes.geometric,arrows.meta,positioning,fit,calc,backgrounds,decorations.pathreplacing}
\usepackage[most]{tcolorbox}
\usepackage{listings}
\usepackage{fancyhdr}
\usepackage[hidelinks]{hyperref}

\setlist{itemsep=1pt,topsep=3pt,parsep=0pt}
\setlength{\parskip}{3pt}
\renewcommand{\arraystretch}{1.12}

\lstdefinestyle{sql}{language=SQL,basicstyle=\ttfamily\small,keywordstyle=\bfseries,
  commentstyle=\itshape\color{black!60},columns=fullflexible,keepspaces=true,
  frame=l,framerule=0.6pt,rulecolor=\color{black!40},xleftmargin=8pt,framexleftmargin=4pt,
  aboveskip=5pt,belowskip=5pt,showstringspaces=false,breaklines=true,
  morekeywords={OVER,PARTITION,RANK,DENSE_RANK,ROW_NUMBER,NTILE,LAG,LEAD,ROWS,RANGE,PRECEDING,FOLLOWING,UNBOUNDED,CURRENT,ROW,NVL,NVL2,REFERENCES,CASCADE,MINUS,INTERSECT,EXISTS,MATERIALIZED,VIEW,REPLACE,RATIO_TO_REPORT,FIRST_VALUE,LAST_VALUE,INTERVAL,NULLS,LAST,FIRST,CHECK,CONSTRAINT,UNIQUE,DEFERRABLE,COMMIT,BEGIN,ROLLBACK,TRANSACTION,INDEX,DEFAULT,ANY,ALL,SOME,ROUND}}
\lstset{style=sql}

\tcbset{
 nb/.style={enhanced,breakable,colback=white,coltitle=black,fonttitle=\bfseries,
   left=3mm,right=3mm,top=1.6mm,bottom=1.6mm,before skip=9pt,after skip=9pt,parbox=false},
 framed/.style={nb,boxrule=0.5pt,colframe=black!70,arc=1.2mm,colbacktitle=white,
   titlerule=0.4pt,titlerule style={black!25},toptitle=1mm,bottomtitle=0.6mm},
 light/.style={framed,colframe=black!35},
 ruled/.style={nb,boxrule=0pt,leftrule=1.1pt,colframe=black!45,sharp corners,left=3.5mm,
   top=0.5mm,bottom=0.5mm,detach title,before upper={\tcbtitle\par\smallskip}},
}
\DeclareTColorBox{keybox}{o}{framed,leftrule=2.2pt,title={\IfValueTF{#1}{#1}{Key idea}}}
\DeclareTColorBox{formulabox}{o}{framed,IfValueTF={#1}{title={#1}}{}}
\DeclareTColorBox{workout}{o}{light,title={Example\IfValueT{#1}{\ \textmd{---}\ #1}}}
\DeclareTColorBox{qa}{m}{light,title={#1}}
\DeclareTColorBox{warnbox}{o}{ruled,leftrule=2.2pt,colframe=black,title={Caution\IfValueT{#1}{\ \textmd{---}\ #1}}}
\DeclareTColorBox{whybox}{o}{ruled,title={\IfValueTF{#1}{#1}{Why it works}}}
\DeclareTColorBox{prof}{o}{ruled,title={Lecture note\IfValueT{#1}{\ \textmd{---}\ #1}}}
\DeclareTColorBox{career}{o}{ruled,title={Career note\IfValueT{#1}{\ \textmd{---}\ #1}}}
\DeclareTColorBox{insight}{o}{ruled,title={In practice\IfValueT{#1}{\ \textmd{---}\ #1}}}

\newcommand{\Step}[1]{\par\smallskip\noindent\textbf{Step #1.}\ }
\newcommand{\Concept}[1]{\hfill{\small[개념: #1]}}
\newcommand{\Use}[1]{\noindent\textit{사용 개념:} #1\par}
\newcommand{\Ans}[1]{\par\smallskip\noindent\fbox{\parbox{\dimexpr\linewidth-2\fboxsep-2\fboxrule\relax}{\textbf{답}\ \ #1}}\par}
\newcommand{\Trap}[1]{\par\smallskip\noindent\textit{흔한 실수:} #1}
\newcommand{\Verify}[1]{\par\smallskip\noindent\textit{검산:} #1}
\newcommand{\Intuit}[1]{\par\smallskip\textbf{직관: #1}}
\newcommand{\pk}[1]{\underline{#1}}
\newcommand{\NULL}{\textsc{null}}
\newcommand{\sqlk}[1]{\texttt{#1}}
\newcommand{\Nul}{{\color{black!50}\textit{null}}}

% relational algebra symbols
\newcommand{\sel}{\sigma}
\newcommand{\proj}{\pi}
\newcommand{\join}{\bowtie}

\pagestyle{fancy}
\fancyhf{}
\fancyhead[R]{\small\itshape\leftmark}
\fancyfoot[C]{\small\thepage}
\renewcommand{\headrulewidth}{0pt}
\renewcommand{\chaptermark}[1]{\markboth{#1}{}}

\usepackage{titlesec}
\titleformat{\chapter}[display]{\bfseries\Large}{\chaptertitlename\ \thechapter}{6pt}{\LARGE}
\titlespacing*{\chapter}{0pt}{-10pt}{14pt}
\titleformat{\section}{\bfseries\large}{\thesection}{0.8em}{}
\titleformat{\subsection}{\bfseries\normalsize}{\thesubsection}{0.8em}{}

% ---------- cheat-sheet / problem-sheet helpers ----------
\newcounter{qn}
\newcommand{\Q}[2]{\par\smallskip\refstepcounter{qn}\noindent\textbf{Q\theqn.}\ {\footnotesize[#1]}\ #2\par}
\newcommand{\A}[1]{\vspace{1pt}\noindent\begin{tcolorbox}[ruled,leftrule=1.6pt,colframe=black!45,before skip=1pt,after skip=4pt,top=0.5mm,bottom=0.5mm]\small\textbf{답} #1\end{tcolorbox}}
\newcommand{\Lab}[3]{\item \textbf{#1}\\ \textit{해 본 것:} #2\\ \textit{\ensuremath{\Rightarrow} 개념:} #3}
\newcommand{\eg}{\textit{예)}\ }
\newcommand{\bi}{\textit{비유)}\ }
  (compile with xelatex)
\usepackage[a4paper,margin=24mm,headheight=14pt]{geometry}
\usepackage{fontspec}
\setmainfont{Charter}[Path=fonts/,Extension=.otf,UprightFont=*-Regular,ItalicFont=*-Italic,
  BoldFont=*-Bold,BoldItalicFont=*-BoldItalic,Ligatures=TeX]
\setsansfont{lmsans10-regular.otf}[BoldFont=lmsans10-bold.otf,ItalicFont=lmsans10-oblique.otf]
\setmonofont{lmmono10-regular.otf}[ItalicFont=lmmono10-italic.otf]
\newfontfamily\krfont{Nanum Myeongjo}[Scale=0.9]
\XeTeXinterchartokenstate=1
\newXeTeXintercharclass\krclass
\count255="AC00 \loop\ifnum\count255<"D7A4 \XeTeXcharclass\count255=\krclass \advance\count255 by 1 \repeat
\count255="3000 \loop\ifnum\count255<"3040 \XeTeXcharclass\count255=\krclass \advance\count255 by 1 \repeat
\count255="3130 \loop\ifnum\count255<"3190 \XeTeXcharclass\count255=\krclass \advance\count255 by 1 \repeat
\count255="1100 \loop\ifnum\count255<"1200 \XeTeXcharclass\count255=\krclass \advance\count255 by 1 \repeat
\XeTeXinterchartoks 0 \krclass = {\begingroup\krfont}
\XeTeXinterchartoks 4095 \krclass = {\begingroup\krfont}
\XeTeXinterchartoks \krclass 0 = {\endgroup}
\XeTeXinterchartoks \krclass 4095 = {\endgroup}
\XeTeXlinebreaklocale "ko"
\XeTeXlinebreakskip=0pt plus 1pt
\usepackage{amsmath,amssymb,amsthm,mathtools,bm}
\usepackage[italic,defaultmathsizes,noasterisk]{mathastext}
\emergencystretch=2.5em

\usepackage{graphicx,booktabs,tabularx,longtable,array,multirow,enumitem,xcolor}
\usepackage{caption}
\captionsetup{font=small,labelfont=bf}
\usepackage{tikz}
\usetikzlibrary{shapes.geometric,arrows.meta,positioning,fit,calc,backgrounds,decorations.pathreplacing}
\usepackage[most]{tcolorbox}
\usepackage{listings}
\usepackage{fancyhdr}
\usepackage[hidelinks]{hyperref}

\setlist{itemsep=1pt,topsep=3pt,parsep=0pt}
\setlength{\parskip}{3pt}
\renewcommand{\arraystretch}{1.12}

\lstdefinestyle{sql}{language=SQL,basicstyle=\ttfamily\small,keywordstyle=\bfseries,
  commentstyle=\itshape\color{black!60},columns=fullflexible,keepspaces=true,
  frame=l,framerule=0.6pt,rulecolor=\color{black!40},xleftmargin=8pt,framexleftmargin=4pt,
  aboveskip=5pt,belowskip=5pt,showstringspaces=false,breaklines=true,
  morekeywords={OVER,PARTITION,RANK,DENSE_RANK,ROW_NUMBER,NTILE,LAG,LEAD,ROWS,RANGE,PRECEDING,FOLLOWING,UNBOUNDED,CURRENT,ROW,NVL,NVL2,REFERENCES,CASCADE,MINUS,INTERSECT,EXISTS,MATERIALIZED,VIEW,REPLACE,RATIO_TO_REPORT,FIRST_VALUE,LAST_VALUE,INTERVAL,NULLS,LAST,FIRST,CHECK,CONSTRAINT,UNIQUE,DEFERRABLE,COMMIT,BEGIN,ROLLBACK,TRANSACTION,INDEX,DEFAULT,ANY,ALL,SOME,ROUND}}
\lstset{style=sql}

\tcbset{
 nb/.style={enhanced,breakable,colback=white,coltitle=black,fonttitle=\bfseries,
   left=3mm,right=3mm,top=1.6mm,bottom=1.6mm,before skip=9pt,after skip=9pt,parbox=false},
 framed/.style={nb,boxrule=0.5pt,colframe=black!70,arc=1.2mm,colbacktitle=white,
   titlerule=0.4pt,titlerule style={black!25},toptitle=1mm,bottomtitle=0.6mm},
 light/.style={framed,colframe=black!35},
 ruled/.style={nb,boxrule=0pt,leftrule=1.1pt,colframe=black!45,sharp corners,left=3.5mm,
   top=0.5mm,bottom=0.5mm,detach title,before upper={\tcbtitle\par\smallskip}},
}
\DeclareTColorBox{keybox}{o}{framed,leftrule=2.2pt,title={\IfValueTF{#1}{#1}{Key idea}}}
\DeclareTColorBox{formulabox}{o}{framed,IfValueTF={#1}{title={#1}}{}}
\DeclareTColorBox{workout}{o}{light,title={Example\IfValueT{#1}{\ \textmd{---}\ #1}}}
\DeclareTColorBox{qa}{m}{light,title={#1}}
\DeclareTColorBox{warnbox}{o}{ruled,leftrule=2.2pt,colframe=black,title={Caution\IfValueT{#1}{\ \textmd{---}\ #1}}}
\DeclareTColorBox{whybox}{o}{ruled,title={\IfValueTF{#1}{#1}{Why it works}}}
\DeclareTColorBox{prof}{o}{ruled,title={Lecture note\IfValueT{#1}{\ \textmd{---}\ #1}}}
\DeclareTColorBox{career}{o}{ruled,title={Career note\IfValueT{#1}{\ \textmd{---}\ #1}}}
\DeclareTColorBox{insight}{o}{ruled,title={In practice\IfValueT{#1}{\ \textmd{---}\ #1}}}

\newcommand{\Step}[1]{\par\smallskip\noindent\textbf{Step #1.}\ }
\newcommand{\Concept}[1]{\hfill{\small[개념: #1]}}
\newcommand{\Use}[1]{\noindent\textit{사용 개념:} #1\par}
\newcommand{\Ans}[1]{\par\smallskip\noindent\fbox{\parbox{\dimexpr\linewidth-2\fboxsep-2\fboxrule\relax}{\textbf{답}\ \ #1}}\par}
\newcommand{\Trap}[1]{\par\smallskip\noindent\textit{흔한 실수:} #1}
\newcommand{\Verify}[1]{\par\smallskip\noindent\textit{검산:} #1}
\newcommand{\Intuit}[1]{\par\smallskip\textbf{직관: #1}}
\newcommand{\pk}[1]{\underline{#1}}
\newcommand{\NULL}{\textsc{null}}
\newcommand{\sqlk}[1]{\texttt{#1}}
\newcommand{\Nul}{{\color{black!50}\textit{null}}}

% relational algebra symbols
\newcommand{\sel}{\sigma}
\newcommand{\proj}{\pi}
\newcommand{\join}{\bowtie}

\pagestyle{fancy}
\fancyhf{}
\fancyhead[R]{\small\itshape\leftmark}
\fancyfoot[C]{\small\thepage}
\renewcommand{\headrulewidth}{0pt}
\renewcommand{\chaptermark}[1]{\markboth{#1}{}}

\usepackage{titlesec}
\titleformat{\chapter}[display]{\bfseries\Large}{\chaptertitlename\ \thechapter}{6pt}{\LARGE}
\titlespacing*{\chapter}{0pt}{-10pt}{14pt}
\titleformat{\section}{\bfseries\large}{\thesection}{0.8em}{}
\titleformat{\subsection}{\bfseries\normalsize}{\thesubsection}{0.8em}{}

% ---------- cheat-sheet / problem-sheet helpers ----------
\newcounter{qn}
\newcommand{\Q}[2]{\par\smallskip\refstepcounter{qn}\noindent\textbf{Q\theqn.}\ {\footnotesize[#1]}\ #2\par}
\newcommand{\A}[1]{\vspace{1pt}\noindent\begin{tcolorbox}[ruled,leftrule=1.6pt,colframe=black!45,before skip=1pt,after skip=4pt,top=0.5mm,bottom=0.5mm]\small\textbf{답} #1\end{tcolorbox}}
\newcommand{\Lab}[3]{\item \textbf{#1}\\ \textit{해 본 것:} #2\\ \textit{\ensuremath{\Rightarrow} 개념:} #3}
\newcommand{\eg}{\textit{예)}\ }
\newcommand{\bi}{\textit{비유)}\ }

\usepackage{float}
\usepackage{multicol}
\newcommand{\Hd}[1]{\par\smallskip\noindent\textbf{#1}\par}
\usetikzlibrary{arrows.meta,positioning,calc}
\usepackage{pgfplots}\pgfplotsset{compat=1.18}
\newcommand{\ft}[1]{\texttt{\detokenize{#1}}}
\newcolumntype{L}{>{\raggedright\arraybackslash}X}
\lstdefinestyle{py}{language=Python,basicstyle=\ttfamily\small,keywordstyle=\bfseries,commentstyle=\itshape\color{black!55},
  frame=leftline,framerule=0.8pt,rulecolor=\color{black!45},xleftmargin=3mm,showstringspaces=false,breaklines=true,columns=fullflexible}
\lstdefinestyle{out}{basicstyle=\ttfamily\footnotesize,frame=single,rulecolor=\color{black!25},xleftmargin=3mm,columns=fullflexible,breaklines=true,numbers=none}
\lstset{style=py}
\begin{document}

\begin{titlepage}\centering
\vspace*{30mm}
{\Huge\bfseries LendingClub 대출 승인 전략\par}\medskip
{\Large 팀 보고서: 문제의식, 방법, 검증, 결론\par}
\vspace{12mm}
{\large 서울대 ABS 14기 팀 프로젝트 내부용\par}
\vspace{4mm}
{\small 2026년 11월 18일 개정 \ensuremath{\cdot} 1--42단계 실험 \ensuremath{\cdot} 숫자는 모두 공식 train 데이터 결과(예시라고 적은 것 제외)\par}
\vfill
{\small 기초(1--4장)는 강의노트와 같은 내용이다. 처음 읽는 사람은 요약을 읽은 뒤 1장부터, 이미 아는 사람은 5장부터 읽으면 된다.\par}
\end{titlepage}

\tableofcontents
\chapter*{핵심 정리}
\addcontentsline{toc}{chapter}{핵심 정리}
\markboth{핵심 정리}{}
\begingroup
\captionsetup{labelformat=empty}

\begin{keybox}[한 문장 결론]
\textbf{LendingClub 대출은 2013--16년에 전부 사도 국채보다 나았다. 그래서 질문은 ``무엇을 고르느냐''다. 우리는 대출마다 부도확률과 금리로 기대수익을 계산해, 승인비율이 몇 \%든 기대수익이 높은 대출부터 사자고 제안한다. train에서 승인비율 10--90\% 9개 모두 같은 양을 무작위로 산 것보다 더 벌었고, 승인비율 60\% 이상이면 덜 흔들리기도 했다.}
\end{keybox}
숫자는 모두 train(2013--16년 발행 36개월 대출, 48개 발행월) 결과이며, test는 아직 열지 않았다. 판정 기준은 test 전에 고정했다(규칙서 v2.3).

\begin{formulabox}[승인비율 $q$란]
매달 그 달에 나온 36개월 대출을 점수 순으로 줄 세워 위에서부터 사는 \textbf{건수의 비율}이다(10\%, 20\%, \dots, 90\%). 투자자가 그 달 시장의 몇 \%를 살 돈이 있는지를 나타내며, 사지 않은 돈은 국채에 둔다고 본다. 대출마다 금액이 달라 \textbf{금액 비중} $w$는 $q$와 조금 다르다(예: 부도확률 순 80\%에서 평균 82.5\%). 선별 가치는 승인한 1달러당 IRR 차이로, Sharpe는 금액 비중 $w$를 곱한 초과수익으로 계산한다.
\end{formulabox}

\begin{formulabox}[선별 가치와 포트폴리오 초과수익은 다르다]
\textbf{선별 가치}는 \textbf{대출에 넣은 돈의 질}을 잰다. 어느 달 대출이 1,000건이고 승인비율이 10\%면, 점수 상위 100건과 아무 100건을 각각 사서 대출이 끝난 뒤 두 묶음의 실제 IRR을 비교한다. 예를 들어 우리 묶음 연 8.0\%, 무작위 묶음 연 5.9\%면 선별 가치는 2.1\%p다. IRR은 ``넣은 1달러가 1년에 몇 \% 벌었나''이므로, 얼마를 넣었는지와 상관없이 \textbf{같은 양을 무작위로 산 것보다 1달러당 몇 \%p 더 버나}를 뜻한다. 48개월 평균이 표의 숫자다. 고르는 실력을 재는 숫자이고, 주 판정에 쓴다.

\textbf{포트폴리오 초과수익}은 \textbf{전체 돈의 성과}다. 사지 않은 돈은 국채에 두어 초과수익이 0이므로, 승인비율 10\%면 돈의 90\%가 국채보다 더 벌지 못한다. 투자자가 실제로 국채보다 얼마나 더 버는지는 이 숫자가 보여 준다.

\smallskip
\centering\small
\begin{tabular}{@{}lccccc@{}}\toprule
기대수익 순(V1), train & 10\% & 50\% & 80\% & 90\% & 전부 승인\\\midrule
선별 가치(1달러당, 무작위 대비) & +2.11 & +0.99 & +0.48 & +0.28 & 0\\
포트폴리오 초과수익(국채 대비, 전체 돈) & 0.65 & 2.98 & 4.34 & 4.67 & 4.86\\\bottomrule
\end{tabular}

\raggedright
\smallskip
\textbf{읽는 법:} 적게 살수록 가장 좋은 대출만 골라 1달러당 차이가 크지만, 돈 대부분이 국채에 있어 전체 수익은 작다. 많이 살수록 평범한 대출까지 사게 되어 1달러당 차이는 줄지만 전체 수익은 커진다. 그래서 ``10\%만 사는 것이 가장 좋다''는 뜻이 아니다. 이 기간에는 대출이 국채보다 나았으므로 많이 살수록 전체 수익은 컸고, 우리 주장은 \textbf{몇 \%를 사든 같은 양을 무작위로 사는 것보다 낫다}는 것이다.
\end{formulabox}


\section*{흐름}
\begin{figure}[H]\centering
\begin{tikzpicture}[x=1mm,y=1mm,box/.style={draw=black!70,rounded corners=1.2mm,align=center,font=\footnotesize,minimum height=12mm,text width=32mm,fill=white,execute at begin node={\XeTeXlinebreakskip=0pt\relax}},ar/.style={-{Stealth},black!65}]
\node[box] (a) at (0,0) {그 달 새 대출\\(36개월 만기)};
\node[box] (b) at (42,0) {부도확률 $p$\\시점 학습 LightGBM\\(그날 알던 정보만)};
\node[box] (c) at (84,0) {기대수익 점수\\$s=(1-p)r-p\cdot\mathrm{LGD}$\\$-$ 국채 금리};
\node[box,line width=0.9pt] (d) at (126,0) {$s$ 높은 순으로\\예산만큼 구매\\남는 돈은 국채};
\draw[ar] (a)--(b); \draw[ar] (b)--(c); \draw[ar] (c)--(d);
\draw[dashed,black!35] (-18,-13)--(144,-13);
\node[font=\scriptsize\itshape,black!60,anchor=south west] at (-18,-13) {위: 승인하는 날(결과를 모름)};
\node[font=\scriptsize\itshape,black!60,anchor=north west] at (-18,-13) {아래: 대출이 끝난 뒤(실제 결과로 채점, $s$는 쓰지 않음)};
\node[box] (e) at (126,-27) {우리가 산 묶음\\실제 현금흐름\\\ensuremath{\to} 코호트 IRR};
\node[box] (f) at (126,-47) {같은 양 무작위(10회)\\실제 현금흐름\\\ensuremath{\to} 코호트 IRR};
\node[box,line width=0.9pt] (g) at (84,-37) {선별 가치\\$=$ 우리 IRR $-$ 무작위 IRR\\(국채는 양쪽에서 지워짐)};
\node[box] (h) at (42,-37) {48개 발행월 평균\\예산 10--90\%\\블록 부트스트랩 구간};
\node[box,line width=0.9pt] (i) at (0,-27) {주 판정(H2)\\95\% 구간 하한 $>0$\\\ensuremath{\to} 제안 유지};
\node[box,dashed] (j) at (0,-47) {보조 지표\\포트폴리오 Sharpe\\(평균 \ensuremath{\div} 흔들림)};
\draw[ar] (d)--(e); \draw[ar] (e)--(g); \draw[ar] (f)--(g); \draw[ar] (g)--(h); \draw[ar] (h)--(i); \draw[ar,dashed] (h)--(j);
\end{tikzpicture}
\caption{제안의 흐름. 위 줄은 승인하는 날 하는 일, 아래 줄은 대출이 끝난 뒤 채점하는 일이다. 코호트 IRR은 빌려준 돈(0개월), 월 납입금, 마지막 달의 조기상환 일시금이나 부도 후 회수금을 묶음 전체로 더해 구하고, 받은 돈에서 수수료 1\%를 뗀다(2장).}
\end{figure}


\section*{핵심 기술과 근거}
\begin{table}[H]\centering\small
\begin{tabularx}{\linewidth}{@{}>{\raggedright\arraybackslash}p{30mm}>{\raggedright\arraybackslash}p{50mm}L@{}}\toprule
기술 & 무엇을, 왜 & 근거(train 결과)\\\midrule
시점 학습(as-of) & 연도 Y 대출의 $p$는 Y년 1월 1일까지 결과가 확정된 대출로만 학습한 모델로 계산. 그때 투자자가 가질 수 있던 정보만 쓰기 위해 & 무작위 분할 AUC 0.73 vs 시점 학습 0.62--0.68. 차이가 ``미래를 본 값''(3장)\\
LightGBM 부도확률 $p$ & 작은 나무를 차례로 쌓아 앞 나무가 틀린 부분을 고침. 결측을 그대로 받고 변수 조합을 자동으로 잡음. 101개 특성, 시드 3개 평균 & 17종 비교에서 XGBoost와 선별 가치 같음(1.03 vs 1.04\%p). AUC 1위 랜덤포레스트는 돈을 덜 범(0.85\%p). 성과가 특정 모델에 달려 있지 않음(7장)\\
$p$는 순위로 사용 & $p$의 크기는 실제보다 낮게 나오지만 순서는 유지 & 2013년 예측 6.1\% vs 실제 12.2\%. 실제 수준으로 맞춰 다시 골라도 승인 대상 88--92\% 같음(5장)\\
기대수익 점수 $s$ & $s=(1-p)r-p\cdot\mathrm{LGD}-$국채. 예상 손실을 빼고도 국채보다 많이 남는 대출을 찾기 위해(9장 ``왜 $p$가 아니라 $s$인가) & 같은 등급 안 $s$ 상위 $-$ 하위 연 +2.38\%p(2.08--2.73), 48개월 모두 양수(7장)\\
발행월 컷 & 매달 그 달 대출끼리 줄 세워 위에서부터 승인비율 $q$만큼 승인 & 시장 상황이 같은 대출끼리 비교\\
코호트 IRR & 같은 달에 산 대출 묶음의 실제 현금흐름(투자, 월 납입금, 조기상환·\allowbreak{}회수금)으로 구한 연수익률, 수수료 1\% & 단순 연환산이나 끝난 대출만 보면 결론이 뒤집힘(2장)\\
선별 가치(주 판정) & 같은 달, 같은 양을 무작위(10회)로 산 묶음과의 IRR 차이. 국채는 양쪽에서 지워짐 & 아래 결과 표\\
포트폴리오 Sharpe(보조) & 매달 (IRR $-$ 국채) \ensuremath{\times} 승인 비중의 평균 \ensuremath{\div} 표준편차. 덜 흔들리는가 & 아래 결과 표. 약 3으로 큰 것은 만기 보유와 좋은 시기 때문이라 전략끼리 비교에만 사용\\
블록 부트스트랩 & 48개월을 6개월 덩어리로 2000번 다시 뽑아 구간 계산. 이웃 달의 상관을 보존 & 모든 전략 공통 인덱스, 짝지은 차이\\
정규화 & LightGBM에는 불필요(순서로만 나눔). 비교용 로지스틱에는 학습 데이터 기준 자르기, 중앙값 채우기, 표준화 & 표준화 폭발(최대 290, AUC 0.508)을 가용성 규칙으로 해결(0.678)\\
사전 고정과 정정 기록 & test 판정 기준을 test 전에 문서로 고정하고 git 태그. 틀린 설명은 정정표에 남김 & 규칙서 v1--v2.3, 정정표(8장)\\\bottomrule
\end{tabularx}
\caption{핵심 기술. 괄호의 장은 본문 위치.}
\end{table}

\section*{결과: 승인비율별 선별 가치와 Sharpe}
\begin{table}[H]\centering\small
\begin{tabular}{@{}rccccccc@{}}\toprule
 & \multicolumn{2}{c}{선별 가치(\%p)} & \multicolumn{3}{c}{포트폴리오 Sharpe} & \multicolumn{2}{c}{V1 $-$ 무작위}\\\cmidrule(lr){2-3}\cmidrule(lr){4-6}\cmidrule(l){7-8}
승인비율 & \textbf{V1(제안)} & V2 & \textbf{V1} & V2 & 무작위 & 선별 가치 90\% & Sharpe 차이 90\%\\\midrule
10\% & \textbf{+2.11} & $-0.32$ & 2.55 & 3.82 & 2.63 & 1.80, 2.45 & $-0.59$, 0.14\\
20\% & \textbf{+1.74} & $-0.04$ & 2.75 & 3.61 & 2.68 & 1.52, 1.99 & $-0.47$, 0.40\\
30\% & \textbf{+1.46} & +0.18 & 2.91 & 3.34 & 2.68 & 1.29, 1.62 & $-0.28$, 0.58\\
40\% & \textbf{+1.19} & +0.30 & 3.01 & 3.17 & 2.70 & 1.06, 1.33 & $-0.12$, 0.61\\
50\% & \textbf{+0.99} & +0.37 & 3.00 & 3.12 & 2.72 & 0.84, 1.12 & $-0.04$, 0.55\\
60\% & \textbf{+0.82} & +0.39 & 3.04 & 3.11 & 2.72 & 0.68, 0.94 & \textbf{0.08, 0.54}\\
70\% & \textbf{+0.63} & +0.37 & 3.05 & 3.07 & 2.73 & 0.51, 0.74 & \textbf{0.15, 0.52}\\
80\% & \textbf{+0.48} & +0.32 & 3.01 & 3.06 & 2.72 & 0.39, 0.56 & \textbf{0.15, 0.46}\\
90\% & \textbf{+0.28} & +0.21 & 2.93 & 2.98 & 2.72 & 0.22, 0.33 & \textbf{0.14, 0.36}\\\midrule
평균 & \textbf{+1.08} & +0.20 & & & & & \\\bottomrule
\end{tabular}
\caption{V1 = 기대수익 순(제안), V2 = 부도확률 낮은 순(대안, 이전 최종안). 발행월 컷, 국채 기준, 수수료 1\%, 무작위 10회 평균, 6개월 블록 부트스트랩 2000회(42단계, \ft{results/42_final_proposal_check.csv}). 선별 가치 구간은 모두 0보다 위, Sharpe 차이의 굵은 구간도 0보다 위.}
\end{table}

\begin{enumerate}[leftmargin=*]
\item \textbf{더 벌기(주 판정):} V1은 9개 승인비율 모두 무작위보다 더 벌었다(구간 모두 0보다 위). 승인한 1달러당 이점(선별 가치)은 승인비율이 작을수록 크다(10\%: +2.11, 80\%: +0.48\%p). 다만 포트폴리오 전체 초과수익은 많이 살수록 크다(10\%: 0.65, 80\%: 4.34\%p).
\item \textbf{덜 흔들리기(보조):} V1의 Sharpe는 승인비율 60\% 이상에서만 무작위보다 확실히 높다. 50\% 이하에서는 구분되지 않고, 10\%는 점추정이 조금 낮다.
\item \textbf{맞교환:} V2는 Sharpe가 모든 승인비율에서 가장 높지만 더 버는 효과는 작다(평균 +0.20\%p). 덜 흔들리는 것이 가장 중요한 투자자에게는 V2를 대안으로 보고한다. 승인비율이 클 때(70--90\%)는 두 전략의 결과가 거의 같다.
\item \textbf{기준 수익률:} 회사채(BAA) 기준에서는 Sharpe도 V1이 높다(승인 30\%: 1.29 vs V2 0.92, 무작위 0.52).
\end{enumerate}

\section*{test에서 확인하는 것}
\begin{itemize}[leftmargin=*]
\item \textbf{주 판정(v1의 H2 그대로):} 기대수익 순위의 선별 곡선(승인비율 10--90\%)이 무작위보다 위에 있는 면적의 95\% 구간 하한이 0보다 큰가. 통과하면 제안 유지, 실패하면 철회.
\item \textbf{보고만:} 위 표의 승인비율별 숫자, V2 80\%의 무작위 대비 Sharpe(이전 판정 H10), 기준 수익률 4단계. 결과가 어떻든 제안을 바꾸지 않는다.
\end{itemize}

\section*{핵심 한계}
\begin{itemize}[leftmargin=*]
\item 모든 숫자는 train이라 낙관적일 수 있다.
\item 승인비율이 클 때는 1달러당 이점이 작고(80\%: +0.48\%p), 승인비율이 작을 때 덜 흔들린다는 근거는 없다.
\item 2013--16년에는 경기침체가 없었다. 불경기에 모델의 위험 순서가 유지되는지는 확인하지 못했다(부도끼리의 상관, 10장).
\item 주 판정을 v1의 H2로 되돌린 것은 train 결과를 본 뒤의 결정이다. 근거는 과제 공지 재확인(Sharpe 필수 아님)이며, 판정 정의는 바꾸지 않았다(9장).
\end{itemize}

\section*{이름표}
\begin{table}[H]\centering\small
\begin{tabularx}{\linewidth}{@{}p{26mm}L@{}}\toprule
이름 & 뜻\\\midrule
V1 & 기대수익 순. \textbf{최종 제안}\\
V2 & 부도확률 낮은 순. 덜 흔들리기 대안(이전 최종안: V2 80\%)\\
V3 & 같은 등급 안에서 최근 24개월 개설 계좌가 적은 순. train을 보고 고른 도전자\\
v1--v2.3 & test 판정 규칙서의 버전. 현재 유효한 것은 v2.3(주 판정 H2)\\
선별 가치 & 같은 달, 같은 양 무작위 대비 IRR 차이(승인 1달러당, 연 \%p)\\
승인비율 $q$ & 그 달 대출 중 사는 건수의 비율. 투자자가 가진 돈에 달려 있고 우리는 정하지 않음. 금액 비중 $w$와 조금 다름\\\bottomrule
\end{tabularx}
\end{table}
\endgroup

\chapter*{요약과 문제의식}
\addcontentsline{toc}{chapter}{요약과 문제의식}
\markboth{요약과 문제의식}{}

\section*{한 쪽 요약}
\begin{keybox}[결론]
\begin{enumerate}[leftmargin=*]
\item \textbf{돈은 같은 등급 안에 있다.} LC 등급이 같아도 대출마다 질이 다르고, 그 차이를 고르면 아무거나 고를 때보다 매년 약 1\%p 더 번다(95\% 구간 0.95--1.11). 등급마다 상위 30\%를 사고 하위 30\%를 판다고 가정한 스프레드는 48개월 \textbf{전부} 양수였다(+2.38\%p).
\item \textbf{수익과 안정성은 다른 축이다.} 기대수익으로 고르면 어떤 기준(허들)에서도 가장 많이 번다. 안정성(과제 지표인 포트폴리오 Sharpe)의 1위는 투자자의 요구수익률에 따라 바뀐다. 국채 수준이면 부도 위험이 낮은 대출 위주가, 회사채(BAA) 수준이면 기대수익 순이 1위다. 국채 기준에서 보수적 선별이 이기는 것은 36개월 고위험 등급(E--G)이 위험만큼 금리를 받지 못했기 때문이다.
\item \textbf{투명한 규칙 하나가 둘 다를 잡을 가능성이 있다.} 등급마다 ``최근 24개월에 새 계좌를 적게 연'' 대출을 고르면, 네 가지 허들 모두에서 Sharpe가 가장 높고 수익은 기대수익 순과 거의 같았다. 다만 같은 데이터에서 보고 고른 규칙이라 test 전에 미리 고정한 \textbf{후보}다.
\item \textbf{재는 방법이 결론을 바꾼다.} 수익률 정의, 학습 시점, 채점 표본(끝난 대출만 보는지)을 바꾸면 같은 데이터의 결론이 뒤집혔다. 이 프로젝트의 가장 큰 기여는 모델이 아니라 \textbf{정직한 성적표}다.
\end{enumerate}
\Intuit{같은 가격표가 붙은 사과 중 더 좋은 사과를 고르면 더 번다. 비싼 사과(고위험 등급)는 가격만큼 좋지 않았다.}
\end{keybox}

\section*{문제의식: 왜 이 프로젝트가 어려운가}
과제는 ``Sharpe ratio를 최대로 만드는 대출 승인 전략''이다. 겉보기에는 부도 예측 문제다. 부도를 잘 맞히는 모델을 만들고, 부도할 것 같은 대출을 거절하면 될 것 같다. 우리는 이 직관에 숨어 있는 함정을 하나씩 확인했고, 그 과정이 곧 이 보고서의 뼈대다.

\begin{table}[H]\centering\small
\begin{tabularx}{\linewidth}{@{}p{38mm}LL@{}}\toprule
함정 & 왜 함정인가 & 우리가 한 일\\\midrule
부도만 맞히면 된다 & LC 금리는 이미 위험의 가격이다. 안전한 대출만 고르면 덜 번다 & 부도확률 $p$를 금리와 합친 \textbf{기대수익 점수}로 고른다(5장)\\
정확도(accuracy)로 평가한다 & 부도 비율이 약 15\%라, 전부 ``정상''이라 해도 85\%가 맞는다. 정확도는 돈과 무관하다 & 실제 현금흐름 수익률(IRR)로 채점한다(2장)\\
수익률은 아무렇게나 잰다 & 같은 부도 대출이 계산법에 따라 0\%, $-27\%$, $-77\%$가 되고, 결론이 뒤집힌다 & 수익률 정의 네 가지를 비교하고 IRR을 쓴다(2장)\\
데이터를 무작위로 나눠 시험한다 & 미래 대출로 과거를 맞혀 성적이 부풀려진다(AUC 0.73 vs 0.62--0.68) & 그때 알던 정보로만 학습한다(시점 학습, 3장)\\
끝난 대출만 보고 채점한다 & 빨리 끝난 대출(빨리 부도 난 대출)이 과대 대표되어 ``부도만 피하는'' 전략이 좋아 보인다 & 일부러 일찍 잘라 보는 시뮬레이션으로 편향을 측정하고 test 해석 규칙을 정했다(3장, 9장)\\
Sharpe 하나로 판단한다 & 대출 건별 Sharpe는 승인만 줄여도 변하고, 포트폴리오 Sharpe는 규모를 무시한다 & Sharpe 정의 세 가지를 모두 계산하고, 수익과 안정성 두 축으로 본다(6장)\\
비교 기준을 하나로 고정한다 & 국채 기준에서 맞는 결론이 회사채 기준에서 뒤집혔다 & 허들을 국채, +1\%p, +2\%p, BAA로 바꿔 가며 결론의 조건을 밝혔다(7장)\\
승인 대상을 자르는 단위를 신경 쓰지 않는다 & 연도 단위로 자르면 달마다 승인 금액이 들쭉날쭉해 Sharpe가 왜곡된다(외부 검증으로 발견) & 발행월 단위로 다시 계산했다(7, 8장)\\
좋은 결과는 믿는다 & 같은 데이터를 수십 번 보면 우연에 맞춰진다 & 판정 기준을 실행 전에 적고, 틀린 판단을 기록하고, test 규칙을 미리 고정했다(8, 9장)\\\bottomrule
\end{tabularx}
\caption{이 프로젝트가 넘어야 했던 함정들.}
\end{table}

\begin{warnbox}[이 보고서의 태도(11월 16일 개정)]
우리 결론 중 일부는 처음 판단과 반대로 바뀌었다. 예를 들어 ``Sharpe는 승인비율에 끌려가므로 주 지표로 부적합하다''는 팀원의 지적으로 범위를 좁혔고, ``LC 등급끼리의 가격은 공정했다''는 36개월 고위험 등급에서 틀렸다. 이런 정정은 숨기지 않고 8장에 모았다. 정정의 기록이 결과의 신뢰도를 만든다고 보기 때문이다.
\end{warnbox}

\section*{읽는 길잡이}
\begin{itemize}[leftmargin=*]
\item \textbf{1--4장(기초):} 과제와 데이터, 수익률을 재는 법, 시간을 지키는 법, EDA와 결측. 처음 보는 사람은 여기부터.
\item \textbf{5장:} 부도확률 $p$가 어떻게 만들어지고 \textbf{왜 믿을 수 있는지}. 이 보고서의 핵심 질문.
\item \textbf{6장:} 기대수익 점수와 IRR의 역할, 왜 ``아무거나 고를 때''와 비교하는지, Sharpe 정의 세 가지.
\item \textbf{7장:} 결과. 두 축, 등급 안 정보, 등급별 위험, 단순 규칙.
\item \textbf{8장:} 검증과 정정의 기록. \textbf{9장:} 한계와 test 계획.
\end{itemize}
각 핵심 박스 끝의 \textbf{직관:} 줄만 순서대로 읽어도 흐름이 잡힌다.

\chapter{과제와 데이터}

\section{과제가 실제로 묻는 것}
과제 문서는 ``Sharpe ratio를 최대로 만드는 대출 승인 전략''을 요구한다. 조교 답변에 따르면 Sharpe는 바꾸거나 다른 식을 써도 되고, accuracy는 중요하지 않으며, \textbf{과정과 논리 도출}이 평가의 중심이다. 그래서 질문을 이렇게 다시 쓴다.

\begin{keybox}[이 프로젝트의 질문]
투자자가 LendingClub(LC) 대출 중 무엇을 사야 하는가? 그리고 그 판단이 ``운이 아니라 정보 때문''이라는 것을 어떻게 보일 수 있는가?
\Intuit{예측을 잘하는 대회가 아니라, 돈이 되는 고르기를 정직하게 증명하는 숙제다.}
\end{keybox}

질문이 두 부분이라는 점이 중요하다. 앞부분(무엇을 사나)은 모델과 규칙의 문제이고, 뒷부분(어떻게 보이나)은 평가와 검증의 문제다. 이 노트의 절반 이상이 뒷부분에 쓰이는 이유다.

\section{LendingClub이라는 시장}
LC는 2007년부터 2020년까지 운영된 미국의 P2P 대출 플랫폼이다. 돈을 빌리려는 사람이 신청하면 LC가 심사해 등급(A--G, 세부등급 A1--G5 35개)과 금리를 붙이고, 투자자는 그 대출 조각(note)을 산다. 대출은 무담보 개인 신용대출이고 만기는 36개월 또는 60개월이다. 2021년 2월 Radius Bank 인수를 마치며 은행이 되었고 P2P 투자 상품은 접었다.

\begin{insight}[투자자 입장에서의 세 가지 위험]
\begin{enumerate}[leftmargin=*]
\item \textbf{부도(default)}: 갚다가 멈춘다. 원금 일부를 잃는다.
\item \textbf{조기상환(prepayment)}: 일찍 다 갚는다. 원금은 돌아오지만 높은 이자를 받을 기간이 줄어든다.
\item \textbf{재투자}: 매달 돌아오는 원리금을 다시 어딘가에 넣어야 한다. 그 수익률이 전체 성과를 바꾼다.
\end{enumerate}
\end{insight}

LC가 매긴 금리는 이미 부도 위험의 가격이다. 그래서 ``부도가 적은 대출 고르기''는 정답이 아니다. 금리가 낮은 A등급만 고르면 부도는 적지만 버는 것도 적다. 이 긴장이 프로젝트 전체를 관통한다.

\section{데이터 한눈에}
\begin{table}[H]\centering\small
\begin{tabularx}{\linewidth}{@{}lL@{}}\toprule
항목 & 값\\\midrule
공식 train & 1,755,294건 \ensuremath{\times} 141개 컬럼, 2007--2020년 발행\\
결과 상태 & Fully Paid 약 90만, Current(진행 중) 약 62만, Charged Off 약 22만, 기타 소수\\
종결 표본 & 결과가 확정된 대출 1,116,156건(LC 신용정책 미충족 대출 제외)\\
모델 특성 & 승인 시점에 알 수 있는 기본 101개(\ft{src/lc/data.py}의 \ft{feature_columns})\\
주 평가 표본 & 2013--2016년 발행 36개월 종결 대출, 발행월 48개\\
공식 test & 조교 답변: 일주일 안에 배포, 상환 결과 포함, 진행 중 대출은 없을 것\\\bottomrule
\end{tabularx}
\caption{데이터 요약.}
\end{table}

\section{컬럼을 시점으로 나누기: Data leakage}
141개 컬럼을 보면 가장 먼저 던질 질문은 하나다.

\begin{keybox}[누수 판별 질문]
``대출을 승인하던 날, 이 값을 알 수 있었나?''
알 수 없었다면 그 컬럼은 학습에 쓰면 안 된다. 쓰면 \textbf{data leakage}다.
\Intuit{시험 답안지를 보고 푼 모의고사 점수는 실력이 아니다.}
\end{keybox}

\begin{table}[H]\centering\small
\begin{tabularx}{\linewidth}{@{}p{33mm}LL@{}}\toprule
묶음 & 예 & 처리\\\midrule
승인 시점 정보 & \ft{loan_amnt}, \ft{int_rate}, \ft{grade}, \ft{annual_inc}, \ft{dti}, \ft{fico_range_low}, \ft{emp_length}, \ft{addr_state} & 특성으로 사용\\
상환 결과 & \ft{total_pymnt}, \ft{recoveries}, \ft{last_pymnt_d}, \ft{out_prncp} & 특성 금지, \textbf{수익 계산에만} 사용\\
사후 갱신 & \ft{last_fico_range_high}, \ft{hardship_*}, \ft{settlement_*}, \ft{debt_settlement_flag} & 삭제(\ft{LEAK}, \ft{LEAK_PREFIX})\\
식별, 설명 & \ft{id}, \ft{url}, \ft{desc}, \ft{title} & 삭제\\\bottomrule
\end{tabularx}
\caption{컬럼 묶음. 실제 목록은 \ft{src/lc/data.py}.}
\end{table}

\begin{warnbox}[미묘한 누수]
\ft{last_fico_range_high}는 이름만 보면 신용점수지만, 대출 후 마지막으로 갱신된 점수다. 부도가 나면 점수가 떨어지므로 결과를 거의 그대로 알려 준다. 이런 컬럼 하나가 AUC를 0.9 넘게 올리는데, 그 점수는 아무 의미가 없다.
\end{warnbox}

\section{결과 라벨 만들기}
\ft{data.parse}는 원본에 다음 열을 붙인다.
\begin{itemize}[leftmargin=*]
\item \ft{outcome}: Fully Paid, Charged Off 등 상태를 정리한 값. \ft{resolved}: 결과가 확정되었는가.
\item \ft{bad}: 부도면 1. Charged Off와 Default를 부도로 본다.
\item \ft{r}: 약정 금리(소수), \ft{term_m}: 만기 개월, \ft{tau}: 발행부터 마지막 납부까지 개월 수.
\item \ft{policy_fail}: ``Does not meet the credit policy''로 표시된 옛 대출. 심사 기준이 달라 분석에서 뺀다.
\end{itemize}

\begin{career}
``누수를 어떻게 막았나요?''는 면접 단골 질문이다. 답의 핵심은 목록이 아니라 \textbf{판별 기준}이다. ``각 컬럼에 대해 결정 시점에 관측 가능한지를 물었고, 결과 컬럼은 라벨과 수익 계산에만 쓰고 특성에서는 뺐다''라고 말하면 된다. 이 프로젝트에서는 한 단계 더 나아가, 컬럼뿐 아니라 \textbf{행(어떤 대출을 학습에 넣나)}도 시점을 지켰다(3장).
\end{career}

\chapter{수익률을 어떻게 잴 것인가}

과제가 Sharpe를 말하는 순간 ``대출 하나의 수익률''이 필요해진다. 그런데 데이터에는 수익률 컬럼이 없다. 있는 것은 약정 금리, 월 상환액, 받은 총액, 마지막 납부일뿐이다. 이 장은 그것들로 수익률을 만드는 법과, 정의를 어떻게 고르느냐에 따라 결론이 뒤집히는 이유를 다룬다.

\section{원리금 균등 상환}
\textbf{질문.} 1만 달러를 연 12\%로 36개월 동안 빌리면 매달 얼마를 갚나?

매달 같은 금액 $A$를 내는데, 그 36번의 납입을 오늘 가치로 할인해 더하면 빌린 돈 $P$가 되어야 한다. 월 이자율 $i=r/12$로 할인하면
\[
P=\sum_{t=1}^{T}\frac{A}{(1+i)^t}=A\cdot\frac{1-(1+i)^{-T}}{i}
\quad\Longrightarrow\quad
A=P\,\frac{i}{1-(1+i)^{-T}}.
\]

\begin{workout}[월 상환액]
\Use{원리금 균등 상환식(이 절), 등비수열의 합}
\Step{1} 월 이자율: $i=0.12/12=0.01$. \Concept{연 금리를 월로}
\Step{2} 할인 인자: $(1.01)^{-36}=0.6989$. 따라서 $1-(1.01)^{-36}=0.3011$.
\Step{3} $A=10000\times 0.01/0.3011=332.14$. \Concept{원리금 균등 상환}
\Ans{매달 \$332.14. 데이터의 \ft{installment}가 바로 이 값이다.}
\Verify{$332.14\times 36=11{,}957$. 이자 총액 약 \$1,957로, 원금이 줄어드는 만큼 12\% \ensuremath{\times} 3년(\$3,600)보다 작다.}
\Trap{연 금리 12\%를 그대로 36으로 곱해 이자를 계산한다. 원금이 매달 줄어드는 것을 잊은 것이다.}
\end{workout}

\section{IRR: 현금흐름이 실제로 낸 수익률}
\textbf{질문.} 대출이 중간에 부도 나면 수익률은 몇인가?

받은 돈의 흐름을 $CF_0=-P$(투자), $CF_1,\dots,CF_H$(받은 돈)로 적자. 이 흐름의 현재가치를 0으로 만드는 월 할인율 $m$이 있으면, 그것이 이 투자가 실제로 낸 월 수익률이다.
\[
\sum_{t=0}^{H}\frac{CF_t}{(1+m)^t}=0,\qquad \mathrm{IRR}=(1+m)^{12}-1.
\]
식이 다항식이라 닫힌 해가 없으므로 수치로 푼다. 코드는 \ft{returns.irr}이며 Brent 방법(\ft{brentq})으로 근을 찾는다.

\begin{keybox}[IRR]
돈을 넣고 돌려받은 흐름 전체를 하나의 ``연 이자율''로 요약한 값이다. 정상 상환이면 약정 금리와 같고, 부도 나면 그보다 낮다.
\Intuit{이 통장에 넣었다 뺀 돈의 흐름을, 매년 몇 \% 이자 통장과 같다고 볼 수 있나?}
\end{keybox}

\begin{workout}[정상 상환 대출의 IRR]
\Use{IRR 정의(이 절), 월 복리}
\Step{1} 현금흐름: $CF_0=-10000$, $CF_1=\dots=CF_{36}=332.14$.
\Step{2} $m=0.01$을 넣으면 \S2.1의 식에 의해 현재가치가 정확히 0이다. 따라서 월 IRR은 1\%. \Concept{IRR 정의}
\Step{3} 연환산: $(1.01)^{12}-1=0.1268$. \Concept{월 복리}
\Ans{연 IRR 12.68\%. 약정 금리 12\%는 단리식 표기(APR)이고, IRR은 복리라 조금 크다.}
\Verify{\ft{checks.py} 2번: 월 0.01000, 연 0.12683.}
\end{workout}

\begin{workout}[10개월 내고 부도 난 대출]
\Use{현금흐름 복원(아래), IRR, 손실률}
\Step{1} 상황(예시): 10번 냈고, 15개월째 회수금 \$500이 들어왔다. 받은 총액 $=10\times332.14+500=3{,}821.43$.
\Step{2} 손실률 $=1-3821.43/10000=0.618$. \Concept{LGD 정의, \S5.6}
\Step{3} 현금흐름 $CF_0=-10000$, $CF_{1..10}=332.14$, $CF_{15}=500$의 IRR을 수치로 풀면 연 $-76.8\%$.
\Step{4} 같은 대출의 단순 연환산(우승 논문식): $(3821.43/10000)^{12/36}-1=-27.4\%$.
\Ans{IRR $-76.8\%$, 단순 연환산 $-27.4\%$. 같은 대출인데 정의에 따라 세 배 가까이 다르다.}
\Verify{\ft{checks.py} 3번. 손실률 0.618은 우연히 2013년 학습 표본의 실제 평균 LGD(0.618)와 같다.}
\Trap{부도 대출의 IRR이 $-100\%$에 가까울 것이라 생각한다. 이미 받은 원리금이 있어 그보다 덜 나쁘다. 반대로 단순 연환산은 손실을 3년에 걸쳐 나눠 덜 나빠 보이게 만든다.}
\end{workout}

\paragraph{현금흐름 복원.} 데이터에는 월별 납입 기록이 없다. 그래서 \ft{returns.schedule}은 이렇게 근사한다: 받은 총액을 월 상환액으로 나눈 몫 $k_f$만큼은 매달 정상 납입했고, 나머지 $\text{rem}=\text{받은 총액}-k_f A$는 마지막 납부 달($\tau$)에 한꺼번에 받았다. 조기상환은 큰 rem이 일찍 들어오는 것으로, 부도 후 회수는 작은 rem으로 표현된다. 이 근사는 연구 노트의 가정 목록에 올라 있다.

\section{포트폴리오 IRR과 수수료}
투자자는 대출 하나가 아니라 여러 개를 산다. 같은 달에 발행된 대출들의 현금흐름을 더해 한 포트폴리오로 보고 IRR을 구한다(\ft{returns.pooled_irr}). LC는 투자자가 받은 돈에서 1\% 서비스 수수료를 뗀다.

\begin{workout}[정상 9건 + 부도 1건]
\Use{pooled IRR(이 절), 수수료}
\Step{1} 각 \$10,000, 연 12\%, 36개월. 9건은 정상 상환, 1건은 앞 예제처럼 부도.
\Step{2} 월별 합계: 1--10개월 $10\times332.14$, 11--36개월 $9\times332.14$, 15개월에 +500. 0개월 $-100{,}000$.
\Step{3} 받은 돈에 0.99를 곱한다. \Concept{투자자 수수료 1\%}
\Step{4} IRR을 풀면 연 6.85\%.
\Ans{부도율 10\%가 연 12.68\%의 수익을 6.85\%로 깎았다.}
\Verify{\ft{checks.py} 4번.}
\end{workout}

\section{네 가지 수익률 정의와 결론의 뒤집힘}
\begin{table}[H]\centering\small
\begin{tabularx}{\linewidth}{@{}p{34mm}LL@{}}\toprule
정의 & 계산 & 문제점\\\midrule
수업식 & 정상이면 약정 금리, 부도면 0 & 부도 손실(원금)을 무시해 과대평가\\
단순 연환산 & (받은/대출)$^{12/T}-1$ & 돈이 일찍 돌아온 것(재투자)을 무시, 조기상환을 손해로 셈\\
재투자 & 받은 돈을 국채 금리로 만기까지 굴림 & 재투자 금리 가정 필요\\
코호트 IRR & 발행월 포트폴리오의 현금흐름 IRR, 수수료 1\% & 현금흐름 시점이 근사\\\bottomrule
\end{tabularx}
\caption{수익률 정의 네 가지(\ft{src/lc/returns.py}). 이 프로젝트의 주 정의는 마지막 줄이다.}
\end{table}

\begin{keybox}[결과 C1--C3]
2013--16년 36개월 대출을 \textbf{전부 승인}하면 코호트 IRR 5.73\%로 같은 기간 3년 국채(0.87\%)를 48개 발행월 모두에서 이긴다(C2). 지난 우승 논문은 ``전부 승인 \ensuremath{<} 국채''라고 결론냈는데, 이는 단순 연환산과 최근 발행분의 종결 표본에서 비롯된 것으로 보인다(C3, 잠정).
\Intuit{같은 데이터라도 자를 대는 방법이 다르면 다른 답이 나온다. 그래서 자부터 정해야 한다.}
\end{keybox}

\section{비교 기준: 무위험수익률은 무엇인가}
Sharpe와 초과수익에는 ``아무 위험 없이 벌 수 있는 수익'' $r_f$가 필요하다.
\begin{itemize}[leftmargin=*]
\item \textbf{만기를 맞춘 국채}(FRED GS3, GS5): 36개월 대출에는 3년 국채, 60개월에는 5년 국채를 발행월 기준으로 붙인다(\ft{returns.attach_treasury}). 투자 기간이 같은 대안이라 가장 정직하다.
\item \textbf{3개월 국채}(TB3MS): 교과서의 $r_f$. 만기가 달라 기간 프리미엄이 섞인다.
\item \textbf{회사채 BAA}: 무위험은 아니지만 ``비슷한 위험의 다른 투자''라는 실무적 비교 기준. BAA로 바꾸면 전부 승인의 우위가 사라지고 선별의 가치가 커진다(C8).
\end{itemize}

\begin{insight}
실무에서 대출 포트폴리오의 성과는 조달 비용(은행이면 예금 금리 + 스프레드) 또는 같은 위험의 벤치마크 대비로 본다. ``무위험''은 출발점일 뿐이라, 이 프로젝트는 국채를 주 기준으로 두고 BAA와 +1--3\%p 프리미엄을 민감도로 함께 보고한다.
\end{insight}

\begin{career}
``Accuracy가 아니라 무엇을 최적화했나?''라는 질문에 이 장이 답이다. 라벨(부도 0/1)이 아니라 \textbf{현금흐름에서 복원한 IRR}을 목표로 삼았고, 정의 선택이 결론을 뒤집는다는 것을 수치로 보였다.
\end{career}

\chapter{시간을 지키는 법: 시점 학습과 편향}

2장에서 수익률의 자를 정했다. 이 장은 그 자로 \textbf{언제의 정보로} 학습하고 \textbf{어떤 대출로} 채점하는지를 다룬다. 누수는 컬럼(1장)에서만 생기지 않는다. 행에서도 생긴다.

\section{무작위 분할은 왜 성적을 부풀리나}
보통 머신러닝 수업에서는 데이터를 무작위로 80:20으로 나눈다. LC 데이터에 그렇게 하면 2016년 대출로 학습해 2013년 대출을 맞히는 일이 생긴다. 2016년의 경기, LC의 심사 기준, 금리 수준은 2013년 투자자가 알 수 없던 정보다.

\begin{keybox}[결과 C5]
무작위 분할 AUC 약 0.73, 시점 학습 AUC 0.62--0.68. 차이 전부가 ``미래를 본 값''이다.
\Intuit{내년 기출문제로 공부하고 올해 시험을 보면 점수는 높지만, 실제로는 그런 일이 불가능하다.}
\end{keybox}

\section{시점 학습(as-of training)}
연도 $Y$에 발행된 대출을 평가할 때, 모델은 \textbf{$Y$년 1월 1일에 이미 결과가 확정된 대출}로만 학습한다(\ft{asof.training_set}). 확정 조건은 두 가지다.

\begin{enumerate}[leftmargin=*]
\item 그 날짜 전에 발행되었다.
\item 그 날짜까지 결과가 났다. 정상 상환은 마지막 납부 달에, 부도는 \textbf{마지막 납부 + 5개월}에 확정된 것으로 본다(\ft{CHARGE_OFF_LAG = 5}).
\end{enumerate}

\begin{figure}[H]\centering
\begin{tikzpicture}[x=1.0cm,y=0.75cm,font=\small]
\draw[-{Stealth}] (0,0)--(9.6,0) node[right]{시간};
\foreach \x/\l in {1/2010,3/2011,5/2012,7/2013,9/2014} {\draw (\x,0.1)--(\x,-0.1) node[below]{\l};}
\draw[very thick] (7,-0.6)--(7,4.3) node[above]{기준일 2013-01-01};
\draw[black!70,line width=2pt] (1.2,1)--(4.5,1); \node[right] at (9.8,1) {A: 정상 상환, 확정};
\draw[black!70,line width=2pt] (2,2)--(5.4,2); \draw[black!35,line width=2pt,dashed] (5.4,2)--(6.4,2); \node[right] at (9.8,2) {B: 부도, 5개월 뒤 확정};
\draw[black!70,line width=2pt] (4,3)--(6.6,3); \draw[black!35,line width=2pt,dashed] (6.6,3)--(7.6,3); \node[right] at (9.8,3) {C: 기준일 뒤 확정, 제외};
\draw[black!70,line width=2pt] (5.5,4)--(9.5,4); \node[right] at (9.8,4) {D: 진행 중, 제외};
\end{tikzpicture}
\caption{시점 학습의 학습 표본. 2013년 대출을 평가할 모델은 A와 B로만 학습한다. 점선은 상각 지연 5개월.}
\end{figure}

\begin{whybox}[왜 5개월인가]
LC는 연체 120일(약 4개월)이 넘으면 상각(charge-off)한다. 마지막 납부 이후 상각이 장부에 잡히기까지 약 5개월이 걸린다. 지연을 0으로 두면 아직 ``연체 중''이던 대출을 이미 부도로 알았던 것처럼 학습해 미래 정보가 섞인다. 6단계에서 지연 0과 5를 비교했고, 가정에 따라 편향 방향이 바뀌는 것을 확인했다.
\end{whybox}

\paragraph{대가.} 시점 학습은 표본이 작다. 2013년 모델의 학습 대출은 11,345건이고, 2016년 모델은 156,153건이다. 이른 해일수록 모델이 약하고 불안정하다(5장의 2013년 결과들이 흔들리는 이유).

\section{종결 표본 편향(resolution bias)}
\textbf{질문.} 평가할 때 ``이미 결과가 난 대출''만 쓰면 무엇이 달라지나?

\begin{workout}[4건짜리 종결 편향]
\Use{censoring, 종결 표본 편향(이 절)}
\Step{1} 같은 달 발행된 36개월 대출 4건(예시):
A 저금리 6\%, 3년 내내 정상. B 고금리 20\%, 3년 내내 정상. C 고금리 20\%, 6개월 뒤 부도. D 저금리 6\%, 1년 안에 조기상환.
\Step{2} 3년 뒤 전체를 보면 부도율 $1/4=25\%$. 고금리 B의 이자가 C의 손실을 메운다.
\Step{3} 발행 1년 뒤에 ``끝난 대출만'' 보면 C와 D만 남는다. 부도율 $1/2=50\%$. \Concept{censoring}
\Step{4} 남은 표본에서 ``고금리를 고르는 규칙''은 C만 고르게 되어 최악이고, ``부도확률이 낮은 것을 고르는 규칙''은 D를 골라 최선이다.
\Ans{일찍 본 표본은 부도를 과대 대표하고, 오래 이자를 받아야 이득인 고금리 정상 대출(B)을 숨긴다.}
\Trap{``종결 대출만 쓰면 깨끗하다''고 생각한다. 깨끗한 것은 라벨이고, 표본 구성은 오히려 왜곡된다.}
\end{workout}

\subsection*{실제 데이터로 시뮬레이션(26단계)}
2013--15년 36개월 대출은 지금 모두 끝나 정답을 안다. 같은 대출을 관측 시점에 맞춰 일부러 잘라, 두 순위 방식의 선별 가치(6장)를 비교했다. LightGBM 시드 10개, 발행월 부트스트랩 1만 번, 판정 기준은 실행 전에 고정했다.

\begin{table}[H]\centering\small
\begin{tabular}{@{}lrrrrr@{}}\toprule
관측 & 남은 비율 & 부도율 & 기대수익 순위 & 부도확률 순위 & 차이(95\% 구간)\\\midrule
발행 후 1--2년 & 약 30\% & 약 18\% & 1.44\%p & 5.28\%p & $-3.83$ ($-4.67$, $-3.04$)\\
발행 후 2--3년 & 약 59\% & 약 18\% & 1.18\%p & 2.21\%p & $-1.03$ ($-1.36$, $-0.70$)\\
전체 & 100\% & 약 14\% & 1.01\%p & 0.10\%p & $+0.91$ ($+0.73$, $+1.10$)\\\bottomrule
\end{tabular}
\caption{종결 편향 시뮬레이션(\ft{scripts/26_resolution_bias.py}). 단위는 무작위 승인 대비 \%p.}
\end{table}

\begin{keybox}[결과 C23과 C6 정정]
관측 기간이 짧을수록 부도확률 순위가 이기는 것처럼 보이고, 전체를 보면 기대수익 순위가 이긴다. 14단계의 2017년 보류 구간에서 부도확률 순위가 이긴 것은 이 편향으로 설명된다.
\Intuit{일찍 채점하면 ``빨리 망한 대출''만 보이므로, 망하지 않는 것만 신경 쓰는 규칙이 똑똑해 보인다.}
\end{keybox}

\begin{warnbox}[고칠 수 있나]
종결 여부는 결과(부도, 조기상환) 자체에 달려 있다. 이런 결측은 MNAR(4장)이라, 특성으로 만든 가중치(inverse probability weighting)로는 고칠 수 없다. 대신 \textbf{test를 보기 전에 읽는 법을 정했다}: 발행연도별 만기 도달 비율부터 확인하고, 비율이 높은 연도만 두 순위 비교의 근거로 쓴다.
\end{warnbox}

\section{생존분석은 왜 안 썼나}
진행 중 대출 62만 건을 버리는 것이 아깝다면, 생존분석(survival analysis)으로 ``아직 안 끝난'' 정보까지 쓸 수 있다. 시도했지만 단순 PD 모델을 이기지 못했다(연구 노트 실패 목록). 수준(확률 크기)에는 편향이 있어도 순위는 크게 무너지지 않았기 때문이다.

\begin{career}
시계열이 섞인 데이터에서 무작위 분할을 쓰지 않은 이유, 그리고 \textbf{평가 표본의 선택 편향}까지 시뮬레이션으로 정량화한 것은 강한 이야깃거리다. ``모델 성능을 올렸다''보다 ``성능이 부풀려지는 경로 두 가지(무작위 분할, 종결 편향)를 찾아 크기를 쟀다''가 훨씬 드물다.
\end{career}

\chapter{EDA, 피처, 결측}

\section{EDA는 무엇을 위한 것인가}
EDA(exploratory data analysis)는 그림을 많이 그리는 일이 아니라, \textbf{모델링 결정을 내리기 위한 질문에 답하는 일}이다. 이 프로젝트의 EDA 질문은 네 개였다.
\begin{enumerate}[leftmargin=*]
\item 어떤 변수가 혼자서 부도를 가르는가? \ensuremath{\to} IV(\S4.2)
\item 어떤 변수들이 서로 같은 정보인가? \ensuremath{\to} 순위상관
\item LC 등급이 이미 설명한 것 말고 남은 정보가 있는가? \ensuremath{\to} 등급 안 분석(\S4.3)
\item 값이 왜 비어 있나? \ensuremath{\to} 결측 메커니즘(\S4.4)
\end{enumerate}
12단계 EDA는 평가 구간을 보지 않도록 2007--12년 대출만 썼다. 20단계 EDA(\ft{results/eda/})는 2007--16년 종결 대출을 썼다.

\section{WOE와 IV: 한 변수의 힘}
\textbf{질문.} 연소득이 부도와 얼마나 관련 있는지를 한 숫자로 말할 수 있을까?

변수를 구간(bin)으로 나누고, 각 구간에 정상 대출과 부도 대출이 얼마나 몰렸는지 비교한다. 전체 정상 $G$건 중 구간 $k$에 $g_k$건, 전체 부도 $B$건 중 $b_k$건이 있으면
\[
\mathrm{WOE}_k=\ln\frac{g_k/G}{b_k/B},\qquad
\mathrm{IV}=\sum_k\Big(\frac{g_k}{G}-\frac{b_k}{B}\Big)\mathrm{WOE}_k .
\]
WOE가 양수면 그 구간은 정상 쪽, 음수면 부도 쪽 증거다. IV는 모든 구간의 ``치우침 \ensuremath{\times} 증거''를 더한 것이라 항상 0 이상이다. 관행상 0.02 미만은 약함, 0.1--0.3 중간, 0.3 이상 강함이다.

\begin{workout}[결측 구간의 WOE]
\Use{WOE, IV(이 절), MNAR(\S4.4)}
\Step{1} 예시: 정상 900건, 부도 100건. 어떤 변수가 결측인 구간과 값이 있는 구간 두 개.
\Step{2} 경우 1: 결측 구간에 정상 270, 부도 30. 비중이 정상 $270/900=0.3$, 부도 $30/100=0.3$로 같다.
$\mathrm{WOE}=\ln(0.3/0.3)=0$, 값 있는 구간도 0. IV $=0$. \Concept{WOE 정의}
\Step{3} 경우 2: 결측 구간에 정상 240, 부도 60. 비중 $0.267$ vs $0.6$.
$\mathrm{WOE}_{\text{결측}}=\ln(0.267/0.6)=-0.811$, $\mathrm{WOE}_{\text{있음}}=\ln(0.733/0.4)=+0.606$.
\Step{4} $\mathrm{IV}=(0.267-0.6)(-0.811)+(0.733-0.4)(0.606)=0.472$. \Concept{IV 정의}
\Ans{경우 1은 결측이 아무 정보도 없고(IV 0), 경우 2는 ``비어 있다는 사실'' 자체가 강한 부도 신호다(IV 0.47).}
\Verify{\ft{checks.py} 8번.}
\Trap{결측을 평균으로 채워 버린다. 경우 2라면 가장 강한 신호를 지운 것이다.}
\end{workout}

\begin{table}[H]\centering\small
\begin{tabular}{@{}lr|lr@{}}\toprule
변수 & IV & 변수 & IV\\\midrule
\ft{sub_grade} & 0.487 & \ft{term} & 0.178\\
\ft{grade} & 0.452 & FICO & 0.112\\
\ft{int_rate} & 0.441 & \ft{annual_inc} & 0.029\\\bottomrule
\end{tabular}
\caption{IV 상위(20단계). 101개 중 67개는 혼자서는 IV 0.02 미만이다.}
\end{table}

등급, 세부등급, 금리가 압도적이다. 셋은 사실상 같은 정보(LC의 판단)다. 순위상관 0.9 이상인 쌍이 12개 있었다(FICO 상한과 하한 1.0, 대출액과 승인액 1.0 등).

\section{등급 안에 남은 정보}
IV만 보면 연소득은 약한 변수(0.029)다. 그런데 \textbf{같은 세부등급 안에서} 비교하면 이야기가 달라진다.

\begin{keybox}[결과 C14, C18]
등급이 설명하지 못한 정보는 소득, 상환 부담(월 상환액/월 소득), 최근 계좌 개설, 부채비율(DTI), 지역에 남아 있다. 같은 세부등급, 같은 금리에서 소득 상위이면서 DTI 하위 20\%인 대출이 IRR이 약 0.9\%p 높다(모델 없이 계산).
\Intuit{LC는 신용점수는 가격에 잘 넣었지만, ``갚을 능력''은 덜 넣었다.}
\end{keybox}

혼자서는 약하고 등급과 함께 볼 때 힘이 나는 변수가 있다는 것, 이것이 후보 A(같은 등급 안 선별)의 출발점이다. 단, 소득 효과는 해마다 줄었다(2.66\%p에서 0.19\%p). 시장이 그 정보를 점점 가격에 넣었을 수 있다.

\section{결측 메커니즘}
\textbf{질문.} 값이 비어 있는 이유가 결과와 관련 있나?

\begin{table}[H]\centering\small
\begin{tabularx}{\linewidth}{@{}p{30mm}LL@{}}\toprule
메커니즘 & 뜻 & LC 예\\\midrule
MCAR & 완전히 무작위로 비어 있음 & 거의 없음(소량 결측 일부)\\
MAR & 관측된 다른 변수(예: 발행 시기)로 설명됨 & 2015년 말 이후 생긴 새 변수 14개(\ft{il_util} 등)\\
MNAR & 비어 있는 값 자체와 관련 있음 & \ft{emp_length}: 무직이거나 밝히기 꺼리면 비움\\
구조적 & 정의상 값이 없음 & 공동신청 변수 15개(단독 신청이면 없음)\\
사건 없음 & 그 사건이 없었음 & \ft{mths_since_last_delinq}: 연체한 적이 없으면 빔\\\bottomrule
\end{tabularx}
\caption{결측 유형(\ft{src/lc/missing.py}의 \ft{classify}).}
\end{table}

\begin{whybox}[사건 없음은 ``아주 오래전''이다]
``마지막 연체 후 몇 개월''이 비었다는 것은 연체한 적이 없다는 뜻이다. 평균(예: 30개월)으로 채우면 ``30개월 전에 연체했다''는 거짓 정보가 된다. 그래서 선형 모델용으로는 학습 데이터의 최댓값으로 채우고 결측 표시를 붙이는 방식을 시험했다(\ft{typed_linear}).
\end{whybox}

\section{결측 처리 결정(22단계)}
결측 처리는 모델 비교보다 \textbf{먼저} 정해야 한다. 그렇지 않으면 ``모델 A가 낫다''가 사실은 ``모델 A에 맞춘 결측 처리가 낫다''일 수 있다. 판정 기준은 실험 전에 정했다: 가장 단순한 방식보다 선별 가치 0.1\%p 초과 또는 AUC 0.002 초과로 나은 연도가 4개 중 3개 이상이면 채택.

\begin{itemize}[leftmargin=*]
\item \textbf{트리 모델}(LightGBM): 결측 그대로 / 중앙값 / 중앙값 + 결측 표시 / 유형별 처리 모두 시드 차이 안. \ensuremath{\to} \textbf{결측 그대로}. LightGBM은 분할마다 결측을 왼쪽, 오른쪽 중 나은 쪽으로 보내는 법을 스스로 배운다.
\item \textbf{선형 모델}(로지스틱): 1차 비교에서 2016년 AUC 0.508로 무너졌다.
\end{itemize}

\begin{warnbox}[로지스틱 붕괴의 진짜 원인]
결측 처리가 문제가 아니었다. 2016년 모델의 학습 기간(2015년까지)에는 새 변수 33개가 거의 비어 있었다. 학습에서 그 변수의 표준편차는 소수 대출로 계산되어 매우 작았고, 평가 연도에는 모든 대출에 값이 있어 표준화 값이 최대 290까지 튀었다. 선형 모델은 큰 입력에 계수를 그대로 곱하므로 예측이 폭발한다.
\end{warnbox}

\begin{keybox}[확정 방침]
\begin{itemize}[leftmargin=*]
\item 트리: 결측 그대로.
\item 선형: \textbf{가용성 규칙}(학습 가용률 50\% 미만 변수 제외) \ensuremath{\to} 1--99\% 분위 자르기 \ensuremath{\to} 중앙값 대치. \ft{emp_length} 결측은 ``모름'' 범주로.
\end{itemize}
두 규칙을 넣자 2016년 AUC가 0.678로 정상화됐다. 이후 비교에서는 중앙값보다 확실히 나은 방식이 없었다.
\Intuit{결측을 어떻게 채우느냐보다, 학습 때 거의 없던 변수를 쓰지 않는 것이 훨씬 중요했다.}
\end{keybox}

\section{Feature engineering}
만든 피처 6종: 상환 부담(PTI = 월 상환액/월 소득), 대출 대비 소득(LTI), 신용 이력 개월, 근속 연수 숫자화, 우편번호 3자리 부도율(out-of-fold target encoding), FICO 평균. 대조군으로 노이즈 피처와 시드 변경을 함께 넣었다.

\begin{keybox}[결과 C15, C16, C21]
LightGBM에서는 모든 피처의 효과($-0.01$--$+0.02$\%p)가 시드만 바꿔도 생기는 차이(최대 0.09\%p)보다 작았다. 트리는 원본 변수(소득, 상환액, 대출액)에서 이미 비율 관계를 잡는다. 선형 모델에서는 피처가 조금 도움이 된다(L2 \ensuremath{\to} L2 + 피처, Lasso \ensuremath{\to} Lasso + 피처).
\Intuit{나무는 나눗셈을 스스로 흉내 내지만, 직선 모델은 누가 나눠 줘야 한다.}
\end{keybox}

\begin{prof}[target encoding의 누수]
우편번호별 부도율을 그 대출 자신이 포함된 데이터로 계산하면 자기 라벨이 특성에 섞인다. 학습 데이터를 5조각으로 나눠 각 조각의 값은 나머지 4조각으로 계산(out-of-fold)하고, 표본이 적은 우편번호는 전체 평균 쪽으로 당겼다(smoothing).
\end{prof}

\begin{career}
``Feature engineering으로 성능을 얼마나 올렸나?''에 ``0이었다, 그리고 그걸 어떻게 확인했나''로 답할 수 있다. 노이즈 피처와 시드 변경을 \textbf{대조군}으로 넣어 효과 크기를 잡음과 비교한 것이 핵심이다.
\end{career}

\chapter{부도확률 $p$는 어떻게 만들어지고, 왜 믿을 수 있나}

이 장은 팀에서 가장 자주 나온 질문에 답한다. ``$p$가 모델에서 나온다는데, 정확히 어떻게 계산되나? 그 숫자를 왜 믿어야 하나?'' 답을 먼저 쓰면 이렇다.

\begin{keybox}[이 장의 답]
\begin{itemize}[leftmargin=*]
\item $p$는 ``이 대출이 만기까지 가는 동안 결국 부도로 끝날 확률''이다. LightGBM이라는 모델이 과거 대출의 결과를 보고 배운 규칙으로 계산한다.
\item 모델은 수백 개의 작은 결정나무를 차례로 쌓는다. 각 나무는 ``앞 나무들이 틀린 만큼''을 고치고, 나무들의 점수를 모두 더한 뒤 S자 함수로 0--1 사이 확률로 바꾼 것이 $p$다.
\item $p$의 \textbf{순서}는 믿을 수 있다. $p$를 10등분하면 실제 부도율이 3.8\%에서 30.6\%까지 차례로 오르고, 같은 등급 안에서도 그렇다. 우연이 아니라는 검증(정답 섞기, 시드, 다른 모델)도 통과했다.
\item $p$의 \textbf{크기}는 그대로 믿으면 안 된다. 실제보다 낮게 나오는 경향이 있다. 그래서 우리는 $p$를 순위를 매기는 데 쓰고, 그 순위가 실제로 돈을 버는지는 IRR로 따로 채점한다.
\end{itemize}
\Intuit{$p$는 정확한 온도계가 아니라 믿을 만한 줄 세우기다. 누가 더 위험한지는 잘 맞히고, 몇 \%인지는 대략 맞힌다.}
\end{keybox}

\section{$p$는 무엇인가}
\paragraph{정답(라벨).} 이미 끝난 대출에 대해 부도(Charged Off, Default)면 1, 다 갚았으면(Fully Paid) 0이다. 진행 중인 대출은 결과를 모르므로 학습에도 채점에도 쓰지 않는다(3장).

\paragraph{입력(특성).} 대출을 승인하던 날 알 수 있던 정보 101개다. 금리, 등급, 대출액, 소득, 부채비율(DTI), FICO, 근속 연수, 거주 주, 최근 계좌 개설 수, 신용 이력 길이 등이다. 결과를 담은 컬럼(받은 총액, 마지막 납부일, 대출 후 갱신된 신용점수)은 쓰지 않는다(1장).

\paragraph{뜻.} $p = 0.15$는 ``이 대출과 비슷한 대출 100개 중 약 15개가 결국 부도났다''는 뜻이다. 1년 안의 부도확률이 아니라 대출 생애 전체 기준이다.

\section{가장 단순한 모델부터: 비율과 로지스틱}
가장 단순한 $p$는 ``학습 데이터의 부도 비율''이다. 2016년 모델의 학습 데이터에서는 약 17\%였다. 모든 대출에 17\%를 주면 아무것도 구분하지 못한다.

다음 단계는 로지스틱 회귀다. 특성의 가중합 $z=\beta_0+\beta_1x_1+\dots$를 만들고 $p=1/(1+e^{-z})$로 바꾼다. $z$는 ``로그 오즈'' $\ln\frac{p}{1-p}$이며, 계수 $\beta_j$는 ``$x_j$가 1 늘면 로그 오즈가 $\beta_j$ 는다''로 읽는다. 문제는 직선 하나로 모든 관계를 표현해야 한다는 점이다. ``소득이 낮으면서 부채비율도 높을 때만 위험이 급증한다'' 같은 조합은 잘 잡지 못한다.

\section{결정나무: 질문으로 나누기}
결정나무는 ``FICO(미국 개인 신용점수)가 690보다 낮은가?''(기준점은 설명용 예시이고, 실제 기준점은 모델이 데이터에서 찾는다) 같은 질문으로 대출을 둘로 나누고, 각 칸의 부도율을 예측으로 쓴다. 질문을 여러 번 하면 조합 효과를 자연스럽게 잡는다. 결측값도 ``값이 없는 쪽''으로 보내는 법을 스스로 배운다(4장). 하지만 나무 하나는 학습 데이터의 우연에 민감해서 불안정하다.

\section{그래디언트 부스팅: 틀린 만큼 고치며 쌓기}
LightGBM은 나무를 \textbf{차례로} 쌓는다. 각 나무는 처음부터 다시 배우지 않고, 앞 나무들이 아직 틀리고 있는 부분만 고친다. 무엇이 ``틀린 부분''인지는 잔차 $y-p$로 잰다. 부도 났는데($y=1$) $p$가 낮으면 잔차가 크게 양수이고, 그 대출의 점수를 올려야 한다는 뜻이다.

\begin{workout}[대출 6건으로 부스팅 두 번]
\Use{로그 오즈, 잔차, 뉴턴 단계(아래 설명), 학습률}
\Step{1} 예시 데이터(숫자는 예시): 대출 6건, 특성 두 개(FICO 낮음 여부, 최근 계좌 개설 많음 여부), 부도 2건.
\begin{center}\small
\begin{tabular}{@{}lcccccc@{}}\toprule
대출 & 1 & 2 & 3 & 4 & 5 & 6\\\midrule
FICO 낮음 & 1 & 1 & 1 & 0 & 0 & 0\\
계좌 개설 많음 & 1 & 0 & 1 & 1 & 0 & 0\\
부도 $y$ & 1 & 0 & 1 & 0 & 0 & 0\\\bottomrule
\end{tabular}\end{center}
\Step{2} 출발점: 전체 부도율 $2/6$의 로그 오즈 $\ln(1/2)=-0.693$을 모든 대출의 점수 $z$로 둔다. $p=0.333$. 로그 손실 0.637. \Concept{로그 오즈}
\Step{3} 나무 1(질문: FICO 낮음?). 잔차 $y-p$는 부도 대출 $+0.667$, 정상 대출 $-0.333$. ``FICO 낮음'' 칸(대출 1--3)의 잔차 합 $+1.0$, 곡률 합 $\sum p(1-p)=0.667$이므로 잎 값 $1.0/0.667=+1.5$. 반대 칸은 $-1.5$. \Concept{뉴턴 단계}
\Step{4} 학습률 0.5를 곱해 점수에 더한다: FICO 낮음 칸 $z=-0.693+0.75=0.057$, 나머지 $-1.443$. 확률은 0.514와 0.191, 로그 손실 0.448.
\Step{5} 나무 2(질문: 계좌 개설 많음?). 새 잔차로 같은 계산: ``많음'' 칸 잎 값 $+1.193$, ``적음'' 칸 $-1.604$. 학습률을 곱해 $+0.597$, $-0.802$를 더한다.
\Step{6} 최종 $p$: 대출 1·\allowbreak{}3(부도) 0.658, 대출 2 0.322, 대출 4 0.300, 대출 5·\allowbreak{}6 0.096. 로그 손실 0.297.
\Ans{나무 두 개만으로 부도 대출의 $p$가 0.33에서 0.66으로 오르고, 안전한 대출은 0.10으로 내려갔다. 실제 모델은 이런 나무를 수백 개 쌓는다.}
\Verify{\ft{docs/notes/lecture/code/boost_toy.py}로 다시 계산할 수 있다. 로그 손실이 0.637 \ensuremath{\to} 0.448 \ensuremath{\to} 0.297로 매번 줄어든다.}
\Trap{나무 하나가 정답을 다 맞히게 하면(학습률 1, 깊은 나무) 학습 데이터의 우연까지 외운다. 그래서 작은 걸음(학습률 0.05)으로 천천히 쌓고, 검증 데이터에서 더 나아지지 않으면 멈춘다.}
\end{workout}

\paragraph{뉴턴 단계란.} 잎 값 = 잔차 합 \ensuremath{\div} 곡률 합. 잔차는 ``어느 방향으로 얼마나 틀렸나'', 곡률 $p(1-p)$는 ``그 점수를 바꾸면 확률이 얼마나 민감하게 변하나''다. $p$가 0.5 근처면 조금만 움직여도 되고, 0이나 1 근처면 크게 움직여야 한다.

\section{실제 설정과 학습 과정}
\begin{figure}[H]\centering
\begin{tikzpicture}[node distance=5mm,box/.style={draw=black!70,rounded corners=1.2mm,align=center,font=\small,minimum height=11mm,text width=27mm}]
\node[box] (a) {Y년 1월 1일까지\\결과 확정 대출};
\node[box,right=of a] (b) {특성 101개\\+ 정답(부도)};
\node[box,right=of b] (c) {LightGBM\\나무 최대 2000개};
\node[box,right=of c] (d) {시드 5개\\$p$ 평균};
\draw[-{Stealth},black!60] (a)--(b); \draw[-{Stealth},black!60] (b)--(c); \draw[-{Stealth},black!60] (c)--(d);
\end{tikzpicture}
\caption{연도 Y의 대출에 $p$를 붙이는 과정(시점 학습).}
\end{figure}

\begin{itemize}[leftmargin=*]
\item \textbf{학습 데이터:} 평가할 대출이 2016년 발행이면, 2016년 1월 1일까지 결과가 확정된 대출만 쓴다(부도는 마지막 납부 + 5개월에 확정, 3장). 2013년 모델은 11,345건, 2016년 모델은 156,153건으로 배운다.
\item \textbf{설정(고정, 결과를 보고 고르지 않음):} 학습률 0.05, 나무당 잎 63개, 행 80\%·\allowbreak{}열 60\% 표본추출, 잎당 최소 200건, L2 규제 5.
\item \textbf{조기 종료:} 학습 데이터의 15\%를 떼어 두고, 그 15\%의 로그 손실이 50번 연속 나아지지 않으면 나무 쌓기를 멈춘다.
\item \textbf{시드 평균:} 표본추출과 조기 종료용 분할이 시드에 따라 달라지므로 시드 5개로 모델 5개를 만들어 $p$를 평균한다.
\end{itemize}

\section{무엇이 $p$를 움직이나}
\begin{table}[H]\centering\small
\begin{tabular}{@{}rlll@{}}\toprule
순위 & 특성 & 방향 & 비중(SHAP)\\\midrule
1 & 금리 & 높을수록 위험 & 15.1\%\\
2 & 연소득 & 높을수록 안전 & 8.0\%\\
3 & 부채비율(DTI) & 높을수록 위험 & 6.1\%\\
4 & 거주 주 & 범주 & 4.1\%\\
5 & 가장 오래된 신용 계좌 시점 & 오래될수록 안전 & 3.1\%\\
6 & 만기 & 비선형 & 2.8\%\\
7 & 대출 목적 & 범주 & 2.5\%\\
8 & 대출액 & 클수록 위험 & 2.5\%\\\bottomrule
\end{tabular}
\caption{LightGBM이 $p$를 계산할 때 많이 쓴 특성(17단계).}
\end{table}
금리가 1위인 것은 자연스럽다. 금리는 LC가 이미 위험을 판단해 붙인 가격이기 때문이다. 중요한 것은 금리 다음에 소득, 부채비율, 신용 이력 같은 ``갚을 능력'' 정보가 따라온다는 점이다. 이것이 7장의 ``등급 안 정보''로 이어진다.

\section{왜 믿을 수 있나: 여섯 가지 근거}
\subsection*{근거 1. 순서가 맞는다(신뢰도 표)}
2013--16년 36개월 대출 52만 건에 시점 학습 $p$를 붙이고, $p$ 크기 순으로 10등분했다.

\begin{table}[H]\centering\small
\begin{tabular}{@{}rrrr@{}}\toprule
분위 & 평균 $p$ & 실제 부도율 & 평균 금리\\\midrule
1(가장 안전) & 2.6\% & 3.8\% & 7.3\%\\
2 & 4.2\% & 6.9\% & 9.2\%\\
3 & 5.5\% & 9.1\% & 10.2\%\\
4 & 6.9\% & 11.0\% & 11.0\%\\
5 & 8.3\% & 13.2\% & 11.6\%\\
6 & 10.1\% & 14.6\% & 12.2\%\\
7 & 12.5\% & 16.5\% & 12.8\%\\
8 & 15.7\% & 19.0\% & 13.6\%\\
9 & 20.8\% & 23.0\% & 14.8\%\\
10(가장 위험) & 33.2\% & 30.6\% & 17.2\%\\\bottomrule
\end{tabular}
\caption{$p$ 10분위별 실제 부도율(36단계). 분위마다 약 5.2만 건.}
\end{table}
실제 부도율이 한 칸도 거꾸로 가지 않고 3.8\%에서 30.6\%까지 오른다. 가장 위험한 10\%는 가장 안전한 10\%보다 8배 자주 부도났다.

\subsection*{근거 2. 같은 등급 안에서도 가른다}
등급이 이미 위험을 반영하므로, $p$가 등급만 따라 한다면 쓸모가 적다. 등급마다 $p$를 셋으로 나눴다.
\begin{table}[H]\centering\small
\begin{tabular}{@{}lrrrr@{}}\toprule
등급 & $p$ 낮음 부도율 & $p$ 중간 & $p$ 높음 & 등급 안 AUC\\\midrule
A & 3.1\% & 5.1\% & 7.9\% & 0.62\\
B & 8.4\% & 11.6\% & 15.2\% & 0.58\\
C & 14.5\% & 18.3\% & 23.4\% & 0.58\\
D & 19.8\% & 24.9\% & 30.6\% & 0.57\\
E & 23.7\% & 31.8\% & 39.6\% & 0.59\\\bottomrule
\end{tabular}
\caption{같은 등급 안에서 $p$로 나눈 실제 부도율(36단계).}
\end{table}
같은 C등급인데 $p$가 낮은 쪽은 14.5\%, 높은 쪽은 23.4\%가 부도났다. 등급이 놓친 차이를 $p$가 잡는다.

\subsection*{근거 3. 미래를 보지 않았다}
연도별 AUC는 2013년 0.631, 2014년 0.663, 2015년 0.671, 2016년 0.684다(0.5는 동전 던지기). 같은 대출로 무작위 분할을 하면 약 0.70(2007--20년 전체 종결 대출로 하면 0.73)으로 이보다 높지만, 그때 알 수 있던 정보만 쓴 \textbf{실제로 가능한 점수}다(3장). 학습 데이터가 늘수록 높아지는 것도 자연스럽다.

\subsection*{근거 4. 우연이 아니다}
\begin{itemize}[leftmargin=*]
\item \textbf{정답 섞기(위약):} 학습 라벨을 무작위로 섞어 만든 가짜 모델은 선별 가치가 0 근처였다. 평가 코드가 가짜 성과를 만들지 않는다.
\item \textbf{시드:} 시드를 바꿔도 결론이 같았다. 시드 차이(최대 0.09\%p)는 효과(약 1\%p)보다 훨씬 작다.
\item \textbf{발행월 부트스트랩:} 선별 가치 1.03\%p의 95\% 구간이 0.95--1.11로 0에서 멀다.
\end{itemize}

\subsection*{근거 5. 특정 모델의 산물이 아니다}
XGBoost는 같은 선별 가치(1.04\%p)를 냈고, 랜덤포레스트·\allowbreak{}로지스틱도 방향이 같았다(23--25단계). 결측 처리 9종, 피처 6종을 바꿔도 거의 같았다. 금리·\allowbreak{}등급을 빼고 학습한 모델도 AUC가 0.01 정도만 낮았다(34단계 F).

\subsection*{근거 6. 돈으로 이어진다}
$p$로 만든 기대수익 점수로 고르면, 실제로 들어온 돈으로 잰 수익률이 무작위 선택보다 높았다(6, 7장). 확률을 잘 맞히는 것과 돈을 버는 것은 다른 문제인데, 둘 다 확인했다.

\section{믿으면 안 되는 부분}
\begin{warnbox}[$p$의 크기는 틀린다]
\begin{itemize}[leftmargin=*]
\item 2013년 모델은 평균 $p$ 6.4\%를 냈지만 실제 부도율은 12.2\%였다.
\item \textbf{주원인은 학습 표본의 구성이다}(외부 검증의 분해). 학습 데이터에 들어가는 직전 해 발행분은 1년 안에 끝난 대출뿐인데, 그중 대부분이 조기상환이라 부도율이 최종의 절반 수준이다. 모델은 ``요즘 대출은 안전하다''고 배운다. 신용 환경 악화(추세)는 부차적이다.
\item 3장의 종결 편향(발행 후 1--3년 된 대출만 보면 부도가 과대 대표)과 방향이 반대로 보이지만 둘 다 사실이다. 나이가 1년 미만인 끝난 대출은 조기상환이, 1--3년인 끝난 대출은 부도가 과대 대표된다. $p$를 낮추는 것은 앞의 것이다.
\item 신뢰도 표의 아래쪽 분위에서 평균 $p$가 실제보다 1--4\%p 낮다.
\item 확률 보정(보정 3종)을 시도했지만 고치지 못했다(9단계). 과거로 만든 보정은 추세 변화를 따라가지 못한다.
\end{itemize}
\end{warnbox}

\begin{whybox}[그런데 왜 결과가 괜찮았나]
\textbf{직접 근거는 두 가지다.} (1) 부도확률 순(V2)은 $p$의 순서만 쓰므로, 크기가 모두 같은 비율로 틀리면 결과가 전혀 바뀌지 않는다. (2) 기대수익 순(V1)은 $s=(1-p)r-p\cdot\mathrm{LGD}-tb$에 금리 $r$이 함께 들어가 $p$의 크기가 순서를 일부 바꾼다. 그래서 $p$를 실제 부도율 수준으로 맞춰 다시 골라 봤고, 승인 대출의 88--92\%가 같았으며 전략 순위도 그대로였다.

\textbf{부수적 관찰(구조적 설명은 아님).} 점수에서 $p$는 손실률과 곱해진다. 손실률 추정(그때까지 부도 난 대출의 평균)은 실제보다 크게 나와 $p$의 과소추정과 같은 방향으로 맞물렸다. 외부 검증에 따르면 2013년에는 기대손실이 거의 맞았고, 이후 연도에는 12--75\% 크게 봤다. 이것은 결과를 사후에 본 관찰이지, 언제나 상쇄된다는 보장이 아니다.

\textbf{언제 문제가 되나.} 사람마다 틀린 정도가 달라 순서까지 망가질 때다. CatBoost가 2016년에 $p$를 평균 3.6\%로 낸 경우(실제 15.9\%)에는 성과가 무너졌다(24단계).
\end{whybox}

\begin{career}
``모델 확률을 왜 믿나요?''라는 질문에는 AUC 한 숫자로 답하지 말고, 신뢰도 표(순서), 등급 안 분리(추가 정보), 위약 검정(우연 아님), 다른 모델 비교(모델 독립), 돈으로의 연결(실제 수익)을 차례로 보여 주면 된다. 그리고 ``크기는 틀리지만 순위로만 쓴다''고 한계를 먼저 말하면 신뢰가 오히려 올라간다.
\end{career}

\chapter{점수와 채점: 무엇으로 고르고 무엇으로 평가하나}

5장의 $p$는 ``누가 더 위험한가''를 알려 준다. 이 장은 그 $p$로 \textbf{어떻게 고르는지}, 고른 결과를 \textbf{어떻게 채점하는지}를 다룬다. 고르는 도구와 채점하는 도구를 섞으면 결과가 무의미해지므로, 둘을 분명히 나눈다.

\section{기대수익 점수: $p$를 돈으로 바꾸기}
\begin{formulabox}[기대수익 점수]
\[ s=(1-p)\times r - p\times \mathrm{LGD} - tb \]
$r$: 약정 금리, LGD: 부도 시 손실률(그 시점까지 확정된 부도 대출에서 ``1 $-$ 받은 총액/대출액''의 평균, 2013년 모델 0.62, 2016년 0.52), $tb$: 같은 달 같은 만기 국채 수익률.
\end{formulabox}
\begin{itemize}[leftmargin=*]
\item $(1-p)\times r$: 안 망하면 받는 이자의 기댓값.
\item $p\times$LGD: 망하면 잃는 원금의 기댓값.
\item $tb$: 이 대출 대신 국채를 샀다면 위험 없이 받았을 수익(기회비용).
\end{itemize}

\begin{workout}[금리와 위험을 함께 봐야 하는 이유]
\Use{기대수익 점수(이 절)}
\Step{1} 예시: LGD 0.6, 국채 1.5\%. 대출 가($p$ 5\%, 금리 8\%), 대출 나($p$ 25\%, 금리 22\%).
\Step{2} 가: $0.95\times8\%-0.05\times0.6-1.5\%=7.6-3.0-1.5=+3.1\%$p.
\Step{3} 나: $0.75\times22\%-0.25\times0.6-1.5\%=16.5-15.0-1.5=0$.
\Step{4} 나의 실제 위험이 15\%라면: $0.85\times22\%-0.15\times0.6-1.5\%=18.7-9.0-1.5=+8.2\%$p.
\Ans{부도확률만 보면 항상 가를 사지만, 위험에 비해 금리가 후한 대출은 나 같은 고금리 대출일 수 있다.}
\Trap{점수를 ``정확한 예상 수익''으로 읽는다. $p$는 생애 기준이고 금리는 연 단위라, 이 식은 순위를 매기는 1년 단위 근사다.}
\end{workout}

점수는 두 곳에 쓴다. (1) 승인 규칙: $s>0$, 즉 국채보다 더 벌 것 같으면 승인. (2) 순위: 점수 상위 10\%, 20\%, $\dots$를 승인해 본다. 같은 달 같은 만기 대출은 국채가 같으므로 순위를 정하는 것은 사실상 ``금리와 위험의 균형''이다.

\section{고르는 도구와 채점하는 도구}
\begin{table}[H]\centering\small
\begin{tabularx}{\linewidth}{@{}p{24mm}LL@{}}\toprule
 & 기대수익 점수 & IRR(현금흐름 수익률)\\\midrule
언제 & 승인하는 날 & 대출이 끝난 뒤\\
무엇으로 & 모델의 $p$, 약정 금리, 평균 손실률, 국채 & 실제로 받은 돈과 받은 시점, 수수료 1\%\\
정체 & 예상 & 실제 결과\\
역할 & 줄 세워 고른다 & 고른 묶음이 실제로 얼마 벌었는지 채점\\\bottomrule
\end{tabularx}
\caption{두 도구의 역할(2장의 IRR 참고).}
\end{table}
IRR로 고를 수는 없다(승인하는 날에는 미래를 모른다). 점수로 채점할 수도 없다(자기 예상으로 자기를 채점하는 셈이다). \textit{비유:} 학원 예상 점수로 지원 대학을 고르고, 실제 수능 성적으로 그 예상이 맞았는지 채점한다.

\section{왜 ``아무거나 고를 때''와 비교하나}
\textbf{질문.} 우리가 고른 대출의 수익률이 7\%였다. 잘한 걸까?
\begin{itemize}[leftmargin=*]
\item 그해 아무거나 샀어도 6\%였다면 우리 실력은 1\%p다.
\item 그해 아무거나 샀어도 8\%였다면 우리는 고르다가 1\%p를 잃었다.
\end{itemize}
수익률 하나로는 구분이 안 된다. 그래서 같은 달, 같은 승인비율로 \textbf{무작위로} 고른 결과를 빼서 ``고르는 실력''만 남긴다. 펀드매니저를 시장 지수 대비로 평가하는 것과 같은 원리다.

\begin{formulabox}[선별 가치]
\[ \mathrm{SV}=\frac19\sum_{q=10\%}^{90\%}\Big[\overline{(\mathrm{IRR}-tb)}_{\text{점수 상위 }q}-\overline{(\mathrm{IRR}-tb)}_{\text{무작위 }q}\Big] \]
승인 대상은 발행월마다 그 달 대출 중 상위 $q$로 자르고, 발행월 묶음 IRR을 계산해 평균한다. 승인비율 하나를 골라 유리하게 보일 수 없도록 9개 비율을 평균한다. (25단계까지는 발행연도마다 잘랐다. 발행월로 바꿔도 1.03 \ensuremath{\to} 1.07\%p로 거의 같다.)
\end{formulabox}

이 기준이 좋은 이유는 세 가지다. (1) 시장 효과가 양쪽에 똑같이 들어가 지워진다. (2) 승인비율을 맞추므로 ``얼마나 샀나''가 아니라 ``무엇을 샀나''만 비교한다. (3) 0이 뜻을 갖는다. 무작위는 정보가 없는 투자자이고, 실제로 정답을 섞은 가짜 모델은 0 근처였다.

\begin{table}[H]\centering\small
\begin{tabularx}{\linewidth}{@{}p{34mm}LL@{}}\toprule
비교 대상 & 답하는 질문 & 주 기준이 아닌/인 이유\\\midrule
국채 & LC 대출에 투자할 가치가 있나? & 2013--16년은 전부 사도 국채를 이겼다. 모델 없이도 양수라 고르는 실력을 못 보여 준다\\
전부 승인 & 모델을 쓰는 게 나은가? & 승인비율이 달라 ``무엇''과 ``얼마나''가 섞인다\\
\textbf{같은 비율 무작위} & \textbf{고르는 방식이 가치를 더하나?} & \textbf{주 기준. 시장과 승인비율을 모두 맞춘다}\\\bottomrule
\end{tabularx}
\end{table}

\section{Sharpe 이야기: 세 가지 정의}
교수님은 ``결과는 Sharpe 극대화''를 말씀하셨고, 조교는 ``0.2--0.4면 옳게 한 것''이라고 했다. 그런데 Sharpe는 무엇을 한 단위로 보느냐에 따라 전혀 다른 숫자가 된다.

\begin{table}[H]\centering\small
\begin{tabularx}{\linewidth}{@{}p{40mm}LLp{20mm}@{}}\toprule
정의 & 계산 & 특징 & 우리 값\\\midrule
건별 Sharpe(거절 = 0) & 모든 대출에 대해, 승인이면 그 대출의 초과수익, 거절이면 0. 평균 \ensuremath{\div} 표준편차 & 승인비율만 줄여도 변한다(아래) & 0.03--0.15\\
\textbf{건별 Sharpe(승인만)} & 승인한 대출들의 초과수익 평균 \ensuremath{\div} 표준편차 & \textbf{조교의 0.2--0.4와 크기가 맞는다} & 0.10--0.34\\
발행월 포트폴리오 Sharpe & 매달 (승인 대출 + 거절 자금 국채) 포트폴리오의 초과수익을 48개월에 걸쳐 평균 \ensuremath{\div} 표준편차 & 승인비율이 지워진다. ``매달 얼마나 고르게 벌었나''를 잰다 & 약 2--3.5\\\bottomrule
\end{tabularx}
\caption{Sharpe 정의 세 가지(32단계). 건별 정의의 대출 수익률은 받은 돈을 국채로 재투자한 만기 수익률이다(대출 한 건의 IRR은 부도 시 $-77\%$처럼 극단값이 커서 쓰지 않는다).}
\end{table}

\begin{workout}[거절을 0으로 세면 생기는 일]
\Use{건별 Sharpe(거절 = 0)}
\Step{1} 예시: 승인한 대출의 초과수익 평균 4\%, 표준편차 20\%. 승인비율 $a$만 바꾼다.
\Step{2} 승인 100\%: $0.04/0.20=0.200$.
\Step{3} 승인 50\%: $\sqrt{0.5}\times0.04/\sqrt{0.04+0.5\times0.0016}=0.140$.
\Step{4} 승인 10\%: 0.062.
\Ans{대출의 질은 같은데 숫자만 1/3이 된다. 이 정의는 ``얼마나 승인했나''에 끌려간다.}
\end{workout}

포트폴리오 Sharpe는 이 문제가 없다(팀원 제안서의 지적, 30단계). 매달 초과수익이 (승인 금액 비중 \ensuremath{\times} 승인 묶음 초과수익)이므로 비중이 평균과 표준편차에서 함께 지워진다. 다만 규모를 무시하고, 위험을 2013--16년(경기침체 없음) 48개월의 출렁임으로만 재며, 그 크기(약 3)는 주식 Sharpe(0.5 안팎)와 비교할 수 없는 \textbf{순위 비교용} 숫자다.

\begin{warnbox}[자르는 단위가 중요하다]
포트폴리오 Sharpe는 매달 승인 금액 비중이 일정하다는 전제에서 승인비율이 지워진다. 승인 대상을 발행연도 단위로 자르면 고금리 대출이 많은 달에는 많이, 적은 달에는 적게 승인되어 비중이 들쭉날쭉해지고, 그 출렁임이 Sharpe를 깎는다. 실제로 기대수익 순의 Sharpe는 연도 단위 2.46, 발행월 단위 2.91이었다(외부 검증, 37단계). 연도 단위는 12월 대출까지 본 뒤 1월 승인을 정하는 셈이라 실전에서도 쓸 수 없다.
\end{warnbox}

\section{그래서 두 축으로 본다}
\begin{keybox}[수익 축과 안정성 축]
\begin{itemize}[leftmargin=*]
\item \textbf{수익:} 선별 가치(IRR 기준), 총 NPV(국채로 할인한 순현재가치, 달러 단위).
\item \textbf{안정성:} 발행월 포트폴리오 Sharpe, 승인 대출 건별 Sharpe, 하위 5\% 평균(꼬리 손실).
\end{itemize}
Sortino(아래쪽 출렁임만 위험으로 보는 Sharpe)도 계산했지만, 9개 승인비율 모두에서 전략 순위가 건별 Sharpe와 같아 새 정보가 없었다(33단계). 포트폴리오 단위 Sortino는 48개월 모두 국채보다 높아 분모가 0이 되어 계산할 수 없다.
\Intuit{더 버는 규칙과 더 고르게 버는 규칙은 다를 수 있다. 하나의 숫자로 우열을 정하지 말고 두 축으로 본다.}
\end{keybox}

\section{성과의 출처를 나누는 두 도구}
\paragraph{Brinson 분해.} 펀드 평가 도구다. 초과수익 = 배분(좋은 등급을 더 많이 샀나) + 선별(같은 등급 안에서 더 나은 대출을 골랐나) + 상호작용. 예를 들어 기준이 A 60\%(4\%)·\allowbreak{}B 40\%(7\%)로 5.2\%이고 모델이 A 70\%(5\%)·\allowbreak{}B 30\%(8\%)로 5.9\%라면, 0.7\%p = 배분 $-0.3$ + 선별 $+1.0$ + 상호작용 0이다.

\paragraph{등급 중립 롱숏.} 주식 팩터 투자에서 신호의 크기를 재는 방법이다. 매달, 등급마다 신호 상위 30\%를 산다고(롱) 하위 30\%를 판다고(숏) 가정하고 두 묶음의 IRR 차이를 기록한다. 등급마다 롱과 숏을 같은 비율로 잡으므로 등급 배분과 시장 효과가 정확히 지워지고 \textbf{등급 안 선별}만 남는다. LC 대출은 공매도가 불가능하므로 실제 전략이 아니라 \textbf{측정 도구}다. 실제로 가능한 형태는 ``상위 30\% $-$ 그 등급 전체''(롱온리)다.

\chapter{결과: 수익, 안정성, 그리고 등급 안의 정보}

모든 결과는 2013--16년 발행 36개월 종결 대출(약 52만 건, 48개 발행월), 시점 학습, 투자자 수수료 1\% 차감 기준이다. 판정 기준은 각 실험 전에 정했다.

\section{출발점: LC 대출 자체는 국채보다 나았다}
2013--16년 36개월 대출을 \textbf{전부} 사면 발행월 묶음 IRR 5.73\%로 같은 기간 3년 국채(0.87\%)보다 높았고, 48개 발행월 \textbf{모두}에서 그랬다. 지난 우승팀은 단순 연환산과 최근 연도의 끝난 대출만 보는 설정에서 반대 결론(전부 사면 국채보다 못하다)을 냈다(2장). 그래서 우리의 질문은 ``LC에 투자할까''가 아니라 ``\textbf{무엇을 고르면 더 나은가}''가 된다.

\section{수익 축: 기대수익 순위가 가장 많이 번다}
\begin{table}[H]\centering\small
\begin{tabular}{@{}llr@{}}\toprule
비교 & 결과 & 근거\\\midrule
기대수익 순위 vs 무작위 & \textbf{+1.03\%p/년} (95\% 구간 0.95--1.11) & 25단계\\
기대수익 순위 vs 부도확률 순위 & +0.91\%p (0.73--1.10, 다 끝난 2013--15년) & 26단계\\
대출마다 같은 금액을 넣는 투자자 기준 & +1.12\%p (1.04--1.20) & 29단계\\
정답을 섞은 가짜 모델 & 0 근처 & 14단계\\\bottomrule
\end{tabular}
\caption{수익 축의 핵심 결과.}
\end{table}

\begin{table}[H]\centering\small
\begin{tabular}{@{}lrrl@{}}\toprule
모델 & AUC & 선별 가치 & 판정\\\midrule
LightGBM(기준) & 0.659 & 1.03\%p & 유지\\
XGBoost & 0.655 & 1.04\%p & 같음($p=0.41$)\\
랜덤포레스트 & \textbf{0.669} & 0.85\%p & AUC는 1위, 돈은 덜 범\\
Lasso + 피처 & 0.649 & 0.76\%p & 선형 중 최선\\
CatBoost & 0.651 & 0.68\%p & 2016년 $p$ 붕괴(평균 3.6\%)\\\bottomrule
\end{tabular}
\caption{모델 17종 중 주요 결과(23--25단계).}
\end{table}
두 가지를 읽을 수 있다. 첫째, \textbf{잘 맞히는 모델이 돈을 더 버는 것은 아니다}(AUC 1위 랜덤포레스트). 둘째, \textbf{1\%p는 모델의 산물이 아니다}. XGBoost가 같은 값을 냈고, 튜닝·\allowbreak{}새 피처·\allowbreak{}결측 처리로 바뀌지 않았다.

\paragraph{돈 단위로(NPV, 국채로 할인).} 승인 30\%에서 총 NPV는 기대수익 순 1억 2,800만 달러, 부도확률 순 1억 1,900만, 무작위 1억 100만 달러다. 9개 승인비율 모두 기대수익 순이 1위다. 다만 승인 50\% 이상에서는 부도확률 순과의 차이가 2--3\%뿐이다(33단계).

\section{등급 안의 정보: 이 프로젝트의 핵심 주장}
\subsection*{Brinson 분해}
기대수익 상위 50\%가 더 번 수익을 나누면 대부분이 \textbf{선별}(같은 등급 안에서 더 나은 대출)에서 나왔다. 세부등급 35개로 나눠도, 금리 구간·\allowbreak{}목적·\allowbreak{}지역·\allowbreak{}금액으로 묶어도, 60개월 대출에서도 같았다(14, 18단계).

\subsection*{등급 중립 롱숏}
\begin{table}[H]\centering\small
\begin{tabular}{@{}lrrr@{}}\toprule
신호 & 롱 $-$ 숏 & 90\% 구간 & 양수 달\\\midrule
기대수익 점수 & \textbf{+2.38\%p} & 2.08--2.73 & 48/48\\
최근 24개월 개설 계좌(적을수록) & +2.32\%p & 1.99--2.59 & 48/48\\
부도확률(낮을수록) & +2.01\%p & 1.72--2.34 & 48/48\\
연소득(높을수록) & +1.09\%p & 0.68--1.74 & 39/48\\
부채비율(낮을수록) & +0.95\%p & 0.76--1.24 & 44/48\\
FICO(높을수록) & +0.50\%p & 0.36--0.71 & 41/48\\
무작위(위약) & $-0.12$\%p & $-0.21$--$+0.03$ & 22/48\\\bottomrule
\end{tabular}
\caption{발행월·\allowbreak{}등급마다 신호 상위 30\% $-$ 하위 30\%의 IRR 차이(33단계). 실제로 가능한 롱온리(상위 30\% $-$ 그 등급 전체)는 기대수익 점수 +1.20\%p.}
\end{table}

\begin{keybox}[이 표가 말하는 것]
\begin{itemize}[leftmargin=*]
\item 등급과 시장 효과를 지워도, 같은 등급 안에서 고른 쪽이 \textbf{48개월 전부} 더 벌었다. 위약은 0이다.
\item \textbf{FICO는 약하다.} LC가 신용점수는 금리에 잘 반영했다.
\item \textbf{최근 계좌 개설 수 하나로 모델만큼의 신호}가 나온다. 짧은 기간에 대출·\allowbreak{}카드를 여러 개 연 사람은 돈이 급하거나 신용을 끌어다 쓰는 중일 가능성이 크고, LC는 이 행동을 금리에 덜 반영했다.
\end{itemize}
\Intuit{LC는 ``얼마나 신용이 좋은가''는 잘 쟀고, ``요즘 돈을 얼마나 찾아다니나''는 덜 쟀다.}
\end{keybox}

\subsection*{이 신호는 견고한가(34단계)}
\begin{table}[H]\centering\small
\begin{tabularx}{\linewidth}{@{}p{50mm}L@{}}\toprule
점검 & 결과\\\midrule
해마다 줄었나 & 기대수익 점수 1.94 \ensuremath{\to} 2.70 \ensuremath{\to} 1.81 \ensuremath{\to} 3.05\%p, 계좌 개설 1.79 \ensuremath{\to} 3.05\%p 근처로 줄지 않았다. 연소득만 2.08 \ensuremath{\to} 0.12\%p로 사라졌다\\
개인이 살 수 있던 대출(조각 대출)만 & +2.58\%p, 48/48개월. 조각 대출 비중은 72\% \ensuremath{\to} 27\%로 줄었다\\
60개월 대출(2013--14) & 계좌 개설 신호 +3.53\%p, 전 개월 양수\\
LC 가격(금리·\allowbreak{}등급) 없이 학습한 모델 & AUC 0.01 정도만 낮고, 등급 중립 롱숏 +2.17\%p\\\bottomrule
\end{tabularx}
\end{table}

\section{안정성 축: Sharpe로 보면 보수적 선별이 이긴다}
\begin{table}[H]\centering\small
\begin{tabular}{@{}lrrrr@{}}\toprule
승인 30\% & 포트폴리오 Sharpe & 평균 초과수익 & 건별 Sharpe(승인만) & 하위 5\% 평균\\\midrule
부도확률 낮은 순 & \textbf{3.24} & 1.66\%p & \textbf{0.26} & \textbf{$-20$\%}\\
금리 낮은 순 & 3.02 & 1.49\%p & 0.22 & $-21$\%\\
기대수익 순 & 2.46 & \textbf{1.92\%p} & 0.17 & $-32$\%\\
무작위 & 2.68 & 1.46\%p & 0.10 & $-34$\%\\\bottomrule
\end{tabular}
\caption{승인 30\%에서 두 축(30, 32, 33단계). 9개 승인비율 모두 같은 방향이고, 적게 승인할수록 차이가 크다.}
\end{table}
기대수익 순은 가장 많이 벌지만 꼬리 손실은 무작위와 거의 같다. 부도확률 순은 조금 덜 벌고 매달 훨씬 고르게 번다.

\subsection*{왜 그런가: 고위험 등급의 가격 실패}
\begin{table}[H]\centering\small
\begin{tabular}{@{}lrrr@{}}\toprule
등급(36개월) & 평균 초과수익 & Sharpe & 국채 이하 달\\\midrule
A & 4.0\%p & 5.2 & 0\\
B & 5.3\%p & 3.4 & 0\\
C & 5.3\%p & 2.3 & 0\\
D & 4.7\%p & 1.6 & 2\\
E & 3.3\%p & 0.6 & 17\\
F & 1.1\%p & 0.1 & 19\\
G & $-9.7$\%p & $-0.6$ & 6(8개월 중)\\\bottomrule
\end{tabular}
\caption{등급별 발행월 초과수익(33단계).}
\end{table}
Sharpe가 A에서 G로 내려가고, E--G는 평균 수익도 낮다. 36개월 고위험 대출은 \textbf{금리가 위험만큼 높지 않았다}. 기대수익 순은 금리가 높은 대출을 좋아하므로 이 영역의 출렁임을 떠안는다.

\begin{warnbox}[60개월은 다르다]
60개월(2013--14)에서는 평균 초과수익이 E(8.7\%p)까지 오르고 Sharpe는 B·\allowbreak{}C(6.7)가 가장 높았다. ``고위험 등급은 손해''는 36개월에만 맞는다. test 결과는 만기별로 따로 봐야 한다.
\end{warnbox}

\section{두 축의 화해}
등급 구성을 무작위와 같게 두고 등급 안에서만 부도확률 낮은 순으로 고르면, 무작위보다 수익은 크게 오르고(+0.34\%p) Sharpe는 조금만 오른다(2.68 \ensuremath{\to} 2.81). 여기에 안전한 등급 쪽으로 쏠리는 효과를 더하면(전체에서 부도확률 낮은 순) Sharpe는 크게 오르고(3.24) 수익은 내려간다.

\begin{keybox}[화해]
\textbf{수익은 같은 등급 안의 선별에서, 안정성은 안전한 등급으로의 배분에서 주로 나온다.} 우리 원래 주장(등급 안 선별)과 팀원 지적(보수적 선별의 Sharpe 우위)은 둘 다 맞았고, 서로 다른 축을 말하고 있었다.
\end{keybox}

\section{투명한 단순 규칙: 둘 다 잡을 수 있나}
``등급마다 최근 24개월 개설 계좌가 적은 대출을 승인''을 시험했다. 처음에는 수익·\allowbreak{}Sharpe 모두 1위였는데 \textbf{동점 버그}였다. 계좌 수가 정수라 경계값 동점이 모두 승인되어 실제 승인비율이 30\%가 아니라 39.5\%였다. 동점을 무작위로 끊어 다시 계산했다.

\begin{table}[H]\centering\small
\begin{tabular}{@{}lrr@{}}\toprule
승인 30\%(50\%) & 포트폴리오 수익 & 포트폴리오 Sharpe\\\midrule
단순 규칙 & 1.86\%p (2.88) & \textbf{3.72 (3.41)}\\
기대수익 순 & \textbf{1.92\%p (3.02)} & 2.46 (2.81)\\
부도확률 순 & 1.66\%p (2.83) & 3.24 (3.05)\\
무작위 & 1.45\%p (2.43) & 2.72 (2.69)\\\bottomrule
\end{tabular}
\caption{단순 규칙(34b단계).}
\end{table}
단순 규칙은 등급 구성을 바꾸지 않는데도 부도확률 순을 수익과 Sharpe \textbf{둘 다}에서 이긴다. ``안정성은 등급 배분에서만 나온다''가 아니었다. 무엇으로 고르느냐에 따라 등급 안 선별만으로 두 축을 다 얻을 수 있다. 다만 33단계에서 같은 데이터로 가장 강한 신호를 보고 고른 규칙이므로 \textbf{과대평가되었을 수 있다}. 그래서 test 규칙 v2에 가설 H7로 미리 고정했다(9장).

\section{투자자 목적별 메뉴}
\begin{table}[H]\centering\small
\begin{tabularx}{\linewidth}{@{}p{45mm}L@{}}\toprule
목적 & 규칙\\\midrule
가장 많이 벌기 & 기대수익 순\\
매달 고르게 벌기(과제의 Sharpe) & 부도확률 낮은 순, 또는 A·\allowbreak{}B등급\\
둘 다(후보, test에서 검증) & 등급 안 계좌 개설 적은 순\\
균형 & 등급 안 부도확률 낮은 순\\\bottomrule
\end{tabularx}
\end{table}

\chapter{검증과 정정의 기록}

\section{우연을 지우는 도구들}
\begin{table}[H]\centering\small
\begin{tabularx}{\linewidth}{@{}p{38mm}LL@{}}\toprule
의심 & 도구 & 결과\\\midrule
운이 좋았던 것 아닌가 & 시드 3--10개, 발행월 단위 부트스트랩(1만 번, 6개월 이동 블록) & 구간이 0에서 멀다\\
평가 코드가 가짜 성과를 만드나 & 정답을 무작위로 섞은 가짜 모델, 무작위 신호 롱숏 & 0 근처(정상)\\
미래 정보가 새지 않았나 & 컬럼 시점 점검, 시점 학습, 상각 지연 5개월, 약한 누수 후보(투자자 펀딩액) 제거 & 결과 동일\\
모델·\allowbreak{}전처리 덕인가 & 모델 17종, 결측 처리 9종, 피처 6종 & 거의 같음\\
채점 표본이 편향됐나 & 다 끝난 대출을 일부러 일찍 자르는 시뮬레이션 & 편향 확인, test 해석 규칙 고정\\
결과를 보고 기준을 바꾸지 않았나 & 판정 기준을 스크립트 맨 위에 먼저 적음, test 규칙 git 태그 & 실패도 실패로 기록\\
여러 번 비교한 탓 아닌가 & Holm 보정, 시도한 조합 수 공개(80개 이상) & 핵심 결과는 보정 후에도 유의\\\bottomrule
\end{tabularx}
\end{table}

\section{처음 판단이 틀렸던 것}
\begin{table}[H]\centering\small
\begin{tabularx}{\linewidth}{@{}LLp{22mm}@{}}\toprule
처음 판단 & 실제 & 어떻게 알았나\\\midrule
부도확률 순위는 무작위와 비슷하다 & 시드 여러 개로 보면 유의하게 양수 & 14단계\\
선별 곡선은 새로운 방법이다 & 수익 기반 점수(2016), EMP(2014)가 이미 있음 & 문헌 검색\\
로지스틱이 무너진 건 결측 처리 탓 & 학습 때 거의 없던 변수의 표준화 폭발 & 22단계\\
CatBoost 붕괴는 새 컬럼이 조기 종결의 대리 변수라서 & 부도율 17.2\% vs 17.6\%로 같아 기각, 해법(가용성 규칙)만 확인 & 24b단계\\
$p$의 크기가 맞아야 돈을 번다 & 극단적으로 틀릴 때만 문제 & 25단계\\
2017년에는 부도확률 순위가 더 낫다 & 끝난 대출만 본 종결 편향 & 26단계\\
\textbf{Sharpe는 승인비율에 끌려가 주 지표로 부적합} & 대출 건별 Sharpe에만 맞음. 포트폴리오 Sharpe에서는 보수적 선별이 우위(팀원 지적) & 30단계\\
\textbf{LC 등급끼리의 가격은 공정했다} & A--D 평균 수익에만 맞음. 36개월 E--G는 금리가 위험을 못 메움 & 33단계\\
고위험 등급은 손해다 & 36개월에만 맞음. 60개월은 반대 & 34단계\\
단순 규칙이 수익·\allowbreak{}Sharpe 모두 1위 & 동점 버그(실제 승인 39.5\%). 고치면 Sharpe 우위만 & 34b단계\\
\textbf{기대수익 순은 무작위보다 Sharpe가 낮다} & 발행연도 단위 컷이 만든 값. 발행월 단위면 2.91로 무작위(2.69)보다 높다. 철회 & 외부 검증, 37단계\\
\textbf{안정성은 부도확률 낮은 순이 1위} & 허들이 국채 수준일 때만. BAA 기준이면 기대수익 순 & 외부 검증, 37--38단계\\
\textbf{1-SE 기준선 안 = 최고와 통계적으로 구분 안 됨} & 1위 하나의 SE로 그은 선. 짝지은 비교로는 10\% $-$ 80\% = 0.76(0.45--1.60)으로 구분됨. 경험적 선택 규칙으로 정정 & 팀원 검토, 39단계 재점검\\
\textbf{승인비율이 달라도 평균과 표준편차는 비례할 뿐} & 대출 구성도 바뀐다(Sharpe 3.82 \ensuremath{\to} 3.06) & 팀원 검토, 41단계 재점검\\
기대수익 순은 작은 예산에서도 무작위보다 덜 흔들린다 & 10\%에서 Sharpe 2.55 vs 2.63, 50\% 이하 구분 안 됨. 60\% 이상에서만 높음 & 42단계\\
과제가 Sharpe를 요구하니 Sharpe 기준으로 최종안을 고른다 & 공지 재확인 결과 필수 아님. 최종안을 기대수익 순(v1 H2)으로 변경 & 공지 재확인, v2.3\\\bottomrule
\end{tabularx}
\caption{정정 기록(\ft{docs/research_log.md} 3절의 요약).}
\end{table}

\begin{insight}[정정에서 배운 것]
정정의 절반 이상이 ``숫자는 맞는데 범위를 너무 넓게 말한'' 경우였다. Sharpe 주장은 건별 정의에만, 등급 가격 주장은 A--D에만, 고위험 손해 주장은 36개월에만 맞았다. 그래서 이제 모든 주장에 \textbf{어느 표본, 어느 정의에서 맞는지}를 붙인다.
\end{insight}

\begin{warnbox}[스스로 검증할 때 왜 놓쳤나]
34단계에서 ``근본 가정 재점검''을 했지만 위의 마지막 두 정정은 외부 검증자가 찾았다. (1) 30--36단계가 같은 승인 함수(발행연도 단위)를 재사용했고, 다시 확인할 때도 그 함수로 돌려 함수 안의 선택이 의심받지 않았다. (2) 본문의 조건 설명(``같은 달'')을 코드가 아니라 의도대로 썼다. (3) 점검 목록에 비교 기준(허들)을 넣지 않았다. 외부 검증자는 본문만 보고 처음부터 따로 구현해서 차이를 찾았다. 그래서 이제 핵심 숫자는 독립 구현으로 한 번 더 확인하고, 문서의 조건은 코드와 대조한다.
\end{warnbox}

\section{데이터 엿보기의 위험}
같은 2013--16년 구간을 1단계부터 36단계까지 계속 다시 썼다. 판정 기준을 미리 정했지만, 수십 번 들여다본 구간에서 좋아 보이는 것은 그 구간에 맞춰졌을 수 있다. 특히 단순 규칙(7장)은 결과를 본 뒤 고른 것이다. 이 위험을 줄이는 유일한 방법은 한 번도 보지 않은 공식 test로 미리 정한 규칙을 한 번만 돌리는 것이다(10장).

\chapter{최종 제안: 기대수익이 높은 대출부터 산다}

\begin{keybox}[최종 제안(test 규칙 v2.3, 42단계)]
\textbf{매달 새로 나온 대출마다 기대수익 점수 $s$를 계산하고, 승인비율이 몇 \%든 $s$가 높은 대출부터 산다. 사지 않은 돈은 국채에 둔다.}
\begin{itemize}[leftmargin=*]
\item 쓰는 전략은 \textbf{기대수익 순 하나}다. 몇 \%를 살지는 정하지 않는다. 승인비율은 투자자가 가진 돈에 달려 있고, 우리는 그 안에서 \textbf{무엇을 살지}를 답한다.
\item \textbf{주장:} 승인비율이 10\%든 90\%든, 같은 양을 무작위로 산 것보다 \textbf{더 번다}.
\item \textbf{주 판정:} 처음 규칙서(v1)의 H2를 정의 그대로 쓴다. Sharpe(흔들림)는 보조 지표로 보고한다.
\end{itemize}
\end{keybox}

\begin{formulabox}[승인비율 $q$란]
매달 그 달에 나온 36개월 대출을 점수 순으로 줄 세워 위에서부터 사는 \textbf{건수의 비율}이다(10\%, 20\%, \dots, 90\%). 투자자가 그 달 시장의 몇 \%를 살 돈이 있는지를 나타내며, 사지 않은 돈은 국채에 둔다고 본다. 대출마다 금액이 달라 \textbf{금액 비중} $w$는 $q$와 조금 다르다(예: 부도확률 순 80\%에서 평균 82.5\%). 선별 가치는 승인한 1달러당 IRR 차이로, Sharpe는 금액 비중 $w$를 곱한 초과수익으로 계산한다.
\end{formulabox}

\begin{formulabox}[선별 가치와 포트폴리오 초과수익은 다르다]
\textbf{선별 가치}는 \textbf{대출에 넣은 돈의 질}을 잰다. 어느 달 대출이 1,000건이고 승인비율이 10\%면, 점수 상위 100건과 아무 100건을 각각 사서 대출이 끝난 뒤 두 묶음의 실제 IRR을 비교한다. 예를 들어 우리 묶음 연 8.0\%, 무작위 묶음 연 5.9\%면 선별 가치는 2.1\%p다. IRR은 ``넣은 1달러가 1년에 몇 \% 벌었나''이므로, 얼마를 넣었는지와 상관없이 \textbf{같은 양을 무작위로 산 것보다 1달러당 몇 \%p 더 버나}를 뜻한다. 48개월 평균이 표의 숫자다. 고르는 실력을 재는 숫자이고, 주 판정에 쓴다.

\textbf{포트폴리오 초과수익}은 \textbf{전체 돈의 성과}다. 사지 않은 돈은 국채에 두어 초과수익이 0이므로, 승인비율 10\%면 돈의 90\%가 국채보다 더 벌지 못한다. 투자자가 실제로 국채보다 얼마나 더 버는지는 이 숫자가 보여 준다.

\smallskip
\centering\small
\begin{tabular}{@{}lccccc@{}}\toprule
기대수익 순(V1), train & 10\% & 50\% & 80\% & 90\% & 전부 승인\\\midrule
선별 가치(1달러당, 무작위 대비) & +2.11 & +0.99 & +0.48 & +0.28 & 0\\
포트폴리오 초과수익(국채 대비, 전체 돈) & 0.65 & 2.98 & 4.34 & 4.67 & 4.86\\\bottomrule
\end{tabular}

\raggedright
\smallskip
\textbf{읽는 법:} 적게 살수록 가장 좋은 대출만 골라 1달러당 차이가 크지만, 돈 대부분이 국채에 있어 전체 수익은 작다. 많이 살수록 평범한 대출까지 사게 되어 1달러당 차이는 줄지만 전체 수익은 커진다. 그래서 ``10\%만 사는 것이 가장 좋다''는 뜻이 아니다. 이 기간에는 대출이 국채보다 나았으므로 많이 살수록 전체 수익은 컸고, 우리 주장은 \textbf{몇 \%를 사든 같은 양을 무작위로 사는 것보다 낫다}는 것이다.
\end{formulabox}

\section{흐름: 고르는 날과 채점하는 날}
\begin{figure}[H]\centering
\begin{tikzpicture}[x=1mm,y=1mm,box/.style={draw=black!70,rounded corners=1.2mm,align=center,font=\footnotesize,minimum height=12mm,text width=32mm,fill=white,execute at begin node={\XeTeXlinebreakskip=0pt\relax}},ar/.style={-{Stealth},black!65}]
\node[box] (a) at (0,0) {그 달 새 대출\\(36개월 만기)};
\node[box] (b) at (42,0) {부도확률 $p$\\시점 학습 LightGBM\\(그날 알던 정보만)};
\node[box] (c) at (84,0) {기대수익 점수\\$s=(1-p)r-p\cdot\mathrm{LGD}$\\$-$ 국채 금리};
\node[box,line width=0.9pt] (d) at (126,0) {$s$ 높은 순으로\\예산만큼 구매\\남는 돈은 국채};
\draw[ar] (a)--(b); \draw[ar] (b)--(c); \draw[ar] (c)--(d);
\draw[dashed,black!35] (-18,-13)--(144,-13);
\node[font=\scriptsize\itshape,black!60,anchor=south west] at (-18,-13) {위: 승인하는 날(결과를 모름)};
\node[font=\scriptsize\itshape,black!60,anchor=north west] at (-18,-13) {아래: 대출이 끝난 뒤(실제 결과로 채점, $s$는 쓰지 않음)};
\node[box] (e) at (126,-27) {우리가 산 묶음\\실제 현금흐름\\\ensuremath{\to} 코호트 IRR};
\node[box] (f) at (126,-47) {같은 양 무작위(10회)\\실제 현금흐름\\\ensuremath{\to} 코호트 IRR};
\node[box,line width=0.9pt] (g) at (84,-37) {선별 가치\\$=$ 우리 IRR $-$ 무작위 IRR\\(국채는 양쪽에서 지워짐)};
\node[box] (h) at (42,-37) {48개 발행월 평균\\예산 10--90\%\\블록 부트스트랩 구간};
\node[box,line width=0.9pt] (i) at (0,-27) {주 판정(H2)\\95\% 구간 하한 $>0$\\\ensuremath{\to} 제안 유지};
\node[box,dashed] (j) at (0,-47) {보조 지표\\포트폴리오 Sharpe\\(평균 \ensuremath{\div} 흔들림)};
\draw[ar] (d)--(e); \draw[ar] (e)--(g); \draw[ar] (f)--(g); \draw[ar] (g)--(h); \draw[ar] (h)--(i); \draw[ar,dashed] (h)--(j);
\end{tikzpicture}
\caption{제안의 흐름. 위 줄은 승인하는 날 하는 일, 아래 줄은 대출이 끝난 뒤 채점하는 일이다. 코호트 IRR은 빌려준 돈(0개월), 월 납입금, 마지막 달의 조기상환 일시금이나 부도 후 회수금을 묶음 전체로 더해 구하고, 받은 돈에서 수수료 1\%를 뗀다(2장).}
\end{figure}


\begin{itemize}[leftmargin=*]
\item \textbf{$s$는 줄 세우기에만 쓴다.} 결과는 실제로 받은 돈으로만 잰다. 예측이 틀렸다면 실제 IRR이 낮게 나와 그대로 드러난다.
\item 부도확률만 보지 않고 금리를 함께 보는 이유: LendingClub은 위험한 대출에 이미 높은 금리를 매긴다. 우리가 찾는 것은 위험에 비해 금리를 더 받는 대출이다.
\item 같은 달 안에서는 국채 금리가 모든 대출에 같으므로, 기준 수익률을 국채 대신 회사채로 바꿔도 줄 세우는 순서는 같다.
\end{itemize}

\section{근거: 승인비율별 결과}
\begin{table}[H]\centering\small
\begin{tabular}{@{}rcccc@{}}\toprule
 & \multicolumn{2}{c}{선별 가치: 무작위보다 더 번 \%p} & \multicolumn{2}{c}{Sharpe: 제안 $-$ 같은 비율 무작위}\\\cmidrule(lr){2-3}\cmidrule(l){4-5}
승인비율 & 값 (90\% 구간) & 더 번 달 & 값 (90\% 구간) & 판단\\\midrule
10\% & \textbf{+2.11} (1.80, 2.45) & 98\% & $-0.07$ ($-0.59$, 0.14) & 구분 안 됨\\
20\% & +1.74 (1.52, 1.99) & 100\% & +0.08 ($-0.47$, 0.40) & 구분 안 됨\\
30\% & +1.46 (1.29, 1.62) & 100\% & +0.23 ($-0.28$, 0.58) & 구분 안 됨\\
40\% & +1.19 (1.06, 1.33) & 100\% & +0.31 ($-0.12$, 0.61) & 구분 안 됨\\
50\% & +0.99 (0.84, 1.12) & 100\% & +0.28 ($-0.04$, 0.55) & 구분 안 됨\\
60\% & +0.82 (0.68, 0.94) & 100\% & +0.32 (0.08, 0.54) & \textbf{높음}\\
70\% & +0.63 (0.51, 0.74) & 100\% & +0.32 (0.15, 0.52) & \textbf{높음}\\
80\% & +0.48 (0.39, 0.56) & 100\% & +0.29 (0.15, 0.46) & \textbf{높음}\\
90\% & +0.28 (0.22, 0.33) & 100\% & +0.21 (0.14, 0.36) & \textbf{높음}\\\midrule
평균 & \textbf{+1.08} & & & \\\bottomrule
\end{tabular}
\caption{발행월 컷, 국채 기준, 수수료 1\%. 선별 가치는 승인 1달러당 연환산 IRR의 차이, 무작위는 10회 평균, 구간은 6개월 블록 부트스트랩 2000회(42단계, \ft{results/42_final_proposal_check.csv}). 평균은 무작위 반복 수에 따라 1.07--1.09\%p로 조금 다르다(38단계 1.07).}
\end{table}
\begin{enumerate}[leftmargin=*]
\item \textbf{모든 승인비율에서 무작위보다 더 벌었다.} 9개 승인비율 모두 구간이 0보다 위이고, 거의 모든 달에 더 벌었다.
\item \textbf{승인비율이 작을수록 1달러당 이점이 크다.} 많이 살수록 결국 거의 다 사게 되므로 차이가 줄어든다(80\%: +0.48\%p). 반대로 포트폴리오 전체 초과수익은 많이 살수록 크다(10\%: 0.65, 80\%: 4.34\%p, 위 상자).
\item \textbf{승인비율이 60\% 이상이면 덜 흔들리기도 한다.} 50\% 이하에서는 Sharpe가 무작위와 구분되지 않는다.
\item \textbf{돈의 출처:} 같은 등급 안에서 기대수익 점수 상위가 하위보다 연 +2.38\%p(90\% 구간 2.08--2.73), 48개월 모두 양수였다(7장).
\item \textbf{특정 모델 덕이 아니다:} XGBoost도 선별 가치가 같았고(1.03 vs 1.04\%p), 정답을 섞은 가짜 모델은 0 근처였다(7장).
\end{enumerate}

\section{부도확률 낮은 순(이전 최종안)과의 관계}
\begin{table}[H]\centering\small
\begin{tabularx}{\linewidth}{@{}>{\raggedright\arraybackslash}p{50mm}cc@{}}\toprule
 & \textbf{기대수익 순(제안)} & 부도확률 낮은 순\\\midrule
무작위보다 더 버나(선별 가치) & 9개 승인비율 모두 & 40\% 이상에서만, 평균 +0.20\%p\\
무작위보다 덜 흔들리나(Sharpe) & 60\% 이상에서 & 9개 승인비율 모두\\
80\%에서 Sharpe / 평균 초과수익 & 3.01 / 4.34\%p & 3.06 / 4.30\%p\\\bottomrule
\end{tabularx}
\end{table}
두 전략은 서로 다른 목표에서 강하다. 기대수익 순은 \textbf{더 벌기}, 부도확률 낮은 순은 \textbf{덜 흔들리기}다. 부도확률 낮은 순은 ``흔들림을 가장 줄이고 싶은 투자자''를 위한 대안으로 보고한다. 승인비율이 클 때(80\%)는 두 전략의 결과가 거의 같다.

\section{어떻게 여기까지 왔나}
\begin{enumerate}[leftmargin=*]
\item \textbf{v1(처음 규칙서):} 주 가설이 기대수익 순위의 선별 가치(H2)였다.
\item \textbf{v2.2(39단계):} ``과제가 Sharpe 최대화를 요구한다''는 전제에서 포트폴리오 Sharpe를 주 지표로 두고, 1-SE 규칙(1위 Sharpe에서 그 SE를 뺀 선 이상 중 승인비율 최대)으로 \textbf{부도확률 낮은 순 80\%}를 골랐다. Sharpe 1위는 10\%(3.82)였고 80\%는 3.06이었다.
\item \textbf{검토에서 나온 약점:} 1-SE 기준선 안의 후보를 ``1위와 통계적으로 구분되지 않는다''고 설명한 것은 틀렸다(10\% $-$ 80\% 차이 0.76, 구간 0.45--1.60). 또 V2의 Sharpe는 40--90\%에서 거의 평평해 80\%가 특별하지 않았다(90\%는 기준선보다 0.001 낮아 빠졌다). ``왜 80\%냐''에 답하기 어려웠다.
\item \textbf{목적 확정(2026-11-17, 18):} Sharpe 최고 튜닝은 목적이 아니고, 과제 공지를 다시 확인한 결과 Sharpe는 필수 지표가 아니었다. 그래서 투자자에게 더 직접적인 ``같은 양을 살 때 더 버는가''를 주 판정으로 두고, 처음 규칙서의 H2로 돌아갔다(test 규칙 v2.3, test를 열기 전 커밋과 태그).
\end{enumerate}
바꾸는 근거는 과제 해석이지만, train 결과를 모두 본 뒤의 결정이라는 사실은 기록한다. 판정 정의는 v1 그대로라 새로 만든 기준이 아니다.

\section{엄격 비판과 대응}
\begin{table}[H]\centering\small
\begin{tabularx}{\linewidth}{@{}>{\raggedright\arraybackslash}p{50mm}L@{}}\toprule
비판 & 점검 결과와 대응\\\midrule
``+1\%p''는 작은 승인비율이 끌어올린 평균이다 & 맞다. 10\%에서 +2.11, 80\%에서 +0.48\%p다(1달러당). 전체 돈의 초과수익은 반대로 10\%가 가장 작다(0.65\%p). 평균 하나가 아니라 승인비율별 표를 보인다. 많이 사는 투자자의 이점은 연 0.3--0.6\%p다.\\
승인비율이 작을 때는 더 흔들릴 수 있다 & 50\% 이하에서 Sharpe는 무작위와 구분되지 않고, 10\%는 점추정이 조금 낮다(2.55 vs 2.63). 금리가 높은 대출을 사기 때문이다. ``승인비율이 작을 때 덜 흔들린다는 근거는 없다''고 쓴다.\\
흔들림만 보면 부도확률 순이 낫다 & 맞다. 맞교환이며 표로 그대로 보인다.\\
$p$의 크기가 틀리면 순서가 바뀐다 & 모델 $p$는 실제보다 낮다(2013년 6.1\% vs 12.2\%). $p$를 실제 수준으로 맞춰 다시 골라도 승인 대상의 88--92\%가 같았다(5장).\\
결과를 보고 전략을 바꾼 것 아닌가 & H2는 v1의 주 가설이었고 정의도 그대로다. 바꾼 근거는 과제 해석이다. train 결과를 본 뒤라는 사실은 기록했다.\\
Sharpe 약 3은 너무 크다 & 만기 보유로 가격 변동이 없고, 2013--16년에 경기침체가 없어 부도가 몰리지 않았다. 우리 전략끼리 비교하는 데만 쓴다.\\
불경기에는? & 직접 확인하지 못했다. 팀 추가 점검(역 스트레스, 저장소 반영 전)에서는 부도가 모든 대출에서 고르게 늘면 약 6배까지 이점이 유지되지만, 모델이 안전하다고 본 절반의 부도만 약 2배가 되면 이점이 사라졌다. 이 제안은 ``불경기에도 모델의 위험 순서가 유지된다''는 가정에 기댄다.\\
살 수 있는 물량 & 개인이 조각으로 살 수 있는 대출 비중이 2013년 72\%에서 2016년 27\%로 줄었다. 상위부터 원하는 만큼 살 수 있다고 가정했다.\\\bottomrule
\end{tabularx}
\end{table}

\section{기준 수익률을 바꾸면}
기준을 국채보다 높게 잡아도 평균 초과수익 1위는 네 경우 모두 기대수익 순이다(38단계). 회사채(BAA) 기준에서는 Sharpe도 기대수익 순이 높다(승인 30\%: 1.29, 부도확률 순 0.92, 무작위 0.52). 부도확률이 낮은 대출은 금리도 낮아서, 기준이 오르면 남는 폭이 작아 조금만 나빠도 기준 아래로 떨어지기 때문이다.

\section{불확실성은 어디까지 셌나}
\begin{itemize}[leftmargin=*]
\item \textbf{센 것:} 기간 변화. 48개월을 6개월 이동 블록으로 다시 뽑았다(2000회, 모든 전략 공통 인덱스, 짝지은 차이). 그 표집 방식 아래의 불확실성일 뿐, 새로운 경기침체에서도 유지된다는 보증이 아니다.
\item \textbf{줄였지만 따로 안 센 것:} 무작위 선택(10회 평균).
\item \textbf{안 센 것:} 모델 seed(3개 평균), 대출 단위 표본오차.
\item 모든 숫자는 train 성과이므로 낙관적으로 평가됐을 가능성이 있다.
\end{itemize}

\chapter{한계와 test 계획}

\section{남은 한계}
\begin{itemize}[leftmargin=*]
\item \textbf{현금흐름은 근사다.} 월별 납입 기록이 없어 받은 총액을 월 상환액으로 나눠 정상 납입 개월을 추정했다(2장).
\item \textbf{기대수익 점수는 1년 단위 근사다.} 조기상환과 부도 시점을 무시한다. 부도와 조기상환을 함께 모형화한 연구(Li 외, 2023, European Journal of Operational Research)가 개선 방향이다.
\item \textbf{안정성은 좋은 시절의 출렁임이다.} 2013--16년은 경기침체가 없었다. 위기 시기(2007--09)는 표본이 작고 A--D 등급만 볼 수 있었다(A가 가장 안정적이었으나 수준은 낮았다).
\item \textbf{소득은 자기 신고 값이고, 그 신호는 사라졌다.}
\item \textbf{개인이 살 수 있는 물량이 줄었다.} 조각 대출 비중이 2013년 72\%에서 2016년 27\%로 줄었다.
\item \textbf{같은 데이터를 오래 봤다}(8장).
\end{itemize}

\section{test 규칙 v1·\allowbreak{}v2 (test 전에 고정, git 태그)}
v1은 그대로 실행하고, v2는 30--34단계에서 나온 후보를 더한다(\ft{docs/test_protocol_v2.md}, 태그 \ft{test-protocol-v2}).

\begin{table}[H]\centering\small
\begin{tabularx}{\linewidth}{@{}p{45mm}L@{}}\toprule
가설 & 지지 기준(주 표본: test의 2013--16년 36개월 종결 대출)\\\midrule
H1 성과 귀속(v1) & 세부등급 분해에서 선별 > 0, 선별 > 배분, 3개 연도 이상 선별 > 0\\
H2 선별 가치(v1) & 기대수익 순위 선별 가치의 95\% 구간 하한 > 0\\
H3 국채 대비(v1) & 전부 승인이 발행월 90\% 이상에서 국채보다 높음\\
H4 회사채 대비(v1) & 선별하면 회사채 이하 발행월이 줄어듦\\
H5 등급 안 신호 & 기대수익 등급 중립 롱숏 90\% 구간 하한 > 0\\
H6 두 축 & 승인 30\%·\allowbreak{}50\% 모두 부도확률 순이 Sharpe 높고 기대수익 순이 수익 높음\\
H7 단순 규칙 & 계좌 개설 규칙이 부도확률 순을 Sharpe와 수익 둘 다에서 이김\\
H8 Sharpe 정의 & 세 정의 모두 부도확률 순 > 기대수익 순\\\bottomrule
\end{tabularx}
\end{table}

\paragraph{해석 규칙(미리 정함).}
\begin{enumerate}[leftmargin=*]
\item 만기 도달 비율이 0.9 미만인 연도는 따로 보고하고 근거로 쓰지 않는다(종결 편향, 3장).
\item 60개월은 36개월과 섞지 않는다.
\item H5가 실패하면 ``LC 가격이 효율화되었다''를 대안 결론으로 쓴다(연소득 신호가 사라진 것과 같은 방향).
\item 크기보다 방향과 구간으로 판정한다.
\item H7이 실패하면 ``train에서 보고 고른 규칙은 과대평가된다''를 그대로 보고한다.
\end{enumerate}

\section{test에서 무엇을 기대하나(판단, 계산 아님)}
\begin{table}[H]\centering\small
\begin{tabularx}{\linewidth}{@{}p{40mm}p{22mm}L@{}}\toprule
주장 & 유지 가능성 & 근거와 위험\\\midrule
등급 안 신호(H5, H2) & 높음 & 48개월 전부 양수, 해마다 커짐, 여러 조건에서 유지. 위험: test가 우리가 안 본 시기에 몰린 경우\\
두 축(H6, H8) & 중상 & 9개 승인비율·\allowbreak{}두 수익률 정의에서 같은 방향. 위험: 60개월 비중, 종결 편향\\
효과 크기 1\%p & 중간 & 같은 기간 안의 결과. 시기가 바뀌면 크기는 쉽게 변함\\
단순 규칙(H7) & 중간 & 사후 선택이라 과대평가 가능\\\bottomrule
\end{tabularx}
\end{table}

\section{발표에서의 한 문장}
\begin{keybox}[한 문장]
``돈은 같은 등급 안에 있다. 기대수익으로 고르면 가장 많이 벌고, 부도 위험이 낮은 쪽으로 고르면 과제의 Sharpe가 가장 높으며, 둘의 차이는 고위험 등급이 위험만큼 금리를 받지 못한 데서 온다. 우리는 이 결론을 미리 정한 규칙으로 test에서 한 번만 확인한다.''
\end{keybox}

\appendix
\chapter{재현 방법과 파일 지도}
\section{실행 순서}
\begin{lstlisting}[language=bash]
python scripts/prepare_data.py              # raw CSV -> parquet
pytest tests                                # hand-checked tests
python scripts/25_paired_bootstrap.py       # selection value and CI
python scripts/26_resolution_bias.py        # resolution-bias simulation
python scripts/30_portfolio_sharpe.py       # two axes
python scripts/33_sortino_npv_longshort.py  # long-short, grade Sharpe
python scripts/34_blind_spots.py            # six blind spots
python scripts/36_p_reliability.py          # p reliability tables
python scripts/35_test_evaluate_v2.py <test file>   # official test, v1 + v2
\end{lstlisting}

\section{어디에 무엇이 있나}
\begin{table}[H]\centering\small
\begin{tabularx}{\linewidth}{@{}p{55mm}L@{}}\toprule
파일 & 내용\\\midrule
\ft{docs/decisions.md} & 단계별로 무엇을 왜 골랐나(1--35단계)\\
\ft{docs/research_log.md} & 주장 장부(C1--C35), 실패, 정정, 가정, 데이터 엿보기 기록\\
\ft{docs/test_protocol.md}, \ft{test_protocol_v2.md} & test 전에 고정한 규칙\\
\ft{docs/컬럼사전_최신.xlsx} & 141개 컬럼의 처리, 결측, IV, 중요도\\
\ft{docs/notes/LC_lecture_notes.pdf} & 기초 강의노트\\
\ft{results/} & 모든 실험의 CSV와 로그\\\bottomrule
\end{tabularx}
\end{table}

\chapter{용어 한 줄 사전}
{\small
\begin{longtable}{@{}p{40mm}p{110mm}@{}}\toprule
용어 & 뜻\\\midrule\endhead
$p$(부도확률) & 대출이 만기까지 가는 동안 결국 부도로 끝날 확률\\
LGD(손실률) & 부도 시 원금 중 잃는 비율의 평균\\
기대수익 점수 & $(1-p)\times$금리 $-\ p\times$LGD $-$ 국채\\
IRR & 넣고 돌려받은 돈의 흐름을 하나의 연 이자율로 바꾼 값\\
코호트 & 같은 달에 발행된 대출 묶음\\
시점 학습 & 그때 알 수 있던 정보로만 학습\\
누수 & 결정 시점에 알 수 없던 정보가 모델에 섞이는 것\\
선별 가치 & 같은 비율로 무작위로 고를 때보다 더 번 연 수익률(\%p)\\
Brinson 분해 & 초과수익을 배분, 선별, 상호작용으로 나눔\\
등급 중립 롱숏 & 등급마다 상위 30\% $-$ 하위 30\%의 수익 차이. 등급 안 신호의 크기\\
포트폴리오 Sharpe & 매달 (승인 대출 + 국채) 초과수익의 평균 \ensuremath{\div} 표준편차\\
건별 Sharpe & 대출 한 건 단위 초과수익의 평균 \ensuremath{\div} 표준편차\\
Sortino & 아래쪽 출렁임만 위험으로 보는 Sharpe\\
NPV & 국채 금리로 할인한 순현재가치(달러)\\
하위 5\% 평균 & 가장 나쁜 5\% 대출 수익률의 평균(꼬리 손실)\\
종결 편향 & 이미 끝난 대출만 보면 빨리 끝난(빨리 부도 난) 대출이 과대 대표됨\\
AUC & 부도 하나와 정상 하나를 뽑았을 때 부도에 더 높은 점수를 줄 확률\\
신뢰도 표 & 예측 확률 구간별로 실제 비율과 비교한 표\\
위약 검정 & 효과가 없어야 할 조건(정답 섞기)에서 효과가 0인지 확인\\
사전 등록 & 결과를 보기 전에 가설과 판정 기준을 고정\\\bottomrule
\end{longtable}}

\chapter{질문과 답}
\newcommand{\QA}[2]{\begin{qa}{Q. #1}#2\end{qa}}

팀 스터디(2026년 11월 16일)에서 보고서를 함께 읽으며 나온 질문과 답이다. 발표와 질의응답 대비용이다. 숫자는 모두 train(2013--16년 36개월 대출) 결과다.

\section{전체 구조}
\QA{결국 우리 과정은 어떻게 흘러가나?}{
\begin{enumerate}[leftmargin=*]
\item 대출마다 부도확률 $p$를 구한다. 모델은 그 시점에 알 수 있던 과거 대출로만 학습한다.
\item 점수로 줄을 세운다. 기대수익 순(V1)은 $s=(1-p)r-p\cdot\mathrm{LGD}-tb$, 부도확률 순(V2)은 $p$가 낮은 순이다.
\item 발행월마다 그 달 대출 중 상위 $q$를 승인한다.
\item 그 달 승인 대출의 현금흐름을 합쳐 발행월 묶음 IRR을 하나 구한다(48개월이면 48개).
\item 매달 초과수익 $=$ (묶음 IRR $-$ 같은 달 같은 만기 국채) \ensuremath{\times} 승인 금액 비중.
\item 48개 초과수익의 평균 \ensuremath{\div} 표준편차 $=$ 포트폴리오 Sharpe.
\item 같은 비율 무작위 선택과 비교해 고르는 실력이 있는지 본다.
\end{enumerate}}

\QA{IRR을 Sharpe의 분자에 넣는 것인가? ``IRR $-$ $s$''를 계산하나?}{IRR 자체가 아니라 (IRR $-$ 국채) \ensuremath{\times} 비중이 분자에 들어간다. $s$는 승인할 때 쓰는 예상 점수(고르는 도구)이고, IRR은 끝난 뒤 실제 현금흐름으로 잰 수익률(채점 도구)이다. 둘을 빼거나 섞지 않는다. 예상 점수로 학생을 뽑고 실제 점수로 그 선택을 채점하는 구조다.}

\QA{기대수익 점수 $s$는 우리가 만든 식인가?}{아니다. 문헌의 기대수익 점수(expected return score)다. 안 망하면 받는 이자, 빼기 망하면 잃는 원금, 빼기 국채 기회비용이다. 1년 단위 근사라 정확한 예상 수익이 아니라 순위용으로만 쓴다.}

\QA{IRR만 보면 되지 왜 Sharpe가 필요한가?}{평균 초과수익만 키우면 답이 ``전부 승인''(4.86\%p)이다. 48개월 내내 대출이 국채보다 나았으니 많이 살수록 평균이 커지고, 부도 위험은 전혀 보지 않게 된다. Sharpe는 흔들림(위험)을 분모에 넣어 준다. 반대로 Sharpe만 최대로 하면 ``가장 안전한 10\%''(3.82)가 답이라 돈의 90\%가 논다. 그래서 Sharpe를 주 지표로, 평균 초과수익을 두 번째 축으로 본다.}

\QA{우리 목표는 무엇인가? Sharpe 극대화인가 수익 극대화인가?}{권하는 해석은 \textbf{``Sharpe를 무작위와 전부 승인보다 확실히 높이면서 수익 규모를 크게 잃지 않는 것''}이다.
\begin{itemize}[leftmargin=*]
\item 평균 초과수익 극대화: 전부 승인. 선별이 필요 없고 위험을 보지 않는다.
\item Sharpe 엄격 최대화: 부도확률 순 10\%. 규모가 작고, 흔들림이 거의 0이라 숫자가 불안정하며, 18개 중 최대값이라 운이 섞였다.
\item Sharpe 개선 $+$ 규모 유지: 부도확률 순 80\%(3.06, 평균 4.30\%p). 무작위(2.72), 전부 승인(2.71)보다 높고 규모는 전부 승인의 약 90\%다.
\end{itemize}
test 판정 H10(같은 비율 무작위보다 Sharpe가 높은가)도 이 목표와 맞는다. 다만 80\%는 v2.2 규칙이 고른 \textbf{현재 후보}이며, 팀이 Sharpe 엄격 최대화로 정하면 test 전에 v2.3으로 다시 고정한다.}

\QA{LC에서 왜 ``포트폴리오''로 보나?}{LC는 플랫폼이다. 투자자는 대출 하나를 통째로 빌려주지 않고 여러 대출에 조각(25달러 단위)으로 나눠 투자한다. LC는 투자자 수수료 1\%를 뗀다. 그래서 대출 묶음의 성과가 투자자 입장의 성적이다.}

\section{요약 문장}
\QA{``매년 약 1\%p 더 번다''는 얼마인가?}{1\%p는 수익률 0.01이다. 연 5\%가 6\%가 되는 정도이고, 1억 원을 1년 굴리면 약 100만 원 차이다. 95\% 구간 0.95--1.11은 운을 감안해도 그 범위라는 뜻이다.}

\QA{롱숏은 등급마다 하나, 전체를 $p$로 줄 세우나?}{\textbf{등급마다 따로} 한다. 매달, 등급마다 기대수익 점수 상위 30\%의 IRR에서 하위 30\%의 IRR을 뺀다. 전체를 줄 세우면 위쪽이 거의 A등급이라 ``좋은 등급을 산 효과''가 섞인다. 등급 안에서 빼야 고르는 실력만 남는다(+2.38\%p, 48개월 전부 양수). LC 대출은 공매도가 안 되므로 실제 전략이 아니라 측정 도구다.}

\QA{회사채(BAA)는 애플 같은 안 망할 회사인가? 국채는 몇 년물인가?}{BAA는 Moody's 투자등급 중 가장 낮은 등급(Baa) 회사채 수익률 지수다. AAA보다 위험하고 금리가 높다. 그래서 ``투자자의 요구수익률이 높을 때''의 허들로 쓴다. 국채는 대출 만기에 맞춰 36개월 대출에 3년물(GS3)을 쓴다.}

\QA{``새 계좌를 적게 연''은 무슨 뜻인가?}{\ft{acc_open_past_24mths}: 신청자가 최근 24개월 동안 새로 연 신용 계좌(카드, 대출 등) 수다. 짧은 기간에 계좌를 많이 열면 돈이 급하다는 신호라 부도와 연결된다. 이 규칙(V3)은 데이터를 보고 찾은 사후 규칙이라 주 제안에서 빼고 test에서 확인할 도전자로만 둔다.}

\QA{등급은 $p$ 계산에만 쓰이나?}{세 곳에 쓰인다. (1) $p$ 모델의 입력(등급, 세부등급, 금리). (2) 롱숏과 V3의 등급 안 비교. (3) 성과 분해(Brinson). V1, V2의 승인 컷은 등급을 나누지 않고 발행월 전체에서 자른다.}

\section{수익률과 IRR}
\QA{2장 예시(1만 달러, 332.14 \ensuremath{\times} 10, 회수 500)는 누구 입장인가?}{투자자(빌려준 쪽)다. 0개월에 10,000을 내주고 1--10개월에 월 상환액 332.14를 받다가 15개월에 회수금 500을 받았다. 이 현금흐름의 IRR이 연 $-76.8\%$다.}

\QA{단순 연환산 $-27.4\%$의 36은 어디서 왔나?}{대출 만기 36개월이다. $(3821.43/10000)^{12/36}-1$로, 실제로는 15개월 만에 끝난 손실을 3년에 걸쳐 늘려서 작아 보인다. 지난 우승 논문의 계산 방식이며, 비판할 때는 논문 출처를 함께 적는다.}

\QA{우리 IRR은 코호트 IRR인가? 근거는?}{발행월 묶음(코호트) IRR이다. 같은 달 승인 대출의 현금흐름을 먼저 더한 뒤 IRR을 한 번 구한다. 근거는 다음과 같다.
\begin{itemize}[leftmargin=*]
\item 투자자 계좌가 실제로 번 수익률이다(금액 가중). 건별 IRR의 평균은 1,000달러와 30,000달러 대출을 똑같이 센다.
\item 건별 IRR은 부도 시 $-77\%$처럼 극단값이 커서 평균이 몇 건에 휘둘린다.
\item 돈이 들어온 시점(조기상환, 부도)을 반영한다.
\item 매달 승인한다는 결정 단위와 맞고, 48개 값이 생겨 흔들림을 잴 수 있다.
\end{itemize}}

\QA{IRR을 구할 때 금리가 필요한가? 명목인가 실질인가?}{IRR 자체에는 시장 금리가 필요 없다. 현금흐름만으로 현재가치 합이 0이 되는 할인율을 역산한다. 모든 값이 명목이고 허들(국채)도 명목이라 초과수익에서 인플레이션은 지워진다.}

\QA{조기상환은 손해인가? 금리가 오르면 오히려 이득 아닌가?}{맞는 지적이다. 조기상환의 손익은 재투자 금리에 달려 있다. 금리가 오르면 일찍 받은 돈을 더 높은 금리에 다시 넣을 수 있어 이득이다. IRR에는 같은 수익률로 재투자한다는 가정이 암묵적으로 들어 있어 이 효과를 따로 보지 않았다. 2013--16년은 저금리가 안정적이라 영향이 작았을 것으로 보지만 한계로 적었다(10장).}

\QA{현금흐름이 추정이라는 게 왜 한계인가? 월 상환액은 일정한데.}{월 상환액은 일정해도 \textbf{실제로 언제 냈는지}는 모른다. 데이터에는 받은 총액과 마지막 납부월만 있다. 연체 후 몰아서 낸 경우, 조기상환 시점, 부도 전 부분 납부는 추정이다. IRR은 시점에 민감하므로 한계로 적었다.}

\section{데이터와 표본}
\QA{결과 라벨(완제, 조기상환, 부도)은 어떻게 만들었나?}{원본 상태는 Fully Paid(완제)와 Charged Off(부도)다. 완제 가운데 만기 전에 끝난 것을 조기상환으로 구분해 분석에 썼다. 진행 중(Current, Late)은 결과를 모르니 학습과 채점에서 뺀다.}

\QA{종결 편향은 무엇인가?}{발행 후 1--3년밖에 안 된 대출 중 이미 끝난 것만 보면, 빨리 망한 대출은 이미 결과가 나와 있고 정상 상환자는 아직 갚는 중이다. 그래서 부도가 실제보다 많이 잡히고, ``부도만 피하는'' 규칙이 똑똑해 보인다. 그래서 만기 도달 비율이 낮은 연도는 따로 보고한다.}

\QA{생존분석은 왜 안 썼나?}{생존분석은 아직 안 끝난 대출을 버리지 않고 쓰는 정석 방법이다. 우리는 시점 학습과 종결 대출만 쓰는 방식을 택했고, 진행 중 대출을 활용한 실험은 효과가 없었다. 시간과 복잡도 때문에 쓰지 않았고, 한계 겸 추후 연구로 적었다(10장).}

\QA{MNAR, WOE는 무엇인가?}{MNAR(missing not at random)은 값이 비어 있는 것 자체가 정보라는 뜻이다. 예를 들어 종결 여부는 부도 여부와 관련이 있다. WOE(weight of evidence)는 변수의 각 구간이 정상과 부도를 얼마나 가르는지 로그 비율로 나타낸 값이다(4장).}

\section{부도확률 $p$}
\QA{FICO는 무엇이고 왜 690점인가?}{FICO는 미국 개인 신용점수(300--850)다. 신청 시점 값이라 누수가 아니다. 690은 설명용 예시이고, 실제 기준점은 모델이 데이터에서 찾는다. 나무가 수백 그루라 기준점도 나무마다 다르다.}

\QA{$p$는 왜 실제보다 낮게 나왔나?}{학습 시점(예: 2013년 1월 1일)에 쓸 수 있는 직전 해 발행분은 이미 끝난 대출뿐이다. 부도는 연체 후 확정까지 오래 걸리므로, 1년 안에 끝난 대출은 대부분 조기상환이다. 모델은 ``요즘 대출은 안전하다''고 배워 $p$가 절반쯤 낮게 나왔다(평균 6.4\% vs 실제 12.2\%). 일찍 하교한 학생만 보고 반 전체가 우등생이라고 판단한 셈이다.}

\QA{종결 편향(부도 과대)과 방향이 반대 아닌가?}{대출의 나이가 다를 뿐 둘 다 사실이다. 나이 1년 미만의 끝난 대출에는 조기상환이 많고(부도는 아직 확정 전), 1--3년의 끝난 대출에는 부도가 많다(정상 상환은 아직 만기 전). $p$를 낮춘 것은 앞의 것이다.}

\QA{$p$가 틀렸는데 왜 결과는 괜찮았나? 오차 상쇄는 구조적 설명인가?}{상쇄는 사후 관찰이지 보장된 구조가 아니다. 직접 근거는 두 가지다. (1) V2는 $p$의 순서만 쓰므로 크기가 같은 비율로 틀리면 결과가 그대로다. (2) V1은 금리가 함께 들어가 $p$의 크기가 순서를 일부 바꾸므로, $p$를 실제 수준으로 맞춰 다시 골랐다. 승인 대상의 88--92\%가 같았다. 순서까지 망가질 만큼 틀리면 문제다(CatBoost 2016년, 평균 3.6\% vs 실제 15.9\%).}

\QA{$p$를 왜 10분위로, 등급 안에서는 왜 셋으로 나눴나?}{10분위는 확률 신뢰도를 볼 때 쓰는 관례다(분위마다 약 5.2만 건). 등급 안 3분할은 등급마다 $p$ 하위, 중간, 상위 1/3이다. 분할 수는 임의이며, 결론은 ``같은 등급 안에서도 $p$가 높을수록 실제 부도율이 오른다''는 방향이다.}

\QA{AUC, seed, 부트스트랩은 무엇인가?}{AUC는 부도인 사람을 정상인 사람보다 위험하게 줄 세울 확률이다(0.5는 동전, 1은 완벽, 우리 0.63--0.68). seed는 무작위 계산의 출발값이고, 바꿔도 결과가 유지되는지 본다. 부트스트랩은 가진 표본을 다시 뽑아 ``운이 달랐다면''을 수천 번 흉내 내 구간을 만드는 방법이다. 우리는 6개월 덩어리로 다시 뽑아 이웃 달의 상관을 보존한다.}

\section{비교와 검증}
\QA{왜 ``같은 비율 무작위''와 비교하나? 무작위는 전체에서인가 등급 안에서인가?}{발행월 전체에서 같은 비율을 무작위로 고른다. 같은 달, 같은 비율이라 그달 시장 상황이 양쪽에 똑같이 들어가 지워지고, 차이에는 ``무엇을 골랐나''만 남는다.}

\QA{``정답을 섞은 가짜 모델은 0''은 무슨 뜻인가?}{부도 여부(정답)를 무작위로 섞어 학습시키면 모델은 아무것도 배울 수 없다. 그 모델로 고르면 선별 가치가 0 근처여야 하고, 실제로 그랬다. 평가 코드가 가짜 성과를 만들지 않는다는 점검이다.}

\QA{선별 가치의 ``9개 비율 평균''의 9는?}{승인 10\%, 20\%, \ldots, 90\%의 9개다. 유리한 비율 하나만 고르지 못하게 평균을 낸다.}

\section{Sharpe}
\QA{``건별''이란 무엇인가? 세 정의는 어떻게 다른가?}{건별은 대출 한 건을 데이터 한 줄로 본다는 뜻이다.
\begin{itemize}[leftmargin=*]
\item 건별, 거절 $=0$: 승인하지 않은 대출을 0으로 센다. 승인비율만 줄여도 값이 작아진다.
\item 건별, 승인만: 승인 대출만 본다. 0.10--0.34로, 조교가 말한 0.2--0.4와 크기가 맞는다.
\item 발행월 포트폴리오(주 지표): 한 달을 한 줄로 본다. 약 2--3.5.
\end{itemize}
어떤 정의로 채점될지는 따로 확인하지 않는다(평가가 과정과 논리 중심이므로). 정의를 왜 골랐는지 밝히고 건별 값을 참고로 함께 적는다.}

\QA{``흔들림도 0.3배''는 무슨 뜻인가? $\sqrt{0.3}$ 아닌가?}{모든 값에 0.3을 곱하면 표준편차도 정확히 0.3배다. 예: 10, 4, 7\% \ensuremath{\to} 3, 1.2, 2.1\%에서 오르내리는 폭이 6\%p \ensuremath{\to} 1.8\%p. 그래서 평균 \ensuremath{\div} 표준편차에서 0.3이 지워진다. $\sqrt{\ }$은 \textbf{대출 건수}가 바뀔 때의 효과다. 묶음에 대출이 적으면 개별 부도 운이 덜 희석되어 묶음 IRR이 더 흔들린다(대략 $1/\sqrt{\text{건수}}$). 그래서 승인 10\%의 Sharpe가 운에 더 민감하다(무작위 seed 간 표준편차 10\%: 0.23, 80\%: 0.02).}

\QA{포트폴리오 Sharpe가 왜 이렇게 크게(2--3.5) 나오나? 공격받지 않을까?}{세 가지 이유다. (1) 대출은 만기까지 보유하고 시장 가격으로 평가되지 않아, 주식처럼 가격이 오르내리는 흔들림이 없다. (2) 수백 건을 묶으면 개별 부도가 희석된다. (3) 2013--16년은 경기침체가 없었다. 그래서 주식 Sharpe(0.5 안팎)와 비교하면 안 되고, 우리 전략끼리 순위를 매기는 숫자로만 쓴다. 발표에서는 값 자체를 내세우지 않고 ``같은 비율 무작위보다 확실히 높다''까지만 주장한다.}

\QA{1년 단위로 묶으면 더 안정적이지 않나?}{1년 단위면 값이 4개(2013--16)뿐이라 표준편차를 믿을 수 없다. 또 묶는 기간이 길수록 흔들림이 작게 잡혀 위험이 과소평가된다. 월 단위는 실제 결정 단위(매달 승인)와도 맞는다.}

\QA{왜 연도 단위가 아니라 발행월마다 자르나?}{매달 승인 비중이 비슷해야 비중이 Sharpe에서 지워진다. 연도 단위로 자르면 고금리 대출이 몰린 달은 많이, 적은 달은 적게 승인되어 비중이 들쭉날쭉하고 Sharpe가 깎인다(2.46 vs 2.91). 12월 대출까지 본 뒤 1월 승인을 정하는 셈이라 실전에서도 불가능하다.}

\QA{Sortino는 무엇인가?}{Sharpe의 분모를 아래로 떨어진 흔들림만으로 바꾼 지표다. 건별로는 순위가 Sharpe와 같아 새 정보가 없었다. 포트폴리오로는 48개월 모두 국채보다 높아 아래쪽 흔들림이 0이라 계산되지 않는다.}

\section{팀 내 비판과 앞으로}
\QA{``좋은 시절만 봤다''는 한계는 어떻게 보완하나? 코로나 시기를 쓰면?}{데이터에 2017--20년 발행분이 있어 2017--19년 대출이 2020년 코로나 시기를 진행 중으로 겪었다. 다만 대부분 아직 진행 중이고, 데이터는 2020년 어느 시점의 상태이며, LC는 2020년에 투자자용 상품을 줄이고 닫았다. 그래서 Sharpe 비교가 아니라 \textbf{보조 스트레스 점검}으로만 권한다. 질문은 ``나쁜 시기에도 점수 상위 대출이 덜 연체·\allowbreak{}부도났나''이고, 판정 기준을 먼저 적는다.}

\QA{다른 팀원이 제안한 단순 규칙과 우리 규칙은 어떻게 다른가?}{그 규칙의 정확한 정의(부도확률 순서인지, 다른 변수인지)를 받아 V2, V3와 같은 표에서 비교해야 한다. 아직 하지 않았다.}

\QA{Brinson, 롱숏 같은 부가 분석은 본문에 다 넣어야 하나?}{큰 줄기만 본문에 두고, 이해가 덜 되거나 시간이 부족하면 부록이나 추후 연구로 돌린다. 롱숏은 핵심 숫자 하나(+2.38\%p, 48개월 전부 양수)만 본문에 남겨도 충분하다.}

\QA{특정 숫자를 집어서 발표해도 되나?}{된다. 다만 (1) train 숫자임을 표시하고, (2) 1등 숫자 하나만 앞세우지 말고 규칙과 함께 보이며, (3) 구간을 함께 쓴다. 예: ``Sharpe 3.06''보다 ``약 3, 무작위보다 +0.34(0.30--0.55)''.}

\part{기술 부록: 기법, 코드, 치트시트}
\chapter{기술 부록 읽는 법}
\begin{warnbox}[먼저 읽기]
이 부록은 따로 만들었던 기술 문서 네 개를 합친 것이다. 기법 정리(피처, 데이터와 검증 기법, 수익률, 모델, 평가 지표), 코드 해설(핵심 코드를 한 줄씩), 개념 치트시트, 코드 치트시트다. 작성 시점이 다르므로(2026년 11월 9--16일) \textbf{최종 제안과 test 판정은 본문(핵심 정리, 9장, 10장)이 우선한다.} 부록에 나오는 1-SE 규칙과 ``부도확률 순 80\%''는 이전 최종안(규칙서 v2.2)의 설명이다. 현재 최종 제안은 기대수익 순(예산 무관), 주 판정은 v1의 H2다(규칙서 v2.3).
\end{warnbox}
\begin{table}[H]\centering\small
\begin{tabularx}{\linewidth}{@{}p{45mm}L@{}}\toprule
부분 & 내용\\\midrule
기법 정리 & 피처 101개의 뜻과 처리, 시점 학습과 결측, 수익률과 비교 기준, 모델과 보정, 평가 지표(AUC, Brier, Sharpe 등), 검증 방법론\\
코드 해설 & \ft{src/lc/}와 \ft{scripts/}의 핵심 코드를 한 줄씩: 데이터, 수익률(IRR), 모델, 선별 곡선, 포트폴리오 Sharpe와 부트스트랩\\
개념 치트시트 & 신용 대출 선별 전략을 검증하는 개념을 두 단 한 쪽으로\\
코드 치트시트 & 복사해서 쓰는 코드 조각(\ft{docs/notes/code/snippets.py})\\\bottomrule
\end{tabularx}
\end{table}

\chapter{한눈에 보기}

\section{과제와 질문}
과제는 투자자 관점에서 LendingClub(LC) 대출 중 무엇을 승인할지 정하는 것이다. 기본 지표는 Sharpe ratio지만 변형이 허용되고, accuracy보다 과정과 논리가 중요하다.

\begin{keybox}[이 문서의 결론 세 줄]
\begin{enumerate}
\item 2013\textasciitilde 2016년 36개월 대출은 전부 승인해도 수수료 차감 IRR 5.73\%로 국채(0.87\%)를 이겼다.
\item 모델의 추가 수익은 좋은 등급을 고르는 데가 아니라 \textbf{같은 등급 안에서} 더 나은 대출을 고르는 데서 나온다(Brinson decomposition).
\item Feature engineering, tuning, 결측 처리는 결과를 거의 바꾸지 못했다. 차이를 만든 것은 \textbf{무엇을 재고 어떻게 검증하느냐}였다.
\end{enumerate}
\Intuit{모델을 더 똑똑하게 만드는 것보다, 성적표를 정직하게 매기는 법이 더 중요했다.}
\end{keybox}

\section{데이터와 평가 설정}
\begin{table}[H]\centering\small
\begin{tabularx}{\linewidth}{@{}lL@{}}\toprule
항목 & 내용\\\midrule
데이터 & train 1,755,294건 \ensuremath{\times} 141 columns, 2007\textasciitilde 2020년 발행\\
결과 상태 & Fully Paid 90만, Current 62만, Charged Off 22만. 종결 표본(정책 미충족 제외) 1,116,156건\\
주 평가 표본 & 2013\textasciitilde 2016년 발행 36개월 종결 대출(종결 비율 99.99\%), 발행월 48개\\
학습 방식 & as-of(시점) 학습: 연도 Y의 대출은 Y년 1월 1일까지 확정된 대출로만 학습\\
수익률 & 발행월 cohort IRR, 투자자 수수료 1\% 차감\\
비교 기준 & 발행월\ensuremath{\cdot}만기를 맞춘 미 국채(GS3, GS5), 보조로 회사채 BAA\\\bottomrule
\end{tabularx}
\caption{기본 평가 설정.}
\end{table}

\begin{center}
\begin{tikzpicture}[node distance=5mm,box/.style={draw=black!70,rounded corners=1.2mm,align=center,font=\small,minimum height=10mm,text width=24mm}]
\node[box] (a) {데이터 정리\\누수 제거};
\node[box,right=of a] (b) {현금흐름 복원\\IRR, 국채};
\node[box,right=of b] (c) {as-of 학습\\PD, LGD};
\node[box,right=of c] (d) {승인 규칙\\기대수익};
\node[box,below=8mm of c] (e) {평가\\cohort, Brinson};
\node[box,below=8mm of b] (f) {검증\\placebo, bootstrap};
\draw[-{Stealth},black!60] (a)--(b); \draw[-{Stealth},black!60] (b)--(c); \draw[-{Stealth},black!60] (c)--(d);
\draw[-{Stealth},black!60] (d) |- (e); \draw[-{Stealth},black!60] (e)--(f);
\end{tikzpicture}
\end{center}

% ======================================================================
\chapter{피처 정리}

\section{컬럼 지도}
141개 컬럼을 의미와 시점으로 나누면 아래와 같다. 판별 질문은 하나다: ``대출을 승인하던 날 이 값을 알 수 있었나?''

\begin{table}[H]\centering\small
\begin{tabularx}{\linewidth}{@{}lrL@{}}\toprule
그룹 & 수 & 대표 컬럼과 역할\\\midrule
대출 조건 & 15 & \ft{int_rate}, \ft{grade}, \ft{sub_grade}, \ft{term}, \ft{loan_amnt}, \ft{purpose}. LC가 매긴 가격 정보 포함\\
차주 기본 정보 & 8 & \ft{annual_inc}, \ft{dti}, \ft{emp_length}, \ft{home_ownership}, \ft{addr_state}\\
신용 이력: 연체\ensuremath{\cdot}공공기록\ensuremath{\cdot}조회 & 24 & \ft{fico_range_low}, \ft{inq_last_6mths}, \ft{delinq_2yrs}, \ft{mths_since_last_delinq}\\
신용 이력: 리볼빙\ensuremath{\cdot}bankcard & 21 & \ft{revol_util}, \ft{bc_util}, \ft{total_bc_limit}, \ft{percent_bc_gt_75}\\
신용 이력: 계좌 수\ensuremath{\cdot}잔액\ensuremath{\cdot}한도 & 15 & \ft{acc_open_past_24mths}, \ft{total_acc}, \ft{tot_cur_bal}\\
신용 이력: 할부 & 9 & \ft{il_util}, \ft{num_il_tl}, \ft{total_il_high_credit_limit}\\
공동신청 & 15 & \ft{sec_app_*}, \ft{annual_inc_joint}. 공동신청일 때만 값이 있음\\
상환 결과(사후) & 16 & \ft{total_pymnt}, \ft{recoveries}, \ft{last_pymnt_d}. \textbf{누수}\\
hardship 제도(사후) & 15 & \ft{hardship_*}. \textbf{누수}\\
식별\ensuremath{\cdot}메타 & 2 & \ft{id}, \ft{url}\\\bottomrule
\end{tabularx}
\caption{컬럼 그룹(컬럼 사전 기준). 모델 입력 검토 104개, 누수 31개, 제외 후보 2개, 확인 필요 1개, 타깃 1개, 시간 분할 기준 1개.}
\end{table}

\begin{warnbox}[누수]
사후 컬럼을 넣으면 AUC가 0.9999까지 오른다(빼면 0.72). 점수는 완벽해 보이지만 실제 심사 때는 그 값이 없어 쓸모가 없다. 코드에서 컬럼 이름 목록과 테스트로 막아 두었다.
\end{warnbox}

\section{핵심 피처}
2016년 1월 1일 기준 as-of LightGBM(시드 5개 평균)을 2016년 발행 대출에 적용해, 각 피처가 부도 예측에 기여한 크기(평균 $|\mathrm{SHAP}|$)를 쟀다.

\begin{table}[H]\centering\small
\begin{tabularx}{\linewidth}{@{}rlLrl@{}}\toprule
순위 & 피처 & 뜻 & 비중 & 방향\\\midrule
1 & \ft{int_rate} & LC가 매긴 이자율 & 15.1\% & 높을수록 위험 \ensuremath{\uparrow}\\
2 & \ft{annual_inc} & 자기 신고 연소득 & 8.0\% & 높을수록 위험 \ensuremath{\downarrow}\\
3 & \ft{dti} & 월 부채 상환액 \ensuremath{\div} 월소득 & 6.1\% & 높을수록 위험 \ensuremath{\uparrow}\\
4 & \ft{addr_state} & 거주 주 & 4.1\% & 범주형\\
5 & \ft{earliest_cr_line} & 첫 신용 계좌 개설 연도 & 3.1\% & 해석 주의(아래)\\
6 & \ft{term} & 만기(36/60개월) & 2.8\% & 비선형\\
7 & \ft{purpose} & 대출 목적 & 2.5\% & 범주형\\
8 & \ft{loan_amnt} & 신청 금액 & 2.5\% & 높을수록 위험 \ensuremath{\uparrow}\\
9 & \ft{sub_grade} & LC 세부 등급 & 2.5\% & 범주형\\
10 & \ft{mo_sin_old_rev_tl_op} & 가장 오래된 리볼빙 계좌 경과 개월 & 2.5\% & 높을수록 위험 \ensuremath{\downarrow}\\
11 & \ft{revol_util} & 리볼빙 한도 대비 사용률 & 2.2\% & 높을수록 위험 \ensuremath{\uparrow}\\
12 & \ft{mths_since_recent_bc} & 최근 bankcard 개설 후 경과 개월 & 2.1\% & 높을수록 위험 \ensuremath{\downarrow}\\
13 & \ft{total_bc_limit} & bankcard 한도 합계 & 2.1\% & 높을수록 위험 \ensuremath{\downarrow}\\
14 & \ft{percent_bc_gt_75} & 한도 75\% 초과 사용 카드 비율 & 2.0\% & 높을수록 위험 \ensuremath{\uparrow}\\
15 & \ft{inq_last_6mths} & 최근 6개월 신용조회 수 & 2.0\% & 높을수록 위험 \ensuremath{\uparrow}\\\bottomrule
\end{tabularx}
\caption{부도 예측 기여 상위 15개(전체 101개 중). 상위 10개가 49\%, 상위 20개가 69\%를 차지한다.}
\end{table}

\begin{keybox}[핵심 피처에서 읽을 것]
\begin{itemize}
\item LC 가격 정보(\ft{int_rate}, \ft{grade}, \ft{sub_grade}, \ft{term}, \ft{installment})가 전체 기여의 23\%다. 나머지 77\%는 다른 변수에 흩어져 있다.
\item \ft{fico_range_low}는 1.5\%로 작다. FICO 정보는 이미 등급과 금리에 반영돼 있기 때문이다.
\item 그룹으로 묶으면 대출 조건 31\%, 차주 기본 정보 23\%, 리볼빙\ensuremath{\cdot}bankcard 20\%, 연체\ensuremath{\cdot}조회 13\% 순이다.
\end{itemize}
\Intuit{LC 금리가 1등이지만, 모델은 금리만 보는 게 아니라 소득\ensuremath{\cdot}부채\ensuremath{\cdot}카드 사용 습관을 함께 본다.}
\end{keybox}

\begin{warnbox}[SHAP 방향 읽기]
방향은 피처 값과 SHAP의 상관으로 정한 요약이다. 인과가 아니다. 예를 들어 \texttt{earliest\_\allowbreak{}cr\_\allowbreak{}line}(연도)은 ``최근일수록 위험 \ensuremath{\downarrow}''로 나와 직관과 반대인데, 다른 신용 이력 변수와 겹치는 정보를 나눠 가지면서 생긴 결과일 수 있다. 방향을 주장할 때는 단일 변수 그림(부분 의존도)으로 따로 확인해야 한다.
\end{warnbox}

\section{등급 안에 남은 정보}
등급이 이미 설명한 부분을 빼고, \textbf{같은 세부등급 안에서} 부도를 가르는 힘을 쟀다(평가 기간을 보지 않도록 2007\textasciitilde 2012년 대출만 사용).

\begin{table}[H]\centering\small
\begin{tabularx}{\linewidth}{@{}lrrL@{}}\toprule
피처 & 전체 AUC$-0.5$ & 등급 안 AUC$-0.5$ & 해석\\\midrule
\ft{annual_inc} & $-0.054$ & $-0.080$ & 등급이 소득을 거의 안 씀. 등급 안에서 더 강함\\
\ft{lti}(대출액 \ensuremath{\div} 소득) & 0.094 & 0.057 & 상환 부담\\
\ft{acc_open_past_24mths} & 0.067 & 0.051 & 최근 신용 확장\\
\ft{pti}(상환액 \ensuremath{\div} 소득) & 0.082 & 0.049 & 상환 부담\\
\ft{dti} & 0.050 & 0.037 & 부채 수준\\
우편번호 부도율 & 0.033 & 0.034 & 등급이 지역을 안 씀\\
\ft{int_rate} & 0.155 & 0.013 & 정보 대부분이 등급과 같음\\
\ft{fico_range_low} & $-0.105$ & $-0.006$ & 정보 대부분이 등급과 같음\\
무작위 노이즈 & 0.002 & 0.004 & 기준선\\\bottomrule
\end{tabularx}
\caption{세부등급 안 AUC는 세부등급별 AUC의 대출 수 가중 평균. 0이면 정보 없음.}
\end{table}

\section{시간에 따른 변화}
\begin{itemize}
\item 금리 중앙값: 2013년 14.3\% \ensuremath{\to} 2016년 12.0\%. 그런데 실제 부도율은 12.2\% \ensuremath{\to} 15.9\%로 올랐다.
\item \ft{dti} 중앙값: 16.85 \ensuremath{\to} 18.25. 차주의 부채 부담이 커졌다.
\item \ft{revol_util}, \ft{bc_util}은 2012년 대비 분포 이동이 가장 크다(2018년 PSI 0.34).
\end{itemize}
\Intuit{위험은 커지는데 가격은 내려갔다. 과거 데이터로 맞춘 확률이 실제보다 낮게 나온 이유다.}

\section{결측 구조}
\begin{table}[H]\centering\small
\begin{tabularx}{\linewidth}{@{}lL@{}}\toprule
패턴 & LC에서 본 것\\\midrule
구조적 결측 & \ft{sec_app_*}는 개인 신청에서 100\% 결측. 단, 공동신청인데도 6\%가 비어 있음\\
사건 없음 & \ft{mths_since_last_delinq}(52\%)는 연체 이력이 없어서 빔. 결측 자체가 정보\\
MNAR 의심 & \ft{emp_length}가 비면 부도율 26.3\%, 있으면 19.0\%\\
시기(MAR) & 13개 컬럼이 함께 29.8\% 결측, 21개 컬럼이 함께 2.4\% 결측. 발행 시기로 설명됨\\
누수 & \ft{hardship_*}가 비어 있으면 부도율이 37\%p 낮음(어려움을 겪은 대출에만 값이 생김)\\\bottomrule
\end{tabularx}
\caption{결측 패턴(자세한 표는 \ft{결측_검증시트.xlsx}).}
\end{table}

\section{Feature engineering 결과}
상환 부담(\ft{pti}, \ft{lti}), 신용 이력 개월 수, 근속 연수 숫자화, 우편번호 부도율을 추가해 as-of 학습으로 다시 평가했다.

\begin{warnbox}[결과: 모두 기각]
모든 후보의 선별 가치 변화가 $-0.01$\textasciitilde $+0.02$\%p였다. 같은 모델을 시드만 바꿔 돌린 차이(최대 0.09\%p)보다 작았다. EDA에서 본 등급 안 정보는 이미 원본 변수(소득, 상환액, 대출액, 주)에 들어 있었다.
\Intuit{피처 하나의 효과보다 모델을 다시 돌렸을 때의 우연한 차이가 더 컸다. 한 번 돌려 본 비교로 피처를 고르면 안 된다.}
\end{warnbox}

% ======================================================================
\chapter{데이터와 검증 기법}

\section{Data leakage}
\begin{keybox}[정의]
Leakage는 예측 시점에 알 수 없는 정보가 모델 입력에 섞이는 것이다. 시험 전에 정답지를 보는 것과 같다.
\Intuit{점수가 너무 좋으면 먼저 누수를 의심한다.}
\end{keybox}

\section{종결 표본 편향과 censoring}
Current 대출은 결과를 모른다(\textbf{censoring}). 종결 대출만 쓰면, 최근 발행분일수록 일찍 끝난 대출(부도 포함)만 남아 부도율이 부풀려진다. 결과를 모르는 대출까지 쓰는 표준 방법이 \textbf{competing risks}(부도 vs 상환) 누적 발생률이다.

\begin{formulabox}[Aalen--Johansen 누적 부도율]
\[ S(t)=\prod_{u\le t}\Bigl(1-\frac{d_1(u)+d_2(u)}{n(u)}\Bigr),\qquad
\mathrm{CIF}_1(t)=\sum_{u\le t} S(u-1)\,\frac{d_1(u)}{n(u)} \]
$n(u)$: $u$개월째 위험 집합, $d_1$: 부도, $d_2$: 상환.
\end{formulabox}

\section{As-of 학습과 시간 분할}
연도 Y의 대출을 고를 때는 Y년 1월 1일 이전에 발행되고 그 전에 결과가 확정된 대출로만 학습한다. 부도는 마지막 납부 후 약 5개월 뒤 상각되므로 그만큼 늦게 확정된다고 본다.

\begin{table}[H]\centering\small
\begin{tabular}{@{}lrr@{}}\toprule
분할 방식 & AUC & 해석\\\midrule
무작위 분할 & 약 0.73 & 미래 정보가 섞여 낙관적\\
as-of 시간 분할 & 0.62\textasciitilde 0.68 & 실제 투자 상황에 가까움\\\bottomrule
\end{tabular}
\caption{같은 LightGBM의 예측력.}
\end{table}

\begin{warnbox}[가정이 결과를 바꾼다]
상각 지연을 5개월로 두면 2016년 초 학습 표본의 부도율은 17.4\%로 최종 18.1\%보다 낮다. 지연을 0으로 두면 22.2\%로 높다. 편향의 방향이 가정에 달려 있다.
\end{warnbox}

\section{결측 메커니즘}
\begin{table}[H]\centering\small
\begin{tabularx}{\linewidth}{@{}lL@{}}\toprule
유형 & 뜻\\\midrule
MCAR & Missing completely at random. 완전히 우연한 결측\\
MAR & Missing at random. 관측된 다른 변수로 설명되는 결측(예: 발행 시기)\\
MNAR & Missing not at random. 빠진 값 자체 때문에 빔. 데이터만으로 증명 불가\\
구조적 & 해당 사항이 없어서 빔(예: 개인 신청의 공동신청 컬럼)\\\bottomrule
\end{tabularx}
\caption{결측 유형.}
\end{table}

\begin{insight}[Target encoding]
우편번호처럼 범주가 많은 변수는 범주별 부도율로 바꿔 쓸 수 있다. 이때 자기 행의 정답을 쓰지 않도록 \textbf{out-of-fold}로 계산하고, 건수가 적은 범주는 전체 평균 쪽으로 \textbf{smoothing}한다.
\end{insight}

% ======================================================================
\chapter{수익률과 비교 기준}

\section{네 가지 수익률 정의}
같은 대출도 정의에 따라 수익률이 크게 달라진다. 2013년 발행분 기준으로 비교한다.

\begin{formulabox}[정의]
\begin{align*}
\text{단순 연환산:}\quad & r=\Bigl(\tfrac{\text{총 상환액}}{\text{대출액}}\Bigr)^{12/T}-1\\
\text{IRR:}\quad & \sum_{k=0}^{H}\frac{CF_k}{(1+m)^k}=0,\qquad \mathrm{IRR}=(1+m)^{12}-1\\
\text{재투자 수익률:}\quad & FV=\sum_k CF_k(1+g)^{T-k},\qquad R=(FV/P)^{12/T}-1
\end{align*}
\end{formulabox}

\begin{table}[H]\centering\small
\begin{tabularx}{\linewidth}{@{}lrL@{}}\toprule
정의 & 2013년 & 가정\\\midrule
수업식(부도면 0) & 11.9\% & 부도의 원금 손실을 무시\\
단순 연환산(우승 논문) & 2.8\% & 매달 돌려받은 원리금이 놀고 있다고 봄\\
국채 재투자 & 3.3\% & 돌려받은 돈을 국채에 넣음\\
통합 IRR(수수료 1\% 후) & 8.9\% & 돌려받은 돈을 같은 수익률로 재투자\\
국채(3년) & 0.8\% & 비교 기준\\\bottomrule
\end{tabularx}
\caption{같은 대출, 네 가지 숫자.}
\end{table}

\begin{workout}[왜 단순 연환산은 낮게 나오나]
\Step{1} 금리 12\%, 36개월, 약정대로 다 갚은 대출의 월 납입금은 원금의 약 3.32\%다. \Concept{원리금 균등 상환}
\Step{2} 총 상환액 \ensuremath{\div} 원금 $\approx 1.196$, 단순 연환산은 $1.196^{1/3}-1\approx 6.1\%$.
\Step{3} 하지만 돈은 매달 돌아오므로 실제 수익률(IRR)은 $1.01^{12}-1=12.7\%$다. \Concept{IRR}
\Ans{단순 연환산은 돌려받은 돈을 다시 투자하지 않는다고 보기 때문에 약정 금리의 절반 정도로 나온다.}
\end{workout}

\section{Cohort(발행월 묶음)}
같은 달에 발행된 대출을 하나의 포트폴리오로 보고 IRR을 계산한다. 발행월 사이의 차이가 곧 ``언제 투자했느냐''의 위험이다.

\section{무위험수익률(rf)의 선택}
과제가 정한 기준은 국채다. 어떤 국채를 쓸지는 우리가 고른다.
\begin{table}[H]\centering\small
\begin{tabularx}{\linewidth}{@{}lrrL@{}}\toprule
기준 & 평균 & 전부 승인 초과수익 & 메모\\\midrule
3개월 국채 & 0.15\% & +5.62\%p & 교과서적 rf\\
만기 맞춘 국채(기본) & 0.94\% & +4.86\%p & 과제 정의에 충실\\
국채 +3\%p & 3.94\% & +1.86\%p & 유동성 보상 가정\\
회사채 BAA & 4.91\% & +0.81\%p & 17개 발행월에서 회사채보다 낮음\\\bottomrule
\end{tabularx}
\caption{비교 기준을 바꾸면 결론의 강도가 바뀐다.}
\end{table}
\begin{warnbox}[회사채는 rf가 아니다]
회사채를 빼서 계산한 비율은 Sharpe가 아니라 \textbf{information ratio}다. BAA는 장기채이고 약속 수익률이라 LC에 불리한 비교다.
\end{warnbox}

% ======================================================================
\chapter{모델과 보정}

\section{Logistic regression과 LightGBM}
\begin{itemize}
\item \textbf{Logistic regression}: 해석이 쉬운 기준 모델. 학습 데이터가 적은 해(2013년)에는 LightGBM보다 나았다(AUC 0.646 vs 0.623).
\item \textbf{LightGBM}: gradient boosting tree. 파라미터는 고정하고(튜닝 효과 약 0.001), 시드 5개의 예측을 평균한다.
\item 실전 조건에서는 금리만 쓴 모델도 비슷한 예측력을 보였다. 복잡한 모델의 가치는 수익보다 \textbf{손실 꼬리를 줄이는 데} 있었다(승인 부도율 14.5\% \ensuremath{\to} 12.4\%).
\end{itemize}

\section{기대수익과 LGD}
\begin{formulabox}[연 기대수익(근사)]
\[ E[r]=(1-\hat p)\,r_{\text{약정}}-\hat p\cdot\mathrm{LGD} \]
$\hat p$: 예측 부도확률, LGD: 부도 시 손실률(그 시점 확정 부도 대출 평균).
\end{formulabox}

\section{확률 보정(calibration)}
as-of 학습 확률은 수준이 틀렸다(2013년 예측 6.1\% vs 실제 12.2\%). 세 방법을 시험했다.
\begin{itemize}
\item \textbf{Isotonic regression}: 예측값과 실제 결과를 단조 함수로 맞춤
\item \textbf{만기 도달 보정}: 그 시점에 이미 만기까지 지켜본 대출로만 맞춤
\item \textbf{Logit shift}: $\operatorname{logit}\hat p'=\operatorname{logit}\hat p+c$, 평균이 생존분석 추정치가 되게 $c$를 고름
\end{itemize}
\begin{warnbox}[결과: 모두 실패]
원인은 표본 편향보다 \textbf{신용 악화}였다. 과거에 맞추는 보정은 추세를 따라가지 못한다. 대신 PD 과소추정과 LGD 과대추정(시점 52\textasciitilde 62\% vs 실제 36.5\%)이 서로 상쇄돼, 보정하지 않은 조합이 투자 성과가 가장 좋았다.
\end{warnbox}

% ======================================================================
\chapter{평가 지표}

\section{Sharpe ratio와 그 함정}
\begin{formulabox}[과제식 Sharpe와 항등식]
\[ \mathrm{SR}=\frac{\bar x}{s_x},\quad x_i=a_i e_i \qquad\Longrightarrow\qquad
\mathrm{SR}=\frac{\sqrt a\, m}{\sqrt{s^2+(1-a)m^2}} \]
$a$: 승인 비율, $m,s$: 승인 대출 초과수익의 평균\ensuremath{\cdot}표준편차, 거절 자금은 초과수익 0.
\end{formulabox}
\begin{whybox}[역U자의 정체]
승인 비율이 줄면 $\sqrt a$ 때문에 SR이 기계적으로 변한다. 무작위로 25\%만 승인해도 SR이 $-0.122$에서 $-0.061$로 약 $\sqrt{0.25}$배가 됐다. 선별 실력과 투자 비중 효과가 섞여 있다.
\end{whybox}

\section{Information ratio}
\[ \mathrm{IR}=\frac{\overline{R-B}}{s_{R-B}} \]
비교 지수 $B$(예: 회사채) 대비 초과수익의 안정성이다. 회사채 기준으로 보면 선별 규칙의 IR은 1.90, 전부 승인은 1.08이었다.

\section{RAROC과 EVA}
\begin{formulabox}[은행 관점]
\[ \mathrm{RAROC}=\frac{r-\mathrm{EL}-f-o}{K},\qquad \mathrm{EVA}=r-\mathrm{EL}-f-o-hK \]
\[ K=\mathrm{LGD}\Bigl[N\Bigl(\tfrac{G(\mathrm{PD})+\sqrt{R}\,G(0.999)}{\sqrt{1-R}}\Bigr)-\mathrm{PD}\Bigr],\quad R=0.03w+0.16(1-w) \]
EL: 예상손실, $f$: 조달비용, $o$: 운영비, $h$: 요구수익률, $K$: Basel 기타 소매 자본, $w=\frac{1-e^{-35\mathrm{PD}}}{1-e^{-35}}$.
\end{formulabox}
\begin{warnbox}[주 지표로 쓰지 않은 이유]
기본 가정에서 RAROC 중앙값이 94\%라 대출의 99\%를 승인했다. 결과가 데이터보다 비용\ensuremath{\cdot}자본 가정에 좌우된다.
\end{warnbox}

\section{꼬리 위험과 효용}
\begin{itemize}
\item \textbf{CVaR(하위 5\% 평균)}: 최악 구간의 평균 손실. 대출 수익은 모든 등급에서 음의 skewness다.
\item \textbf{Sortino ratio}: 하방 변동성만 위험으로 본다. 이번 데이터에서는 Sharpe와 순위가 거의 같았다.
\item \textbf{확실성 등가}: $\mathrm{CE}=\bar X-\tfrac{\gamma}{2}\mathrm{Var}(X)$. 위험회피 $\gamma=50$에서도 최적은 전부 승인이었다. 분산된 포트폴리오에서는 규모가 위험보다 중요하다.
\end{itemize}

\section{선별 곡선}
\begin{formulabox}[선별 가치]
\[ \mathrm{SV}=\frac{1}{|Q|}\sum_{q\in Q}\bigl[\mathrm{ER}_{\text{model}}(q)-\mathrm{ER}_{\text{random}}(q)\bigr],\quad Q=\{0.1,\dots,0.9\} \]
$\mathrm{ER}(q)$: 점수 상위 $q$를 승인했을 때 승인 대출의 국채 대비 초과 IRR.
\end{formulabox}
기대수익 순위의 SV는 1.03\%p(95\% 구간 0.95\textasciitilde 1.12), 부도확률 순위는 0.22\%p였다. 임계값을 고르는 자의성이 없다. 같은 생각은 profit scoring(Serrano-Cinca \& Gutiérrez-Nieto, 2016), EMP(Verbraken 외, 2014)에 이미 있다.

\section{Brinson decomposition}
\begin{formulabox}[배분 \ensuremath{\cdot} 선별 \ensuremath{\cdot} 상호작용]
\begin{align*}
\text{배분}&=\sum_g (w^s_g-w^b_g)(R^b_g-R^b), \qquad
\text{선별}=\sum_g w^b_g(R^s_g-R^b_g),\\
\text{상호작용}&=\sum_g (w^s_g-w^b_g)(R^s_g-R^b_g)
\end{align*}
$g$: 등급(또는 세부등급), $w$: 금액 비중, $R$: pooled IRR, $s$: 전략, $b$: 전부 승인.
\end{formulabox}
\begin{keybox}[결과]
세부등급 35개 기준 등급 안 선별 효과가 2014\textasciitilde 2016년 0.69\textasciitilde 1.11\%p, 2017년 1.64\%p였고, 배분 효과는 0\textasciitilde 0.23\%p였다. 2013년만 배분이 더 컸다.
\Intuit{펀드매니저로 치면 ``좋은 업종을 많이 산'' 사람이 아니라 ``같은 업종 안에서 좋은 종목을 고른'' 사람이다.}
\end{keybox}

\section{예측 지표}
\begin{itemize}
\item \textbf{AUC}: 부도와 정상을 순위로 얼마나 잘 가르는지. 0.5는 동전 던지기.
\item \textbf{Brier score}: $\frac1n\sum(\hat p_i-y_i)^2$. 확률 수준까지 본다.
\item \textbf{ECE}: 예측 구간별 (평균 예측 $-$ 실제 비율)의 가중 평균. 보정이 맞는지 본다.
\item \textbf{PSI}: $\sum(b-a)\ln(b/a)$. 분포가 시간에 따라 얼마나 이동했는지. 0.1 이상이면 의미 있는 변화.
\end{itemize}

% ======================================================================
\chapter{검증 방법론}

\begin{table}[H]\centering\small
\begin{tabularx}{\linewidth}{@{}lL@{}}\toprule
기법 & 하는 일\\\midrule
Placebo(무작위 승인) & 선별 실력이 없으면 성과가 0이 나와야 함. 평가 방법이 가짜 성과를 만들지 않는지 확인\\
Placebo(라벨 섞기) & 정답을 섞어 학습한 모델은 선별 가치가 0 근처여야 함(실제 0.01\%p)\\
시드 반복 & 같은 모델을 시드 5개로 돌려 결론이 우연인지 확인\\
Block bootstrap & 발행월 단위로 다시 뽑아 95\% 신뢰구간 계산(같은 달 대출은 함께 움직임)\\
보류 구간 & 한 번도 안 쓴 구간(2017년)에서 같은 방향이 나오는지 확인\\
사전 등록 & test를 받기 전에 가설\ensuremath{\cdot}규칙\ensuremath{\cdot}성공 기준을 고정(태그 \ft{test-protocol-v1})\\
주장 장부 & 주장마다 근거와 상태(지지/잠정/기각)를 기록(\ft{docs/research_log.md})\\\bottomrule
\end{tabularx}
\caption{결론을 지키는 장치들.}
\end{table}

\begin{warnbox}[데이터 엿보기]
같은 평가 구간을 여러 번 쓰면, 그 구간에 맞춘 결론을 고르게 된다. 지금까지 비교한 규칙\ensuremath{\cdot}지표 조합은 60개 이상이다. 그래서 최종 판단은 아직 안 본 공식 test로만 한다.
\end{warnbox}

\begin{qa}{Q. 결론이 틀렸던 적은?}
있다. ``부도확률 순위는 무작위와 비슷하다''는 시드 하나로 본 착각이었다. 시드 5개와 부트스트랩으로 보니 유의하게 양수였고, 해마다 커졌다(2017년 2.00\%p). 그래서 기각하고 기록했다.
\end{qa}

% ======================================================================
\chapter{앞으로 쓸 기법}

\section{test 평가(고정됨)}
\begin{table}[H]\centering\small
\begin{tabularx}{\linewidth}{@{}lL@{}}\toprule
가설 & 지지 기준\\\midrule
H1 성과 귀속 & 2013\textasciitilde 16년 통합 등급 안 선별 $>0$이고 배분보다 크며, 4개 연도 중 3개 이상 양수\\
H2 선별 가치 & 기대수익 순위 SV의 95\% 구간 하한 $>0$\\
H3 국채 대비 & 전부 승인이 발행월의 90\% 이상에서 국채보다 높음\\
H4 회사채 대비 & 선별 규칙의 BAA 이하 발행월 수 $<$ 전부 승인\\\bottomrule
\end{tabularx}
\caption{주 승인 규칙은 ``기대수익 $>$ 만기 국채''(R1).}
\end{table}

\section{남은 후보}
\begin{itemize}
\item \textbf{거시 변수}: 실업률 등으로 신용 악화를 설명할 수 있는지. 독립 관측이 48개월뿐이라 약할 수 있다.
\item \textbf{60개월 대출 확장}: 종결 편향을 처리해야 한다.
\item \textbf{조기상환 모형}: 정상 상환의 76\%가 조기상환이다. 부도 다음으로 큰 수익 변수다.
\item \textbf{Partial dependence}: 핵심 피처의 방향을 단일 변수 그림으로 확인한다.
\end{itemize}

\begin{insight}[이 문서의 숫자]
모든 숫자의 근거 파일은 \ft{docs/decisions.md}와 \ft{results/} 폴더에 있다. 피처 중요도는 \ft{scripts/17_feature_importance.py}, 결측 표는 \ft{docs/결측_검증시트.xlsx}로 다시 만들 수 있다.
\end{insight}


\chapter{시작 전에: 지도와 실행법}

\section{이 문서를 읽는 법}
코드를 처음부터 끝까지 외울 필요는 없다. 이 프로젝트의 코드는 결국 \textbf{여섯 단계를 한 줄로 이은 것}이다. 각 장은 한 단계를 맡고, 같은 순서로 쓴다.
\begin{enumerate}[leftmargin=*]
\item \textbf{무엇을 하려는가}(한 문단).
\item \textbf{실제 코드}(줄 번호 포함).
\item \textbf{줄별 설명}(``3번 줄'' 식으로 짚는다).
\item \textbf{따라 해 보기}: 데이터 없이 도는 작은 예제와 그 출력.
\item \textbf{흔한 실수}.
\end{enumerate}

\begin{keybox}[전체 흐름 한 장]
\small\begin{tabularx}{\linewidth}{@{}p{5mm}Lp{40mm}@{}}\toprule
단계 & 하는 일 & 파일\\\midrule
1 & CSV를 읽어 상태를 ``상환/부도''로 정리, 미래 정보 제거 & \ft{src/lc/data.py}\\
2 & 그 시점에 결과를 알 수 있던 대출만 학습에 쓰기 & \ft{src/lc/asof.py}\\
3 & LightGBM으로 부도확률 $p$ 만들기 & scripts 23, 30번\\
4 & $p$를 기대수익 점수로 바꾸고 발행월마다 상위 $q$ 승인 & scripts 30, 38, 39번\\
5 & 승인한 묶음의 실제 현금흐름으로 IRR 채점 & \ft{src/lc/returns.py}\\
6 & 매달 초과수익 \ensuremath{\to} Sharpe, 부트스트랩, 1-SE 선택, 파레토 & scripts 39, 41번\\\bottomrule
\end{tabularx}
\Intuit{예상(점수)으로 고르고, 실제 결과(IRR)로 채점하고, 그 채점표를 48개월 모아 Sharpe를 낸다.}
\end{keybox}

\section{폴더 구조}
\begin{table}[H]\centering\small
\begin{tabularx}{\linewidth}{@{}p{38mm}L@{}}\toprule
위치 & 들어 있는 것\\\midrule
\ft{src/lc/} & 여러 스크립트가 같이 쓰는 \textbf{함수 모음}(라이브러리). 데이터 정리, 시점 학습, 수익률, 비교 기준\\
\ft{scripts/NN_이름.py} & 실험 하나 = 파일 하나. 번호가 실험 순서다. 실행하면 \ft{results/}에 표와 로그를 남긴다\\
\ft{results/} & 실행 결과(\ft{.csv}, \ft{.log}). 문서의 숫자는 모두 여기서 나온다\\
\ft{docs/} & 결정 기록(\ft{decisions.md}), 연구 노트(\ft{research_log.md}), test 규칙(\ft{test_protocol*.md}), 보고서\\
\ft{data/} & 원본 CSV, 변환한 parquet, 국채·\allowbreak{}회사채 금리(\ft{external/})\\\bottomrule
\end{tabularx}
\end{table}

\section{실행하는 법}
터미널(맥은 ``터미널'' 앱)을 열고 저장소 폴더로 이동한 뒤 실행한다. 이 컴퓨터는 conda 환경 \ft{cmapss}의 파이썬을 쓴다.
\begin{lstlisting}[language=bash,numbers=none]
cd ~/Desktop/learningClub
conda activate cmapss                       # once per terminal window
python scripts/39_final_choice.py           # one experiment (a few minutes)
python docs/notes/code/snippets.py irr      # one practice example (instant)
\end{lstlisting}
\begin{warnbox}[python을 못 찾는다고 나오면]
\ft{command not found: python}은 conda 환경을 켜지 않았다는 뜻이다. \ft{conda activate cmapss}를 먼저 하거나, 전체 경로 \ft{/opt/homebrew/Caskroom/miniconda/base/envs/cmapss/bin/python}으로 실행한다. 맥 기본 \ft{python3}에는 LightGBM이 없다.
\end{warnbox}

\section{스크립트 맨 위의 두 가지 특이한 줄}
거의 모든 스크립트가 이렇게 시작한다.
\begin{lstlisting}
import _path  # noqa: F401
from pathlib import Path

SRC = (Path(__file__).parent / '30_portfolio_sharpe.py').read_text()
exec(SRC[:SRC.index("sharpe = lambda")])
\end{lstlisting}
\begin{itemize}[leftmargin=*]
\item \textbf{1번 줄} \ft{import _path}: \ft{scripts/_path.py}는 ``\ft{src/} 폴더도 찾아봐라''라고 파이썬에 알려 주는 세 줄짜리 파일이다. 그래서 다음 줄부터 \ft{from lc import returns}처럼 라이브러리를 불러올 수 있다. \ft{# noqa}는 ``안 쓰는 import라고 경고하지 마''라는 표시일 뿐이다.
\item \textbf{4번 줄}: 30번 스크립트의 \textbf{내용을 글자로} 읽는다. \ft{Path(__file__).parent}는 ``지금 이 파일이 있는 폴더''다.
\item \textbf{5번 줄}: 그 글자 중 \ft{"sharpe = lambda"}가 나오기 \textbf{직전까지만} 잘라(\ft{SRC[:위치]}) 실행한다(\ft{exec}). 30번의 앞부분은 데이터를 읽고 모델을 학습해 \ft{ev} 표를 만드는 부분이다. 결국 ``30번이 만든 \ft{ev}를 그대로 이어받아 쓴다''는 뜻이다.
\end{itemize}
\begin{warnbox}[왜 이렇게 하나, 무엇을 조심하나]
같은 학습 코드를 복사하지 않으려는 지름길이다. 대신 \textbf{앞 스크립트의 변수 이름이 갑자기 나타난다}(\ft{ev}, \ft{months}, \ft{TB} 등). 처음 보는 변수가 있으면 \ft{exec}한 파일에서 찾는다. 39번은 30번을, 41번은 39번을, 40번은 35번을, 35번은 15번을 이어받는다. 또 앞 파일에서 자르는 기준 문자열이 바뀌면 \ft{ValueError: substring not found}가 난다.
\end{warnbox}

\chapter{파이썬, numpy, pandas: 이 코드에 나오는 것만}

이 장은 문법 교과서가 아니다. 프로젝트 코드에 \textbf{실제로 나오는 것만} 모았다. 모르는 기호가 나오면 이 장으로 돌아온다.

\section{파이썬 기본 여섯 가지}
\begin{lstlisting}
QS = [0.1, 0.2, 0.3]                  # list: ordered values, QS[0] is 0.1
hurdles = {'tb': 0.015, 'baa': 0.05}  # dict: name -> value, hurdles['tb']
for q in QS:                          # loop: q takes each value in turn
    print(f'approve {q:.0%}')         # f-string: {q:.0%} prints 10%
sharpe = lambda x: x.mean() / x.std(ddof=1)   # one-line function
def cut(score, q):                    # normal function
    return score >= 0.5
\end{lstlisting}
\begin{itemize}[leftmargin=*]
\item \textbf{1번 줄} 리스트: 대괄호, 순서가 있다. 번호는 0부터 센다.
\item \textbf{2번 줄} 딕셔너리: 중괄호, ``이름: 값'' 짝이다. \ft{hurdles.items()}는 (이름, 값) 짝을 하나씩 준다.
\item \textbf{3--4번 줄} \ft{for}: 들여쓰기(공백 4칸)한 줄들이 반복된다. 파이썬은 \textbf{들여쓰기가 문법}이라, 칸이 틀리면 오류가 난다.
\item \textbf{4번 줄} f-문자열: 따옴표 앞 \ft{f}, 중괄호 안에 변수. \ft{:.2f}는 소수 둘째 자리, \texttt{:.0\%}는 백분율이다.
\item \textbf{5번 줄} \ft{lambda}: 이름 없는 한 줄 함수. ``\ft{x}를 받아 평균 \ensuremath{\div} 표준편차를 돌려준다''.
\item \textbf{6--7번 줄} \ft{def}: 여러 줄 함수. \ft{return} 뒤가 결과다.
\end{itemize}
그 밖에 자주 나오는 것:
\begin{itemize}[leftmargin=*]
\item \ft{[f(x) for x in xs]}: ``xs의 각 x에 f를 적용한 리스트''(리스트 컴프리헨션).
\item \ft{a, b = 1, 2}: 한 번에 두 변수.
\item \ft{x if cond else y}: 조건에 따라 x 또는 y.
\item \ft{**dict}: 딕셔너리를 펼쳐 넣기. \texttt{\{'a': 1, **row\}}는 \ft{row}의 내용을 합친 새 딕셔너리다.
\end{itemize}

\section{numpy: 숫자 배열을 한 번에}
\ft{np.array}는 숫자를 줄지어 담은 배열이다. 핵심은 \textbf{반복문 없이 전체에 한 번에 계산한다}는 것이다.
\begin{lstlisting}
import numpy as np
r = np.array([0.07, 0.11, 0.15])
bad = np.array([0, 1, 0])
r * 100                      # [7, 11, 15]  every element at once
np.where(bad == 1, 0.0, r)   # [0.07, 0, 0.15]  if default then 0 else rate
np.quantile(r, 0.7)          # 70th percentile
r >= np.quantile(r, 0.7)     # [False, False, True]  boolean array
r[r > 0.1]                   # [0.11, 0.15]  keep where True
rng = np.random.default_rng(11)   # random generator with a fixed seed
rng.random(3)                # 3 numbers in [0, 1), same every run
\end{lstlisting}
\begin{itemize}[leftmargin=*]
\item \textbf{4번 줄}: 배열 \ensuremath{\times} 숫자는 원소마다 곱한다.
\item \textbf{5번 줄} \ft{np.where(조건, 참일 때, 거짓일 때)}: 엑셀의 IF를 배열 전체에 한 번에 한다.
\item \textbf{6--7번 줄} 분위수: \ft{np.quantile(s, 1 - q)} 이상이면 ``상위 $q$''다. 상위 30\%를 고르려면 70\% 분위수 이상을 고른다. 이 프로젝트의 모든 승인 규칙이 이 한 줄이다.
\item \textbf{8번 줄}: 참/거짓 배열을 대괄호에 넣으면 참인 자리만 남는다.
\item \textbf{9--10번 줄}: seed(씨앗)를 고정하면 ``무작위''도 매번 같은 값이 나와 결과를 재현할 수 있다. 코드의 \ft{default_rng(2026)}, \ft{default_rng(500 + k)}가 모두 이것이다.
\end{itemize}
\paragraph{자주 나오는 함수.} \ft{x.mean()} 평균, \ft{x.std(ddof=1)} 표본 표준편차(\ft{ddof=1}은 $n-1$로 나눈다는 뜻), \ft{x.sum()}, \ft{np.zeros(n, bool)} 거짓 $n$개, \ft{np.ones}, \ft{np.arange(5)}는 0--4, \ft{np.isclose(a, b)}는 소수 비교(0.8과 0.80000001을 같다고 본다).
\paragraph{2차원 배열과 \ft{axis}.} \ft{x[BIX]}처럼 번호 표로 뽑으면 표 모양 배열이 된다. \ft{.mean(axis=1)}은 \textbf{행마다} 평균, \ft{axis=0}은 \textbf{열마다} 평균이다. 10장의 부트스트랩에서 쓴다.

\section{pandas: 엑셀 표를 코드로}
\ft{pd.DataFrame}은 행과 열이 있는 표, 열 하나는 \ft{Series}다.
\begin{lstlisting}
import pandas as pd
df['r']                          # one column (Series)
df[['id', 'r']]                  # several columns (DataFrame)
df['bad'] = (df['status'] == 'Charged Off').astype(int)   # new column
df[df['grade'] == 'A']           # rows where condition is True
df.loc[rows, 'r']                # by label     df.iloc[3]  by position
df['r'].to_numpy()               # Series -> numpy array
df.groupby('ym')['amt'].sum()    # sum per issue month
df.groupby('ym').indices         # {month: array of row positions}
d.merge(tb, on=['ym', 'term_m'], how='left')   # VLOOKUP
\end{lstlisting}
\begin{itemize}[leftmargin=*]
\item \textbf{4번 줄}: 비교(\ft{==})는 참/거짓 열을 만들고, \ft{.astype(int)}가 1/0으로 바꾼다. 부도 여부 \ft{bad}가 이렇게 만들어진다.
\item \textbf{5번 줄}: 조건을 여러 개 묶을 때는 \texttt{\&}(그리고), \ft{|}(또는), \texttt{\textasciitilde}(아님)을 쓰고 \textbf{각 조건을 괄호로 감싼다}. \ft{and}/\ft{or}를 쓰면 오류가 난다.
\item \textbf{6번 줄} \ft{loc}은 이름표(인덱스) 기준, \ft{iloc}은 0부터 센 위치 기준이다. 코드가 \ft{reset_index(drop=True)}를 자주 하는 이유는 이름표를 0, 1, 2, \ldots{}로 다시 매겨 둘을 같게 만들기 위해서다.
\item \textbf{7번 줄}: 계산을 빠르게 하려고 열을 numpy 배열로 바꾼다. 예를 들어 \ft{P_amt}는 \ft{funded_amnt} 열을 배열로 바꾼 것이다.
\item \textbf{9번 줄} \ft{groupby().indices}: \textbf{이 프로젝트에서 가장 중요한 pandas 표현}이다. ``발행월마다 그 달 대출의 행 번호''를 준다. 발행월마다 자르기(8장), 발행월마다 IRR(9장)이 모두 이것을 돈다.
\item \textbf{10번 줄} \ft{merge}: 엑셀의 VLOOKUP. 발행월과 만기가 같은 국채 금리를 대출 옆에 붙인다(6장).
\end{itemize}

\begin{workout}[따라 해 보기: \ft{snippets.py groupby}]
6건짜리 표로 달마다 합계, 등급별 부도율, 달마다 행 번호를 본다. 출력의 마지막 두 줄:
\begin{lstlisting}[style=out]
2013-01 -> rows [0 1 2]
2013-02 -> rows [3 4 5]
\end{lstlisting}
\Ans{1월 대출은 0--2번 행, 2월 대출은 3--5번 행이다. 이후 모든 ``발행월마다''는 이 행 번호 묶음을 하나씩 도는 것이다.}
\end{workout}
\Trap{\ft{df[df['r'] > 0.1 and df['grade'] == 'A']}. pandas에서는 \texttt{(df['r'] > 0.1) \& (df['grade'] == 'A')}로 써야 한다.}

\chapter{데이터 준비와 시점 학습}

\section{상태 정리: \texttt{data.py}}
\textbf{하려는 것.} 원본의 \ft{loan_status}는 ``Fully Paid'', ``Charged Off'', ``Late (31-120 days)'' 등 문자열이다. 이것을 상환(paid), 부도(default), 연체(late), 진행 중(current) 네 가지로 정리하고, 숫자로 쓸 열(금리, 만기, 발행월)을 만든다.
\begin{lstlisting}
def classify(status):
    s = status.astype(str)
    return pd.Series(np.select(
        [s.str.contains('Fully Paid'), s.str.contains('Charged Off') | s.eq('Default'),
         s.str.startswith('Late') | s.eq('In Grace Period')],
        ['paid', 'default', 'late'], 'current'), index=status.index)

def parse(df):
    d = df.copy()
    d['ym'] = pd.to_datetime(d['issue_d'], format='%b-%Y')
    d['issue_year'] = d['ym'].dt.year
    d['r'] = pd.to_numeric(d['int_rate'].astype('string').str.strip().str.rstrip('%')) / 100
    d['term_m'] = d['term'].astype('string').str.extract(r'(\d+)')[0].astype(int)
    lp = pd.to_datetime(d['last_pymnt_d'], format='%b-%Y', errors='coerce')
    d['tau'] = months_between(d['ym'], lp).fillna(1).clip(lower=1).astype(int)
    d['outcome'] = classify(d['loan_status'])
    d['resolved'] = d['outcome'].isin(['paid', 'default'])
    d['bad'] = (d['outcome'] == 'default').astype('int8')
    return d
\end{lstlisting}
\begin{itemize}[leftmargin=*]
\item \textbf{3--6번 줄} \ft{np.select([조건1, 조건2, 조건3], [값1, 값2, 값3], 기본값)}: 엑셀의 중첩 IF. 위에서부터 처음 참인 조건의 값을 쓰고, 아무것도 아니면 \ft{'current'}다. \ft{.str.contains}는 ``그 글자를 포함하나'', \ft{.str.startswith}는 ``그 글자로 시작하나''다.
\item \textbf{9번 줄} \ft{df.copy()}: 원본을 건드리지 않으려고 복사본에서 작업한다.
\item \textbf{10번 줄}: ``Dec-2015'' 같은 글자를 날짜로 바꾼다. \texttt{\%b}는 월 약자, \texttt{\%Y}는 네 자리 연도다. 이 \ft{ym}(발행월)이 이후 모든 ``발행월마다''의 기준이다.
\item \textbf{12번 줄}: ``13.56\%'' \ensuremath{\to} 공백 제거 \ensuremath{\to} \% 제거 \ensuremath{\to} 숫자 13.56 \ensuremath{\to} 100으로 나눠 0.1356.
\item \textbf{13번 줄}: `` 36 months''에서 숫자만 뽑는다. \texttt{r'(\textbackslash d+)'}는 ``연속된 숫자''라는 정규식이다.
\item \textbf{14--15번 줄} \ft{tau}: 발행부터 마지막 납부까지 개월 수. 현금흐름 복원(7장)과 시점 학습(아래)에 쓴다. \ft{errors='coerce'}는 ``변환이 안 되면 빈 값으로''이고, \ft{fillna(1)}은 빈 값을 1로 채운다.
\item \textbf{17--18번 줄}: 결과가 확정된(상환 또는 부도) 대출만 \ft{resolved}=참. 부도면 \ft{bad}=1이 모델의 정답이다.
\end{itemize}

\paragraph{누수 목록 \ft{LEAK}.} \ft{total_pymnt}(받은 총액), \ft{recoveries}(회수금), \ft{last_pymnt_d}처럼 \textbf{승인한 뒤에야 알 수 있는} 열들이다. \ft{feature_columns()}가 이것을 모델 입력에서 뺀다. 이 열들은 채점(IRR)에는 쓰지만 고르기(모델)에는 절대 쓰지 않는다.
\begin{warnbox}[누수]
\ft{total_pymnt}를 모델에 넣으면 AUC가 0.99 가까이 나온다. 이미 결과를 보고 맞히는 것이라 실전에서는 쓸 수 없다. 새 열을 넣을 때마다 ``승인하는 날 이 값을 알 수 있나?''를 묻는다.
\end{warnbox}

\section{시점 학습: \texttt{asof.py}}
\textbf{하려는 것.} 2016년 1월 1일에 승인 결정을 한다면, 그날까지 \textbf{결과가 확정된} 대출로만 배워야 한다. 2015년 9월에 나간 대출이 2016년 3월에 부도났다는 사실은 2016년 1월에는 모른다.
\begin{lstlisting}
CHARGE_OFF_LAG = 5

def months_to(d, cutoff):
    c = pd.Timestamp(cutoff)
    return ((c.year - d['ym'].dt.year) * 12 + c.month - d['ym'].dt.month).to_numpy()

def known_at(d, cutoff, lag=CHARGE_OFF_LAG):
    age = months_to(d, cutoff)
    settle = d['tau'].to_numpy() + np.where(d['bad'].to_numpy() == 1, lag, 0)
    return (age > 0) & d['resolved'].to_numpy() & (settle <= age)

def training_set(d, cutoff, lag=CHARGE_OFF_LAG):
    return d[(d['ym'] < pd.Timestamp(cutoff)).to_numpy() & known_at(d, cutoff, lag)]
\end{lstlisting}
\begin{itemize}[leftmargin=*]
\item \textbf{1번 줄}: 부도는 마지막 납부 뒤 약 150일(5개월) 연체 후에 상각(Charged Off)되므로, 5개월 늦게 확정된다고 본다.
\item \textbf{5번 줄} \ft{age}: 발행월부터 기준일까지 개월 수. (연도 차이 \ensuremath{\times} 12) + 월 차이.
\item \textbf{9번 줄} \ft{settle}: 결과가 확정된 개월. 상환이면 \ft{tau}, 부도면 \ft{tau} + 5.
\item \textbf{10번 줄}: 세 조건을 모두 만족해야 학습에 쓸 수 있다. 기준일 전에 발행됐고, 결과가 확정됐고, 확정 시점이 기준일 이전이다.
\end{itemize}
\begin{workout}[따라 해 보기: \ft{snippets.py asof}]
기준일 2016-01-01, 대출 4건:
\begin{lstlisting}[style=out]
          ym  tau  bad  resolved  age  settle  known
0 2014-01-01   12    0      True   24      12   True
1 2014-06-01    6    1      True   19      11   True
2 2015-03-01    3    0      True   10       3   True
3 2015-09-01    2    1      True    4       7  False
\end{lstlisting}
\Step{1} 3번 대출: 2015년 9월 발행, 2개월 납부 뒤 부도. 확정은 $2+5=7$개월째인데 기준일까지는 4개월뿐이다.
\Ans{3번은 2016년 1월에는 ``아직 모르는'' 대출이라 학습에서 빠진다. 이 규칙이 없으면 미래를 보고 배우는 셈이다.}
\end{workout}
\Trap{기준일 이전에 \emph{발행된} 대출을 모두 쓰는 것. 그러면 빨리 끝난 대출(대부분 조기상환)만 남거나, 아직 모를 부도가 섞인다. 둘 다 결과를 왜곡한다(보고서 3장의 종결 편향).}

\chapter{부도확률 모델과 기대수익 점수}

\section{LightGBM 학습: \texttt{fit\_lgb}}
\textbf{하려는 것.} 신청 정보(특성 \ft{X})로 부도 여부(\ft{y}=\ft{bad})를 맞히는 나무 모델을 학습하고, 평가 연도 대출의 부도확률 $p$를 낸다.
\begin{lstlisting}
LGB_PARAMS = dict(n_estimators=2000, learning_rate=0.05, num_leaves=63, subsample=0.8, subsample_freq=1,
                  colsample_bytree=0.6, min_child_samples=200, reg_lambda=5, cat_smooth=50, n_jobs=-1, verbose=-1)

def fit_lgb(Xt, yt, seed, extra=None):
    p = dict(models.LGB_PARAMS)
    p.update(extra or {})
    u = np.random.default_rng(seed).random(len(Xt)) < 0.15
    m = lgb.LGBMClassifier(random_state=seed, **p)
    m.fit(Xt[~u], yt[~u], eval_set=[(Xt[u], yt[u])], callbacks=[lgb.early_stopping(50, verbose=False)])
    return m
\end{lstlisting}
\begin{itemize}[leftmargin=*]
\item \textbf{1--2번 줄} 설정값(고정, 튜닝 안 함). 핵심만 보면 이렇다.
  \begin{itemize}
  \item \ft{n_estimators=2000}: 나무를 최대 2000그루까지 쌓는다.
  \item \ft{learning_rate=0.05}: 나무 하나가 앞 나무들의 실수를 5\%씩만 고친다. 작을수록 천천히, 안정적으로 배운다.
  \item \ft{num_leaves=63}: 나무 하나가 손님을 최대 63개 무리로 나눈다.
  \item \ft{min_child_samples=200}: 한 무리에 최소 200건. 몇 건에 맞춘 규칙(과적합)을 막는다.
  \item \ft{subsample}, \ft{colsample_bytree}: 나무마다 행 80\%, 열 60\%만 무작위로 써서 나무들이 서로 다르게 배우게 한다.
  \end{itemize}
\item \textbf{7번 줄}: 학습 표본의 15\%를 무작위로 떼어(\ft{u}=참) 검증용으로 둔다.
\item \textbf{9번 줄}: 나머지 85\%(\texttt{\textasciitilde u})로 학습한다. 검증 15\%의 성적이 50그루 연속 나아지지 않으면 멈춘다(조기 종료). 2000은 상한일 뿐이고 실제로는 수백 그루에서 멈춘다.
\end{itemize}

\section{연도마다 학습하고 $p$, 점수 만들기: 30번의 앞부분}
\begin{lstlisting}
EV = []
for Y in range(2013, 2017):
    tr = asof.training_set(d, f'{Y}-01-01')
    tri, ytr = tr.index, tr['bad'].to_numpy()
    usable = [c for c in feats if X0.loc[tri, c].nunique(dropna=True) > 1]
    evi = d.index[(d['issue_year'] == Y) & d['resolved'] & (d['term_m'] == 36)]
    ev = returns.attach_treasury(d.loc[evi]).reset_index(drop=True)
    lgd = asof.loss_given_default(tr)
    p = np.mean([fit_lgb(X0.loc[tri, usable], ytr, s).predict_proba(X0.loc[evi, usable])[:, 1]
                 for s in SEEDS], axis=0)
    ev['p'] = p
    ev['s_er'] = (1 - p) * ev['r'] - p * lgd - ev['tb']
    EV.append(ev)
ev = pd.concat(EV, ignore_index=True)
\end{lstlisting}
\begin{itemize}[leftmargin=*]
\item \textbf{2번 줄}: 평가 연도 2013, 2014, 2015, 2016을 하나씩 돈다(\ft{range(2013, 2017)}은 2017을 포함하지 않는다).
\item \textbf{3번 줄}: 그해 1월 1일 기준으로 알 수 있던 대출만으로 학습 표본을 만든다(2장).
\item \textbf{5번 줄}: 학습 표본에서 값이 한 가지뿐인 열(예: 아직 생기지 않은 컬럼)은 뺀다.
\item \textbf{6번 줄} 평가 대상: 그해 발행된, 결과가 확정된, 36개월 대출.
\item \textbf{7번 줄}: 국채 금리 \ft{tb}를 붙이고 행 번호를 0부터 다시 매긴다.
\item \textbf{8번 줄} LGD: 학습 표본 부도 대출의 평균 손실률, 즉 $1-$(받은 총액/대출액)의 평균이다.
\item \textbf{9--10번 줄}: seed 0, 1, 2로 모델을 세 번 학습해 $p$를 평균 낸다. 무작위 요소(검증 분할, 행·\allowbreak{}열 추출) 때문에 생기는 흔들림을 줄인다. \ft{predict_proba}는 [부도 아닐 확률, 부도일 확률] 두 열을 주고, \ft{[:, 1]}이 두 번째 열, 즉 부도확률이다.
\item \textbf{12번 줄} \textbf{기대수익 점수}: $s=(1-p)\,r-p\cdot\mathrm{LGD}-tb$. 안 망하면 받는 이자, 빼기 망하면 잃는 원금, 빼기 국채 기회비용.
\item \textbf{14번 줄}: 4개 연도 표를 위아래로 붙여 하나의 \ft{ev}로 만든다. 이후 모든 장의 \ft{ev}가 이것이다.
\end{itemize}

\begin{workout}[따라 해 보기: \ft{snippets.py lgbm}]
가짜 데이터 5000건(부도율 약 17\%)으로 같은 설정을 seed 세 개로 학습한다.
\begin{lstlisting}[style=out]
seed 0: 158 trees, AUC 0.982
seed 1: 182 trees, AUC 0.985
seed 2: 76 trees, AUC 0.985
3-seed mean AUC 0.984, mean p 0.161 vs actual 0.166
\end{lstlisting}
\Ans{2000그루 상한이지만 조기 종료로 76--182그루에서 멈췄다. seed마다 나무 수가 다르고, 평균을 내면 조금 더 안정적이다. 가짜 데이터라 AUC가 높고, 실제 LC 데이터에서는 0.63--0.68이다.}
\end{workout}
\Trap{\ft{predict_proba(X)}를 그대로 $p$로 쓰는 것. 두 열짜리 표이므로 \ft{[:, 1]}을 빠뜨리면 모양이 맞지 않는다. 또 최신 LightGBM은 \ft{eval_set} 대신 \ft{eval_X}, \ft{eval_y}를 권한다는 경고를 낼 수 있는데, 결과는 같다.}

\section{국채 붙이기: \texttt{attach\_treasury}}
\begin{lstlisting}
def attach_treasury(d):
    m = d.merge(treasury(), on=['ym', 'term_m'], how='left', validate='many_to_one')
    m.index = d.index
    assert m['tb'].notna().all()
    return m
\end{lstlisting}
\ft{treasury()}는 FRED에서 받은 GS3(3년), GS5(5년) 월별 금리 표다. \textbf{2번 줄}은 발행월과 만기가 같은 금리를 VLOOKUP한다. \ft{validate='many_to_one'}은 ``대출 여러 건에 금리 한 줄''이 맞는지 검사한다. \textbf{4번 줄} \ft{assert}는 ``빈 금리가 하나라도 있으면 즉시 멈춰라''라는 안전장치다.

\chapter{채점: 현금흐름, IRR, 고르기}

\section{현금흐름 복원: \texttt{schedule}, \texttt{cashflows}}
\textbf{하려는 것.} 데이터에는 월별 납부 기록이 없고 ``받은 총액''(\ft{total_pymnt}), ``월 상환액''(\ft{installment}), ``마지막 납부월''(\ft{tau})만 있다. 이것으로 ``몇 달은 정상 납부, 마지막 달에 나머지 일시금''이라고 근사한다.
\begin{lstlisting}
def schedule(d):
    inst = d['installment'].to_numpy(float)
    tot = d['total_pymnt'].to_numpy(float)
    n = d['tau'].to_numpy(int)
    kf = np.minimum(n, np.floor(tot / np.maximum(inst, 1e-9))).astype(int)
    rem = np.maximum(tot - kf * inst, 0.0)
    return kf, rem, n
\end{lstlisting}
\begin{itemize}[leftmargin=*]
\item \textbf{5번 줄} \ft{kf}(정상 납부 개월): 받은 총액 \ensuremath{\div} 월 상환액을 내림한 값. 단 마지막 납부월 \ft{n}을 넘을 수 없다. \ft{np.maximum(inst, 1e-9)}는 0으로 나누기를 막는다.
\item \textbf{6번 줄} \ft{rem}(일시금): 받은 총액에서 정상 납부분을 뺀 나머지. 조기상환이면 남은 원금이고, 부도면 회수금이다.
\end{itemize}
예: 월 33.21달러짜리 대출에서 총 500달러를 받고 마지막 납부가 20개월째라면 다음과 같다.
\[ kf=\min(20,\ \lfloor 500/33.21\rfloor)=15,\qquad rem=500-15\times33.21=1.85 \]
즉 15개월은 33.21달러씩, 20개월째에 1.85달러를 받는다.

\ft{cashflows()}는 이것을 \textbf{묶음 전체의 월별 현금흐름 배열}(0--99개월)로 만든다. 0개월째는 대출액 합계의 음수(투자)이고, 1--\ft{kf}개월은 월 상환액, \ft{n}개월째에 일시금이 들어온다. 수수료 1\%는 받은 돈마다 \ft{(1 - fee)}를 곱한다. 코드의 \ft{np.cumsum}, \ft{np.add.at}은 ``대출마다 반복문을 돌지 않고 한 번에 더하기'' 위한 기술이다. 뜻은 위 문장과 같다.

\section{IRR: \texttt{irr}}
\begin{lstlisting}
def irr(cf):
    cf = np.where(np.abs(cf) < 1e-6, 0.0, cf)
    f = lambda m: (cf / (1 + m) ** np.arange(len(cf))).sum()
    try:
        m = brentq(f, -0.2, 0.2)
    except ValueError:
        m = brentq(f, -0.95, 1.0)
    return (1 + m) ** 12 - 1
\end{lstlisting}
\begin{itemize}[leftmargin=*]
\item \textbf{3번 줄} \ft{f(m)}: 월 할인율 $m$으로 모든 현금흐름을 현재가치로 바꿔 더한 값(NPV). $t$개월째 돈은 $(1+m)^t$로 나눈다.
\item \textbf{5번 줄} \ft{brentq(f, a, b)}: $a$와 $b$ 사이에서 $f(m)=0$이 되는 $m$을 찾는다(엑셀 목표값 찾기). NPV를 0으로 만드는 할인율이 IRR이다.
\item \textbf{6--7번 줄}: 손실이 아주 크면 $-20\%$--$20\%$ 범위에 답이 없어서, 범위를 넓혀 다시 찾는다.
\item \textbf{8번 줄}: 월 수익률을 연 수익률로 바꾼다(복리).
\end{itemize}
\begin{workout}[따라 해 보기: \ft{snippets.py irr}]
1000달러, 연 12\%, 36개월.
\begin{lstlisting}[style=out]
monthly payment 33.21
fully paid: IRR 0.1268
1% fee: IRR 0.1191
default after 10: IRR -0.6814
pool of two: IRR -0.1405
\end{lstlisting}
\Step{1} 다 갚으면 12.68\%다. 12\%가 아닌 이유는 월 1\%를 12번 복리로 쌓으면 $1.01^{12}-1=12.68\%$이기 때문이다.
\Step{2} 수수료 1\%가 연 수익률을 약 0.8\%p 깎는다.
\Step{3} 한 건짜리 부도는 $-68\%$처럼 극단값이 된다. 그래서 대출 하나하나의 IRR로 Sharpe를 내지 않는다(보고서 6장).
\Step{4} 두 대출을 \textbf{현금흐름을 합쳐} 한 묶음으로 보면 $-14\%$다. 두 IRR의 평균($(12.68-68.14)/2=-27.7\%$)과 다르다.
\Ans{묶음 IRR = 현금흐름을 먼저 합치고 IRR 한 번. 이 프로젝트의 채점은 모두 ``발행월 묶음 IRR''이다(\ft{pooled_irr}).}
\end{workout}

\section{고르기: 발행월마다 상위 $q$}
\begin{lstlisting}
def cut(score, q, keys):
    a = np.zeros(len(ev), bool)
    for _, ix in ev.groupby(keys, observed=True).indices.items():
        s = score[ix]
        a[ix] = s >= np.quantile(s, 1 - q)
    return a

jit = np.random.default_rng(11).random(len(ev)) * 1e-6
er = ev['s_er'].to_numpy() + jit            # V1: expected return
lowpd = -ev['p'].to_numpy() + jit           # V2: low default probability
a = cut(er, 0.3, 'ym')                      # approve top 30% within each month
a3 = cut(-acc + jit, 0.3, ['ym', 'grade'])  # V3: within month x grade
\end{lstlisting}
\begin{itemize}[leftmargin=*]
\item \textbf{2번 줄}: 승인 여부 배열을 전부 ``거절''로 시작한다.
\item \textbf{3번 줄}: 발행월(\ft{keys='ym'})마다 그 달 대출의 행 번호 \ft{ix}를 돈다. \ft{_}는 ``이름(그 달 날짜)은 안 쓴다''는 표시다.
\item \textbf{4--5번 줄}: 그 달 점수 중 상위 $q$를 승인으로 바꾼다. \textbf{발행월마다} 자르므로 매달 승인 금액 비중이 거의 일정하다. 연도 단위로 자르면 이 비중이 들쭉날쭉해 Sharpe가 왜곡된다(37단계의 교훈).
\item \textbf{8번 줄} 동점 깨기: 점수에 0--0.000001 사이의 아주 작은 무작위 값을 더한다. 순위는 거의 그대로이고, 동점만 무작위로 갈린다.
\item \textbf{10번 줄}: ``부도확률이 낮을수록 좋다''를 ``점수가 높을수록 좋다''로 맞추려고 마이너스를 붙인다.
\item \textbf{12번 줄}: \ft{keys}에 두 열을 주면 ``발행월 \ensuremath{\times} 등급''마다 자른다. 같은 등급 안에서 고르는 규칙이다.
\end{itemize}
\begin{workout}[따라 해 보기: \ft{snippets.py cut}]
\begin{lstlisting}[style=out]
top 50%: [1 1 1 0 0 1 1 0]
ties as is: [1 1 1 1 1 1 1 1] -> approves all, not 50%
ties + tiny noise: [0 1 1 0 1 1 0 0]
\end{lstlisting}
\Step{1} 첫 줄: 1월 점수 0.9, 0.5, 0.5, 0.1에서 50\% 분위수는 0.5다. 0.5 이상이 셋이라 75\%가 승인됐다. 동점이 있으면 정확히 $q$가 안 된다.
\Step{2} 둘째 줄: 전부 동점이면 전부 승인된다. 34단계에서 실제로 이 버그가 나서, 계좌 개설 수(정수) 규칙의 실제 승인비율이 30\%가 아니라 39.5\%였다.
\Ans{셋째 줄처럼 작은 무작위 값을 더하면 매달 정확히 절반씩 승인된다.}
\end{workout}
\Trap{점수가 정수이거나 같은 값이 많으면 꼭 동점을 깬다. 그리고 고른 뒤에는 \ft{a.mean()}으로 실제 승인비율을 확인한다.}

\paragraph{무작위 기준선.} \ft{cut(np.random.default_rng(500 + k).random(len(ev)), q, 'ym')}. 점수 대신 무작위 수로 같은 비율을 고른다. seed 500--509로 10번(학습 쪽은 5번) 해서 평균 낸다. 이것이 ``아무거나 고를 때''다.

\chapter{Sharpe, 불확실성, 최종 선택, test}

\section{발행월 포트폴리오 초과수익: \texttt{port\_series}}
\begin{lstlisting}
def port_series(a, hurdle):
    out = []
    for m in months:
        ix = grp[m]
        sel = ix[a[ix]]
        r = returns.pooled_irr(ev.iloc[sel], fee=FEE) - hurdle[ix].mean() if len(sel) else 0.0
        out.append(P_amt[sel].sum() / P_amt[ix].sum() * r)
    return np.array(out)

sharpe = lambda x: x.mean(axis=-1) / x.std(axis=-1, ddof=1)
\end{lstlisting}
\begin{itemize}[leftmargin=*]
\item \textbf{3--4번 줄}: 48개 발행월을 차례로 돈다. \ft{grp = ev.groupby('ym').indices}(30번에서 만듦)라서 \ft{grp[m]}은 그 달 행 번호다.
\item \textbf{5번 줄}: 그 달 대출 중 승인된 것만 남긴다. \ft{a[ix]}는 그 달 대출들의 승인 여부이고, 그것으로 \ft{ix}를 거른다.
\item \textbf{6번 줄}: 승인 묶음의 IRR(수수료 1\%)에서 그 달 허들(국채 등)을 뺀다. 승인이 하나도 없으면 0이다.
\item \textbf{7번 줄} \textbf{핵심}: 그 달 포트폴리오 초과수익 = (승인 금액 \ensuremath{\div} 그 달 전체 금액) \ensuremath{\times} 초과수익. 거절한 돈은 국채에 넣었다고 보므로 그 부분의 초과수익은 0이다.
\item \textbf{10번 줄}: Sharpe = 48개 값의 평균 \ensuremath{\div} 표준편차. \ft{axis=-1}은 ``마지막 축으로''라는 뜻이다. 1차원이면 그냥 전체이고, 2차원(부트스트랩 표)이면 행마다 계산한다.
\end{itemize}
\begin{workout}[따라 해 보기: \ft{snippets.py sharpe}]
\begin{lstlisting}[style=out]
share 0.3: mean 0.0124, std 0.0053, Sharpe 2.33
share 0.8: mean 0.0331, std 0.0142, Sharpe 2.33
\end{lstlisting}
\Ans{비중을 0.3에서 0.8로 바꾸면 평균과 표준편차가 모두 같은 배수로 커져 Sharpe는 그대로다. 그래서 승인비율 사이의 Sharpe 차이는 ``얼마나 샀나''가 아니라 ``얼마나 까다롭게 골랐나''의 차이다. 같은 이유로 파레토 비교는 같은 승인비율 안에서 한다(아래).}
\end{workout}

\section{이동 블록 부트스트랩: 표준오차와 구간}
\textbf{하려는 것.} 48개월은 짧다. ``다른 48개월이 왔다면 Sharpe가 얼마나 달랐을까''를 알려고, 있는 달들을 다시 뽑아 가상의 48개월을 2000번 만든다. 이웃한 달끼리 비슷하므로(경기 흐름) 한 달씩이 아니라 \textbf{6개월 덩어리}로 뽑는다.
\begin{lstlisting}
n = len(months)                                   # 48
st = np.random.default_rng(2026).integers(0, n - BLOCK + 1, size=(B, int(np.ceil(n / BLOCK))))
BIX = np.concatenate([np.arange(BLOCK)[None, :] + st[:, [j]] for j in range(st.shape[1])], axis=1)[:, :n]

se = sharpe(port[BIX]).std(ddof=1)                # standard error of Sharpe
diff_bs = sharpe(ch[BIX]) - np.mean([sharpe(r[BIX]) for r in rnd], axis=0)
lo, hi = np.quantile(diff_bs, 0.05), np.quantile(diff_bs, 0.95)   # 90% interval
\end{lstlisting}
\begin{itemize}[leftmargin=*]
\item \textbf{2번 줄}: 덩어리 시작점을 무작위로 뽑는다. 크기는 (2000번) \ensuremath{\times} (덩어리 8개)이고, 시작점은 0--42 사이다(덩어리가 48을 넘지 않게).
\item \textbf{3번 줄}: 시작점 $s$를 $s, s+1, \ldots, s+5$로 늘여 이어 붙인다. 결과 \ft{BIX}는 2000행 \ensuremath{\times} 48열의 ``달 번호 표''다. 행 하나가 가상의 48개월 하나다.
\item \textbf{5번 줄}: \ft{port[BIX]}는 2000 \ensuremath{\times} 48 표다. 행마다 Sharpe를 내면 2000개의 Sharpe가 나오고, 그 흩어짐(표준편차)이 \textbf{표준오차(SE)}다.
\item \textbf{6번 줄} \textbf{짝지은 비교}: 같은 가상 48개월(같은 \ft{BIX} 행)에서 선택안과 무작위를 함께 계산해 차이를 낸다. 시장 전체가 좋았던 달 같은 공통 흔들림이 차이에서 지워져, 차이의 구간이 각각의 SE보다 훨씬 좁다.
\item \textbf{7번 줄}: 2000개 차이의 하위 5\%와 상위 5\%가 90\% 구간이다. 하한이 0보다 크면 ``거의 확실히 무작위보다 낫다''.
\end{itemize}
\begin{workout}[따라 해 보기: \ft{snippets.py boot}]
\ft{n=12, block=3, B=3}으로 줄여 본 \ft{BIX}:
\begin{lstlisting}[style=out]
[[ 8  9 10  1  2  3  0  1  2  6  7  8]
 [ 3  4  5  4  5  6  0  1  2  3  4  5]
 [ 6  7  8  3  4  5  8  9 10  7  8  9]]
\end{lstlisting}
\Ans{행마다 세 달씩 연속인 덩어리 네 개가 이어져 있다. 같은 달이 두 번 뽑히기도 하고(복원 추출), 안 뽑히는 달도 있다.}
\end{workout}

\section{1-SE 규칙으로 최종 제안 고르기: 39번}
\begin{lstlisting}
best = C.loc[C['Sharpe'].idxmax()]
ok = C[C['Sharpe'] >= best['Sharpe'] - best['SE']]
pick = ok.sort_values(['q', 'Sharpe'], ascending=[False, False]).iloc[0]
\end{lstlisting}
\begin{itemize}[leftmargin=*]
\item \ft{C}는 후보 18개(전략 2 \ensuremath{\times} 승인비율 9)의 Sharpe와 SE 표다.
\item \textbf{1번 줄} \ft{idxmax()}: Sharpe가 가장 큰 행의 이름표. \ft{loc}으로 그 행을 꺼낸다.
\item \textbf{2번 줄}: (최고 Sharpe $-$ 그 SE) 이상인 후보만 남긴다. 최고와 통계적으로 구분되지 않는 후보들이다.
\item \textbf{3번 줄}: 승인비율 큰 순, 같으면 Sharpe 큰 순으로 정렬해 첫 줄(\ft{iloc[0]})을 고른다.
\end{itemize}
\begin{workout}[따라 해 보기: \ft{snippets.py onese}]
\begin{lstlisting}[style=out]
best A 0.1: 3.8 - SE 0.8 = bar 3.0
pass: [['A', 0.1], ['A', 0.5], ['B', 0.5]]
pick: A 0.5
\end{lstlisting}
\Ans{최고는 A 10\%지만 기준선 3.0 이상인 후보 중 승인비율이 가장 큰 50\%에서, Sharpe가 더 높은 A를 고른다. 실제 결과(train)도 같은 구조였다. 최고는 부도확률 순 10\%(3.82, SE 0.84), 기준선은 2.98, 최종 제안은 부도확률 순 80\%(3.06)다.}
\end{workout}
\begin{whybox}[왜 최고값을 그냥 고르지 않나]
18개를 비교하면 운 좋게 높게 나온 후보가 1위가 되기 쉽다(다중 비교). SE가 0.7--0.8인데 1위와 2위의 차이가 0.2라면 그 순위는 의미가 없다. 구분되지 않는다면 많이 승인하는 쪽이 금액과 NPV가 크고, 모델이 극단 꼬리를 맞혀야 하는 부담도 적다. 머신러닝에서 규제 강도를 고를 때 쓰는 ``one-standard-error rule''과 같은 생각이다.
\end{whybox}

\section{파레토: 41번의 \texttt{dominated}}
\begin{lstlisting}
def dominated(mean, std):
    m, s = np.asarray(mean), np.asarray(std)
    return np.array([bool(((m >= m[i]) & (s <= s[i]) & ((m > m[i]) | (s < s[i]))).any())
                     for i in range(len(m))])
\end{lstlisting}
\textbf{3번 줄}은 후보 $i$마다 ``평균은 같거나 높고, 표준편차는 같거나 낮고, 둘 중 하나는 엄격히 나은 다른 후보가 하나라도(\ft{.any()}) 있나''를 묻는다. 있으면 $i$는 지배당한다(둘 다 나쁜 조합이라 뺀다).
\begin{workout}[따라 해 보기: \ft{snippets.py pareto}, 숫자는 train 승인 30\%]
\begin{lstlisting}[style=out]
ER frontier
lowPD dominated
simple frontier
random dominated
\end{lstlisting}
\Ans{단순 규칙(평균 1.84, 표준편차 0.48)이 부도확률 순(1.66, 0.50)보다 수익도 높고 흔들림도 작아 부도확률 순을 지배한다. 기대수익 순(1.88, 0.65)은 수익이 가장 높아 경계에 남는다. 다만 단순 규칙은 train을 보고 고른 규칙이라 test에서 확인해야 한다.}
\end{workout}

\section{test 실행: 40 \ensuremath{\to} 35 \ensuremath{\to} 15}
\begin{lstlisting}[language=bash,numbers=none]
python scripts/40_test_evaluate_v2_2.py --dry-run          # rehearsal on train
python scripts/40_test_evaluate_v2_2.py data/raw/test.csv  # the real test, once
\end{lstlisting}
40번은 35번 전체를 실행하고(35번은 다시 15번을 실행한다), 그 위에 v2.2 판정을 덧붙인다. 그래서 한 번 실행으로 v1, v2.1, v2.2가 모두 나온다. 결과는 \ft{results/test_v2_2/}에 생긴다.
\begin{table}[H]\centering\small
\begin{tabularx}{\linewidth}{@{}p{45mm}L@{}}\toprule
파일 & 무엇을 보나\\\midrule
\ft{summary.md} & 전체 요약. 맨 아래가 v2.2 판정\\
\ft{verdict.csv} & v1 가설 H1--H4 지지 여부\\
\ft{verdict_v2.csv} & v2.1 가설 H5--H9, 위약 점검\\
\ft{verdict_v2_2.csv} & H10(제안 유지), H11, H12, V3 도전자, H13(파레토)\\
\ft{V5_candidates_v2_2.csv} & test에서 18개 후보의 Sharpe와 순위\\
\ft{V6_dominance_v2_2.csv} & 같은 승인비율 지배 관계\\\bottomrule
\end{tabularx}
\end{table}
\begin{warnbox}[test는 한 번만]
결과를 보고 코드나 규칙을 고치면 test는 더 이상 증거가 아니다. 고쳐야 할 것이 생기면 \ft{docs/test_protocol_v2_2.md}의 ``이탈 기록''에 날짜, 내용, 이유를 적고, 고치기 전 결과도 함께 보고한다.
\end{warnbox}

\section{막힐 때}
\begin{table}[H]\centering\small
\begin{tabularx}{\linewidth}{@{}p{55mm}L@{}}\toprule
오류 메시지 & 뜻과 해결\\\midrule
\ft{command not found: python} & conda 환경이 꺼져 있다. \ft{conda activate cmapss}\\
\ft{ModuleNotFoundError: lc} & \ft{import _path}가 빠졌거나 \ft{scripts/} 밖에서 파일을 옮겼다\\
\ft{KeyError: 'xxx'} & 그런 열 이름이 없다. \ft{print(df.columns)}로 확인\\
\ft{ValueError: substring not found} & \ft{exec}로 자르는 기준 문자열이 앞 파일에서 바뀌었다\\
\ft{IndentationError} & 들여쓰기 칸 수가 섞였다(탭과 공백). 공백 4칸으로 통일\\
결과 숫자가 문서와 조금 다름 & seed, 승인 컷 단위(발행월인지 연도인지), 허들을 확인\\\bottomrule
\end{tabularx}
\end{table}

\chapter{개념 치트시트: 신용 대출 선별 전략을 검증하는 법}
\begingroup\small\setlist{nosep,leftmargin=*}\setlength{\columnsep}{6mm}

\raggedcolumns
\titlespacing*{\subsection}{0pt}{6pt}{2pt}
\begin{multicols}{2}\raggedright\small

\begin{keybox}[딱 이것만]
\begin{enumerate}
\item \textbf{예상으로 고르고, 실제로 채점한다.} 점수(승인일에 아는 정보)와 IRR(끝난 뒤 실제 현금흐름)을 섞지 않는다.
\item \textbf{같은 비율 무작위와 비교한다.} 시장과 승인비율을 맞춰야 ``고르는 실력''만 남는다.
\item \textbf{그때 알던 것만 쓴다.} 학습도, 승인 컷도 시점을 지킨다.
\item \textbf{판정 기준은 결과 전에 적고 커밋한다.} test는 한 번만 본다.
\end{enumerate}
\end{keybox}

\subsection*{1. 고르는 점수}
\begin{formulabox}[기대수익 점수]
$s=(1-p)\,r-p\cdot\mathrm{LGD}-tb$\\
$p$ 부도확률, $r$ 약정 금리, LGD 부도 시 손실률(약 0.5--0.6), $tb$ 같은 달 같은 만기 국채.
\end{formulabox}
\begin{itemize}
\item 부도확률만 보면 늘 안전한 대출을 사지만, 위험에 비해 금리가 후한 고금리 대출을 놓친다.
\item 1년 단위 근사라 \textbf{순위용}이다. 정확한 예상 수익으로 읽지 않는다.
\end{itemize}

\subsection*{2. 채점: IRR}
\begin{itemize}
\item 현금흐름 $c_0<0, c_1, \ldots$에서 $\sum_t c_t/(1+m)^t=0$인 월 $m$, 연 $(1+m)^{12}-1$.
\item \textbf{묶음 IRR}: 현금흐름을 먼저 합치고 IRR을 한 번 구한다. IRR들의 평균과 다르다.
\item 한 건짜리 IRR은 부도 시 $-70\%$처럼 극단이라, 건별 통계에 그대로 쓰지 않는다.
\item 수수료 1\%(받은 돈마다)가 연 약 0.8\%p를 깎는다.
\end{itemize}

\subsection*{3. 선별 가치(주 지표)}
\begin{formulabox}
$\mathrm{SV}=\frac19\sum_{q=10\%}^{90\%}\big[\overline{(\mathrm{IRR}-tb)}_{\text{상위 }q}-\overline{(\mathrm{IRR}-tb)}_{\text{무작위 }q}\big]$
\end{formulabox}
\begin{itemize}
\item 발행월마다 자르고, 9개 비율을 평균한다(유리한 비율 하나를 고르지 못하게).
\item 결과: \textbf{+1.07\%p}(90\% 0.94--1.21), 48개월 전부 양수.
\item \textbf{위약 점검}: 무작위 점수를 넣으면 0 근처여야 한다.
\end{itemize}

\subsection*{4. 성과를 나누는 도구}
\begin{itemize}
\item \textbf{Brinson}: 초과 = 배분(좋은 등급을 많이 샀나) + 선별(같은 등급 안에서 잘 골랐나) + 상호작용.
\item \textbf{등급 중립 롱숏}: 매달 등급마다 신호 상위 30\% $-$ 하위 30\%. 등급과 시장이 지워져 \textbf{등급 안 선별}만 남는다(측정 도구일 뿐 실제 전략은 아님). 기대수익 신호 +2.38\%p.
\end{itemize}

\subsection*{5. Sharpe는 정의부터}
\begin{tabularx}{\linewidth}{@{}p{24mm}X@{}}\toprule
정의 & 특징(우리 값)\\\midrule
건별, 거절 = 0 & 승인비율만 줄여도 값이 변한다. 선택 기준으로 부적합(0.03--0.15)\\
건별, 승인만 & 조교 기준 0.2--0.4와 크기가 맞는다(0.10--0.34)\\
\textbf{발행월 포트폴리오} & 매달 비중 $w$ \ensuremath{\times} (묶음 IRR $-$ 허들), 48개월 평균 \ensuremath{\div} 표준편차. $w$가 지워진다(약 2--4)\\\bottomrule
\end{tabularx}
\begin{warnbox}[자르는 단위]
연도 단위로 상위 $q$를 자르면 달마다 $w$가 들쭉날쭉해 포트폴리오 Sharpe가 깎인다(2.46 vs 발행월 2.91). 실전에서도 12월을 보고 1월을 정할 수 없다. \textbf{발행월마다 자른다.}
\end{warnbox}
\begin{itemize}
\item \textbf{허들}에 따라 1위가 바뀐다. 국채면 부도확률 순, BAA면 기대수익 순이 Sharpe 1위다. 허들별로 모두 보고한다.
\item Sortino는 순위가 Sharpe와 같아 새 정보가 없었다.
\end{itemize}

\subsection*{6. 시점 학습(as-of)}
\begin{itemize}
\item 기준일 $C$의 학습 표본 = $C$ 전에 발행 $\wedge$ 결과 확정 $\wedge$ 확정 시점 $\le C$.
\item 확정 시점: 상환은 마지막 납부월, 부도는 그 뒤 5개월(150일 연체 후 상각).
\item \textbf{종결 편향}: 최근 발행분 중 이미 끝난 대출은 대부분 조기상환이다. 만기 도달 비율이 0.9 미만인 해는 따로 보고한다.
\item \textbf{누수}: 받은 총액, 회수금, 마지막 납부일 등 승인 후 정보는 채점에만 쓰고 모델에는 넣지 않는다.
\end{itemize}

\subsection*{7. 불확실성}
\begin{itemize}
\item \textbf{이동 블록 부트스트랩}: 6개월 덩어리로 48개월을 2000번 다시 뽑는다. 이웃 달의 상관을 보존한다.
\item \textbf{짝지은 비교}: 같은 가상 표본에서 두 전략의 차이를 계산한다. 공통 흔들림이 지워져 구간이 좁다(Sharpe SE 0.8 vs 차이 구간 폭 0.25).
\item 판정은 \textbf{90\% 구간 하한 > 0}으로, 크기보다 방향과 구간을 본다.
\end{itemize}

\subsection*{8. 여러 후보 중 고르기: 1-SE 규칙}
\begin{enumerate}
\item 후보마다 지표와 SE를 구한다.
\item 최고값 $-$ 그 SE 이상인 후보 = 허용 범위(동등성 검정 아님).
\item 그중 \textbf{가장 단순하거나 규모가 큰} 것을 고른다(여기서는 승인비율 최대).
\end{enumerate}
결과: 최고 부도확률 순 10\%(3.82, SE 0.84) \ensuremath{\to} 기준선 2.98 \ensuremath{\to} \textbf{부도확률 순 80\%(3.06)}. 무작위보다 +0.34(0.30--0.55). 짝지어 보면 10\%가 0.76 높다(구분됨): Sharpe를 내주고 규모(약 8.5배)를 얻는 선택이다. (이전 최종안. v2.3부터 최종 제안은 기대수익 순, 예산 무관.)

\subsection*{9. 두 축과 파레토}
\begin{itemize}
\item Sharpe = 평균 \ensuremath{\div} 표준편차이므로 \textbf{세 번째 목표가 아니다}. 축은 수익(평균)과 흔들림(표준편차) 둘이다.
\item \textbf{지배}: 수익은 같거나 높고 흔들림은 같거나 낮으며, 하나는 엄격히 나음. 둘 다 나쁜 조합은 뺀다.
\item \textbf{같은 승인비율 안에서} 비교한다. 비율이 바뀌면 규모와 구성이 함께 바뀌고, 비율을 섞으면 거의 다 비지배라 걸러 내지 못한다.
\item 결과: 기대수익 순 대 부도확률 순은 진짜 맞교환(서로 지배 안 함). 단순 규칙은 30--50\%에서 부도확률 순을 지배(사후 규칙, test로 확인).
\end{itemize}

\subsection*{10. 정직한 검증 습관}
\begin{itemize}
\item \textbf{사전 등록}: 판정 기준을 실험 전에 적고 git 태그로 고정한다(v1 \ensuremath{\to} v2 \ensuremath{\to} v2.1 \ensuremath{\to} v2.2). 개정은 test 전에만 하고, 이유를 남긴다.
\item \textbf{주장 장부}: 주장마다 근거 파일, 상태(잠정, 철회), 확인 방법을 적는다.
\item \textbf{사후 규칙}: 데이터를 보고 찾은 규칙은 과대평가된다. 후보에서 빼고 도전자로만 검증한다.
\item \textbf{독립 재구현}: 핵심 숫자는 같은 함수를 재사용하지 말고 다르게 한 번 더 계산한다(연도 컷 실수를 이렇게 놓쳤다).
\item \textbf{동점 점검}: 정수 점수는 동점이 많다. 고른 뒤 실제 승인비율을 확인한다.
\end{itemize}

\subsection*{11. 다른 과제에 옮겨 쓰기}
\begin{tabularx}{\linewidth}{@{}p{22mm}X@{}}\toprule
상황 & 바꿔 끼울 것\\\midrule
다른 대출 데이터 & 상태 분류, LGD, 상각 지연 개월\\
주식 선별 & 점수 = 예상 수익, IRR \ensuremath{\to} 월 수익률, 등급 \ensuremath{\to} 업종\\
마케팅 타깃 & 점수 = 반응확률 \ensuremath{\times} 이익 $-$ 비용, 무작위 발송과 비교\\
어떤 모델이든 & 같은 비율 무작위 기준선, 시점 분할, 블록 부트스트랩, 사전 등록\\\bottomrule
\end{tabularx}
\end{multicols}

\endgroup

\chapter{코드 치트시트: 복사해서 쓰는 조각}
모두 \ft{docs/notes/code/snippets.py}에 실행 가능한 형태로 있다. \ft{python snippets.py irr}처럼 하나씩 실행한다.
\begingroup\setlist{nosep,leftmargin=*}\setlength{\columnsep}{6mm}\lstset{basicstyle=\ttfamily\scriptsize,numbers=none,aboveskip=3pt,belowskip=3pt}

\begin{multicols}{2}\raggedright\small

\Hd{0. 불러오기와 실행}
\begin{lstlisting}
import numpy as np
import pandas as pd
from scipy.optimize import brentq
\end{lstlisting}
\begin{lstlisting}[language=bash]
conda activate cmapss
python scripts/39_final_choice.py > results/39.log 2>&1
\end{lstlisting}
\ft{> 파일 2>&1}: 화면 출력과 오류를 함께 파일에 저장한다.

\Hd{1. pandas 기본 동작}
\begin{lstlisting}
df = pd.read_csv('loans.csv')          # or pd.read_parquet
df['bad'] = (df['status'] == 'Charged Off').astype(int)
df['ym'] = pd.to_datetime(df['issue_d'], format='%b-%Y')
df['r'] = pd.to_numeric(df['int_rate'].str.rstrip('%')) / 100
m = (df['grade'] == 'A') | (df['r'] > 0.14)   # & | ~ with brackets
sub = df[m]                                   # filter rows
df.groupby('grade')['bad'].mean()             # rate per group
idx = df.groupby('ym').indices                # {month: row positions}
df = df.merge(tb, on=['ym', 'term_m'], how='left')   # VLOOKUP
df = df.reset_index(drop=True)                # rows 0..n-1 again
x = df['amt'].to_numpy(float)                 # fast numpy array
\end{lstlisting}

\Hd{2. numpy 기본 동작}
\begin{lstlisting}
np.where(bad == 1, 0.0, r)        # IF for whole array
np.quantile(s, 0.7)               # 70th percentile
s >= np.quantile(s, 1 - q)        # top q as True/False
rng = np.random.default_rng(11)   # fixed seed = reproducible
x.mean(); x.std(ddof=1)           # sample std
X2d.mean(axis=1)                  # mean of each row
\end{lstlisting}

\Hd{3. 시점 학습 표본}
\begin{lstlisting}
def known_at(d, cutoff, lag=5):
    c = pd.Timestamp(cutoff)
    age = (c.year - d['ym'].dt.year) * 12 + c.month - d['ym'].dt.month
    settle = d['tau'] + np.where(d['bad'] == 1, lag, 0)
    return (age > 0) & d['resolved'] & (settle <= age)

train = d[(d['ym'] < cutoff) & known_at(d, cutoff)]
\end{lstlisting}

\Hd{4. LightGBM 부도확률(seed 평균)}
\begin{lstlisting}
import lightgbm as lgb
def fit_lgb(X, y, seed):
    u = np.random.default_rng(seed).random(len(X)) < 0.15
    m = lgb.LGBMClassifier(n_estimators=2000, learning_rate=0.05,
            num_leaves=63, min_child_samples=200, subsample=0.8,
            subsample_freq=1, colsample_bytree=0.6, reg_lambda=5,
            random_state=seed, verbose=-1)
    m.fit(X[~u], y[~u], eval_set=[(X[u], y[u])],
          callbacks=[lgb.early_stopping(50, verbose=False)])
    return m
p = np.mean([fit_lgb(Xtr, ytr, s).predict_proba(Xte)[:, 1]
             for s in [0, 1, 2]], axis=0)
\end{lstlisting}
\ft{[:, 1]} = 부도(1)일 확률 열. 빠뜨리면 두 열짜리 표가 된다.

\Hd{5. 기대수익 점수와 LGD}
\begin{lstlisting}
b = train[train['bad'] == 1]
lgd = 1 - (b['total_pymnt'] / b['funded_amnt']).mean()
ev['s_er'] = (1 - p) * ev['r'] - p * lgd - ev['tb']
\end{lstlisting}

\Hd{6. IRR}
\begin{lstlisting}
def irr(cf):                       # cf[0] < 0, monthly
    f = lambda m: (cf / (1 + m) ** np.arange(len(cf))).sum()
    m = brentq(f, -0.95, 1.0)
    return (1 + m) ** 12 - 1

def annuity(P, rate, T):           # level monthly payment
    g = rate / 12
    return P * g / (1 - (1 + g) ** -T)

inst = annuity(1000, 0.12, 36)                 # 33.21
irr(np.r_[-1000, np.full(36, inst)])           # 0.1268
irr(np.r_[-1000, np.full(36, inst) * 0.99])    # 1% fee: 0.1191
\end{lstlisting}
묶음 IRR: 대출들의 현금흐름 배열을 \textbf{먼저 더하고} \ft{irr}을 한 번.

\Hd{7. 그룹마다 상위 $q$(동점 깨기)}
\begin{lstlisting}
def cut(score, groups, q):
    a = np.zeros(len(score), bool)
    for _, ix in pd.Series(groups).groupby(groups).indices.items():
        s = score[ix]
        a[ix] = s >= np.quantile(s, 1 - q)
    return a

jit = np.random.default_rng(11).random(len(s)) * 1e-6
a = cut(s + jit, ym, 0.3)              # by month
a = cut(s + jit, list(zip(ym, grade)), 0.3)   # month x grade
print(a.mean())                        # check real share
rand = cut(np.random.default_rng(500).random(len(s)), ym, 0.3)
\end{lstlisting}

\Hd{8. 발행월 포트폴리오 초과수익과 Sharpe}
\begin{lstlisting}
def port_series(a, hurdle):
    out = []
    for m_, ix in sorted(df.groupby('ym').indices.items()):
        sel = ix[a[ix]]
        r = pooled_irr(df.iloc[sel]) - hurdle[ix].mean() if len(sel) else 0.0
        out.append(amt[sel].sum() / amt[ix].sum() * r)
    return np.array(out)

sharpe = lambda x: x.mean(axis=-1) / x.std(axis=-1, ddof=1)
\end{lstlisting}

\Hd{9. 이동 블록 부트스트랩}
\begin{lstlisting}
def block_index(n, block=6, B=2000, seed=2026):
    st = np.random.default_rng(seed).integers(
        0, n - block + 1, size=(B, int(np.ceil(n / block))))
    return np.concatenate([np.arange(block)[None, :] + st[:, [j]]
        for j in range(st.shape[1])], axis=1)[:, :n]

BIX = block_index(48)
se = sharpe(x[BIX]).std(ddof=1)                # SE of Sharpe
d = sharpe(x[BIX]) - sharpe(y[BIX])            # paired difference
lo, hi = np.quantile(d, [0.05, 0.95])          # 90% interval
\end{lstlisting}

\Hd{10. 1-SE 규칙}
\begin{lstlisting}
best = C.loc[C['Sharpe'].idxmax()]
ok = C[C['Sharpe'] >= best['Sharpe'] - best['SE']]
pick = ok.sort_values(['q', 'Sharpe'], ascending=[False, False]).iloc[0]
\end{lstlisting}

\Hd{11. 파레토 지배}
\begin{lstlisting}
def dominated(mean, std):
    m, s = np.asarray(mean), np.asarray(std)
    return np.array([bool(((m >= m[i]) & (s <= s[i])
        & ((m > m[i]) | (s < s[i]))).any()) for i in range(len(m))])
\end{lstlisting}

\Hd{12. 결과 저장과 표 보기}
\begin{lstlisting}
T = pd.DataFrame(rows)                       # rows = list of dicts
T.round(4).to_csv('results/x.csv', index=False)
print(T.pivot(index='q', columns='rule', values='Sharpe').round(2))
\end{lstlisting}
\ft{pivot}: 긴 표를 엑셀 피벗처럼 행 \ensuremath{\times} 열로 펼친다.

\Hd{13. 자주 나는 오류}
\begin{tabularx}{\linewidth}{@{}p{42mm}X@{}}\toprule
\ft{KeyError} & 열 이름 오타. \ft{print(df.columns)}\\
\ft{ValueError} (truth value) & \ft{and}/\ft{or} 대신 \texttt{\&}/\texttt{|}, 조건마다 괄호\\
\ft{SettingWithCopyWarning} & 거른 표에 쓰기 전 \ft{.copy()}\\
\ft{IndentationError} & 공백 4칸으로 통일\\
\ft{brentq} 부호 오류 & 범위를 \ft{(-0.95, 1.0)}으로 넓힘\\\bottomrule
\end{tabularx}
\end{multicols}

\endgroup

\end{document}
